\documentclass[11pt]{article}
\usepackage[utf8]{inputenc}
\usepackage[T1]{fontenc}
\usepackage{amsmath,amssymb,amsthm}
\usepackage{graphicx}
\usepackage{booktabs}
\usepackage{array}
\usepackage[margin=2.4cm]{geometry}
\usepackage{hyperref}

\graphicspath{{figures/}}
\DeclareMathOperator{\Tr}{Tr}
\newcommand{\C}{\mathbb{C}}
\newcommand{\M}{\mathcal{M}}
\newtheorem{proposition}{Proposition}
\newtheorem{theorem}[proposition]{Theorem}
\newtheorem{lemma}[proposition]{Lemma}

% link to a program in man/article/listings of the repository
\newcommand{\repofile}[1]{\href{https://github.com/andre-a-no/von_neumann_type_i/blob/main/man/article/listings/#1}{\texttt{listings/\allowbreak\detokenize{#1}}}}

% numbers and tables produced by scripts/*.py (bash scripts/run_all.sh [quick|full] [device])
\newcommand{\gen}[1]{\IfFileExists{generated/#1.tex}{\input{generated/#1.tex}}{}}
\gen{validation}\gen{search}\gen{bounds}\gen{chains}\gen{numerics}\gen{baselines}\gen{benchmark}\gen{renyi_threshold}\gen{connes_stormer}\gen{permutational}\gen{trace_characterizations}\gen{symmetry_resolved}
% Captions that use these numbers have an empty short caption: the short caption is expanded into main.aux,
% where a number whose file is missing would be an undefined control sequence.
\providecommand{\ValidationRows}{\multicolumn{6}{c}{\emph{run scripts/validate\_known\_results.py}}\\}
\providecommand{\SearchRows}{\multicolumn{9}{c}{\emph{run scripts/inequality\_search.py}}\\}
\providecommand{\BoundsRows}{\multicolumn{7}{c}{\emph{run scripts/bounds\_search.py}}\\}
\providecommand{\BenchmarkRows}{\multicolumn{5}{c}{\emph{run scripts/benchmark\_library.py}}\\}
\providecommand{\BenchHasGPU}{0}
\providecommand{\TCRows}{\multicolumn{10}{c}{\emph{run scripts/trace\_characterizations.py}}\\}
\providecommand{\NumTimeRows}{\multicolumn{9}{c}{\emph{run scripts/numerics.py}}\\}
\providecommand{\NumAccRows}{\multicolumn{5}{c}{\emph{run scripts/numerics.py}}\\}
\providecommand{\NumTimeHeader}{generate & eigh $f(A)$ & SVD $|A|$ & channel & expm & opt.\ step}
\providecommand{\BaseMCRows}{\multicolumn{7}{c}{\emph{run scripts/baselines.py}}\\}
\providecommand{\BaseGSRows}{\multicolumn{7}{c}{\emph{run scripts/baselines.py}}\\}
\providecommand{\ValTraceRows}{\multicolumn{5}{c}{\emph{run scripts/validate\_known\_results.py}}\\}
\providecommand{\ChainKrylovRows}{\multicolumn{6}{c}{\emph{run scripts/chains.py}}\\}
\providecommand{\BaseMinKS}{--}\providecommand{\BoundsStChecked}{--}\providecommand{\BoundsStOutside}{--}\providecommand{\BoundsShortfall}{--}
\providecommand{\ChainRingRows}{\multicolumn{6}{c}{\emph{run scripts/chains.py}}\\}
\providecommand{\RenyiRows}{\multicolumn{5}{c}{\emph{run scripts/renyi\_threshold.py}}\\}
\providecommand{\BoundsSupDevOne}{(not generated)}\providecommand{\BoundsSupDevThree}{(not generated)}
\providecommand{\RenyiMode}{--}\providecommand{\RenyiDevice}{--}\providecommand{\RenyiTol}{--}
\providecommand{\CSTwoRows}{\multicolumn{6}{c}{\emph{run scripts/connes\_stormer.py}}\\}
\providecommand{\CSRows}{\multicolumn{10}{c}{\emph{run scripts/connes\_stormer.py}}\\}
\newcommand{\fig}[2]{\IfFileExists{figures/#1}{\includegraphics[width=#2]{#1}}{\fbox{\small figure \detokenize{#1} not generated}}}

\title{Finite-Dimensional Type~I von Neumann Algebras in PyTorch:\\
A GPU-Accelerated Framework for Random Block-Diagonal Operators}
\author{Irina Nikolaeva \and Andrej Novikov}
\date{}

\begin{document}
\maketitle

\begin{abstract}
\texttt{torch\_vn\_algebra} is a PyTorch library for batched, differentiable computations on CPUs and GPUs in the
finite-dimensional Type~I von Neumann algebras $\bigoplus_cM_{k_c}(\C)$ of quantum systems with conserved charges or
superselection rules: random operators with prescribed spectra, states, channels between sectors, spin chains, Lindblad
dynamics, inclusions of algebras, symmetry reduction and constrained optimisation. Validated against exact results, it
serves here as a research tool in four use cases. We determine how strongly trace-characterising inequalities fail
for non-tracial functionals: several violation envelopes are sharp functions of the Michelson contrast of the
density, given by the Kantorovich constant ($\kappa-1$), the eigenvalue ratio ($a/b-1$) and the Specht ratio, and the
inequalities detect only the contrast inside a block. For two maximal abelian subalgebras of $M_2$ we prove a closed
form for the Connes--St{\o}rmer entropy $H(\Delta_2|U\Delta_2U^*)$, stated as open by Bakshi and Guin and by Petz,
Sz\'ant\'o and Weiner; for $n\ge3$ we observe partitions of unity that beat the best orthonormal bases found. For a
boundary time crystal with local emission or dephasing we compute finite-$N$ lifetimes up to $N=80$ and observe that
the collective and local damping rates approximately add. For the XXZ chain at $L=24$, batched over eight anisotropies, the
sector blocks give the symmetry-resolved entanglement directly, and the charge statistics give the Luttinger parameter
to about one per cent away from the isotropic point. Constrained optimisation also replaces the Monte Carlo sampling of our
earlier work on trace inequalities, whose envelopes we prove for sectors of every size, and rediscovers the
data-processing thresholds $\alpha=\frac12$ and $\alpha=2$ of the Rényi divergences.
\end{abstract}

\section*{Program summary}
\begin{description}\itemsep0pt
\item[Program title:] \texttt{torch\_vn\_algebra}
\item[Developer's repository link:] \url{https://github.com/andre-a-no/von_neumann_type_i}
\item[Licensing provisions:] MIT
\item[Programming language:] Python ($\ge3.8$), PyTorch ($\ge1.13$), NumPy ($\ge1.21$)
\item[Nature of problem:] Numerical experiments in finite-dimensional Type~I von Neumann algebras
$\bigoplus_cM_{k_c}(\C)$, which describe quantum systems with superselection rules or conserved charges:
random operators with prescribed spectra, several trace functionals, states, channels within and between
sectors, symmetric tensor products, spin chains, Hamiltonian and Lindblad dynamics, extremal problems for
trace inequalities under spectral constraints, inclusions of algebras, symmetry reduction, and quantum Rényi
divergences.
\item[Solution method:] Batched block-diagonal tensors processed by PyTorch linear algebra on CPUs and GPUs;
exact Haar sampling; spectral functional calculus; Kraus representations; Lanczos and Krylov methods for
sparse sector Hamiltonians; differentiable parametrisations with batched multistart gradient optimisation;
run-time and memory estimates before expensive steps; inclusion matrices; randomised Wedderburn decomposition
from random elements of the commutant; highest-weight (Yamanouchi) bases and exact Clebsch--Gordan coefficients.
\item[Additional comments:] more than 800 unit tests, run in continuous integration; every number, table and figure of
this paper is produced by the scripts and examples in the repository and regenerated by \texttt{scripts/run\_all.sh}.
\end{description}

\section{Introduction}\label{sec:intro}

\paragraph{Where block-diagonal algebras arise.} Whenever a conserved quantity, a symmetry or a superselection
rule splits a Hilbert space into sectors, the physically relevant operators are block diagonal. Their algebra
is a finite-dimensional Type~I von Neumann algebra
\begin{equation}\label{eq:M}
  \M=\bigoplus_{c=1}^{C}M_{k_c}(\C)\quad\text{acting on}\quad\mathcal H=\bigoplus_{c=1}^C\C^{k_c},
\end{equation}
where $C$ is the number of sectors (direct summands) and $M_{k}(\C)$ is the algebra of complex $k\times k$
matrices. Every finite-dimensional von Neumann algebra is $*$-isomorphic to one of this form~\cite{takesaki2002}, and its centre is
spanned by the sector projections. The library also accepts nominal factor sizes $n_c\ge k_c$: this describes
computations in the corner $p\M_0p\cong\bigoplus_cM_{k_c}(\C)$ of a larger algebra $\M_0=\bigoplus_cM_{n_c}(\C)$,
with $p=\bigoplus_cp_c$ and $p_c$ a projection of rank $k_c$. The corner is again of the form~\eqref{eq:M}; all
operators, traces and normalisations refer to it, and in most applications $k_c=n_c$.
\paragraph{From a symmetry to the algebra.} The link between physics and~\eqref{eq:M} is the commutant. If a group
$G$ (or a set of conserved quantities) acts unitarily on $\mathcal H$ and $\mathcal H\cong\bigoplus_\lambda
V_\lambda\otimes\C^{m_\lambda}$ is its decomposition into irreducible representations $V_\lambda$ with
multiplicities $m_\lambda$, then the operators commuting with the action form the algebra
$\bigoplus_\lambda \mathbb 1_{V_\lambda}\otimes M_{m_\lambda}(\C)\cong\bigoplus_\lambda M_{m_\lambda}(\C)$. For one
conserved self-adjoint charge $Q$ (particle number, magnetisation, Hamming weight) the irreducibles are
one-dimensional and $m_\lambda$ is the dimension of the eigenspace of $Q$ with eigenvalue $\lambda$. A Hamiltonian
that conserves $Q$, a time evolution generated by it, and every observable compatible with a superselection rule
for $Q$ lie in this commutant; the sectors are the eigenvalues of $Q$, and the centre consists of the functions
of $Q$. Physical modelling with a symmetry is therefore computation in an algebra of the form~\eqref{eq:M}; for
non-abelian symmetries the block sizes $m_\lambda$ are smaller than the dimensions of the charge sectors, and
Section~\ref{sec:symmetries} computes the decomposition numerically from the group or its generators.

Examples are abundant:
\begin{itemize}\itemsep0pt
\item superselection rules for charge, parity or particle number~\cite{wick1952,bartlett2007} and the
classical-quantum structure produced by decoherence and einselection~\cite{zurek2003};
\item many-body Hamiltonians with conserved charges, whose exact diagonalisation and dynamics are carried out
sector by sector~\cite{sandvik2010,weinberg2017}, and symmetry-resolved entanglement~\cite{goldstein2018,xavier2018};
\item quantum information: decoherence-free subspaces and noiseless subsystems~\cite{zanardi1997,knill2000},
operator-algebra quantum error correction~\cite{beny2007}, and quantum algorithms that preserve a conserved
quantity, such as the XY mixers of the quantum alternating operator ansatz, which keep the Hamming
weight~\cite{hadfield2019};
\item random matrix theory, where symmetries decompose ensembles into blocks and spectral statistics are
meaningful only inside a block~\cite{dyson1962,mehta2004};
\item machine learning: equivariant layers are the linear maps that commute with a group
representation~\cite{cohen2016,finzi2021}; by Schur's lemma~\cite{fulton1991}, for a complex unitary representation
on one space they form an algebra of the form~\eqref{eq:M}, whose block sizes are the multiplicities of the
irreducible representations, and equivariant layers between two representations are the corresponding off-diagonal
blocks (for real representations the blocks can also be real or quaternionic~\cite[\S3.5]{fulton1991}); random-matrix methods are used to analyse trained
weights~\cite{martin2021};
\item operator theory, our original motivation: inequalities between traces of products characterise the trace
and the centre of an operator algebra~\cite{gardner1979,petz1988,novikov2015,virosztek2017}, and the size of
their failure was studied numerically against spectral quantities such as the Michelson contrast~\cite{abed2024}.
\end{itemize}
On $\M$ the trace is not unique: every positive weighting of the sectors gives one, and statements such as
``trace preserving'' or ``the dual map'' depend on that choice as soon as maps change the sector structure.

\paragraph{What existing software offers.} Table~\ref{tab:software} compares the features relevant here.
Dense and sparse linear algebra (NumPy/SciPy~\cite{harris2020,virtanen2020}, with ARPACK~\cite{lehoucq1998} for
sparse eigenproblems) handles single matrices; QuTiP~\cite{johansson2013,lambert2026} and the quantum-information module of
Qiskit~\cite{qiskit2019} represent states, operators and channels on one Hilbert space, and PIQS, part of QuTiP,
exploits permutational symmetry for the Lindblad dynamics of $N$ two-level systems~\cite{shammah2018}; \texttt{dynamiqs}~\cite{dynamiqs}
adds batching, GPUs and automatic differentiation for Schrödinger and Lindblad dynamics; QuSpin~\cite{weinberg2017}
builds symmetry-reduced bases for exact diagonalisation; tensor-network libraries (ITensor~\cite{fishman2022},
TeNPy~\cite{hauschild2018}) store charge-conserving tensors in block-sparse form, and YASTN~\cite{rams2025}, the
closest to this work in its software stack, does so on a PyTorch backend; RTNI~\cite{fukuda2019} integrates
Haar-random tensor networks symbolically. None of them treats the algebra~\eqref{eq:M} itself, with its centre,
several traces, states and maps between algebras of different sector structure, as the basic object, and none
combines it with batched sampling and optimisation over operators with prescribed spectra.

\begin{table}[h]
\centering\footnotesize
\setlength{\tabcolsep}{3pt}
\begin{tabular}{lccccccc}
\toprule
 & NumPy/ & QuTiP & Qiskit & \texttt{dynamiqs} & QuSpin & ITensor/ & this\\
 & SciPy & & q.\,info & & & TeNPy & work\\
\midrule
GPU execution & -- & ($\checkmark$) & -- & $\checkmark$ & -- & ($\checkmark$) & $\checkmark$\\
batched over many operators & ($\checkmark$) & -- & -- & $\checkmark$ & -- & -- & $\checkmark$\\
automatic differentiation & -- & ($\checkmark$) & -- & $\checkmark$ & -- & ($\checkmark$) & $\checkmark$\\
symmetry sectors of many-body systems & -- & -- & -- & -- & $\checkmark$ & $\checkmark$ & $\checkmark$\\
algebra $\bigoplus_cM_{k_c}$ with centre and several traces & -- & -- & -- & -- & -- & -- & $\checkmark$\\
channels within and between sectors & -- & ($\checkmark$) & ($\checkmark$) & -- & -- & -- & $\checkmark$\\
prescribed spectra, Haar/COE/CSE per sector & ($\checkmark$) & ($\checkmark$) & ($\checkmark$) & -- & -- & -- & $\checkmark$\\
optimisation under spectral constraints & -- & -- & -- & -- & -- & -- & $\checkmark$\\
large single systems & $\checkmark$ & $\checkmark$ & $\checkmark$ & $\checkmark$ & $\checkmark$ & $\checkmark$ & ($\checkmark$)\\
\bottomrule
\end{tabular}
\caption{Features relevant to computations in the algebra~\eqref{eq:M}, to the best of our knowledge at the
time of writing ($\checkmark$: supported; ($\checkmark$): partly, through plug-ins, or with user code;
--: not supported). The last row is where the specialised packages are the better choice. The QuTiP column refers to
QuTiP~5~\cite{lambert2026}.
YASTN~\cite{rams2025} has the marks of ITensor/TeNPy, with GPU execution and automatic differentiation through its
PyTorch backend.}
\label{tab:software}
\end{table}

\paragraph{Why not one large matrix?} Nothing forces one to treat $\M$ as a separate object: it embeds into
$M_D(\C)$, $D=\sum_ck_c$, as the block-diagonal matrices, and every package above handles a single $D\times D$
matrix (the case $C=1$ of~\eqref{eq:M}). The embedding is exact for linear algebra, so the question is only what it
costs and what has to be supplied by hand.
\begin{itemize}\itemsep0pt
\item \emph{Cost.} A product, an inverse or a spectral decomposition in $M_D$ costs $O(D^3)$ time and $O(D^2)$ memory,
in $\M$ it costs $\sum_c k_c^3$ and $\sum_c k_c^2$. For $C$ equal blocks this is a factor $C^2$ in time and $C$ in
memory; for the half-filling sector of a chain of $L=16$ spins ($D=2^{16}$, largest block $12870$) the dense $M_D$
computation is out of reach, the block computation is routine.
\item \emph{Measures.} The unitary group of $\M$ is $\prod_cU(k_c)$, for $C\ge2$ a subgroup of $U(D)$ of measure zero. A Haar-random
element of $U(D)$ is not in $\M$, and a random element of $\M$ has to be produced block by block. The same holds for
random states, for the circular ensembles and for any optimisation: an unconstrained gradient step in $M_D$ leaves $\M$.
\item \emph{Traces.} On $M_D$ the trace is unique up to a factor; on $\M$ the faithful traces form a $C$-parameter
family $\sum_cw_c\Tr_c$. Densities, ``trace preserving'' and the dual of a map between algebras with different sectors
depend on the choice of the weights; on the embedded matrices this information is no longer visible.
\item \emph{Centre.} The centre of $\M$ (the sector projections) is what physicists call the classical or superselected
part of the system; conditioning on a sector, post-selection and sector-resolved quantities are operations with the
centre, and they have to be reconstructed from the block indices when $\M$ is stored as one matrix.
\end{itemize}
The library keeps these four pieces of structure inside the data type. For $C=1$ it reduces to batched
linear algebra in $M_k(\C)$, and all its functions accept that case. The first version of this
library~\cite{v1} shows that the bookkeeping is error-prone: its complex Haar sampler was not Haar distributed, its
circular ensembles were not the circular ensembles, and several conclusions drawn from Monte Carlo samples were wrong;
all of this is corrected and tested here.

\paragraph{Challenges.} Turning the algebra into an efficient and reliable tool raises concrete problems, each
addressed by a design decision described in Section~\ref{sec:design}: (i) heterogeneous block sizes have to be
mapped onto batched GPU kernels; (ii) trace conventions and dual maps must stay consistent when maps change the
sectors; (iii) random ensembles must be exactly Haar (or COE/CSE) in every block; (iv) gradients propagated through
eigen- or singular-vector decompositions (as in $f(A)$ or $|X|$) are numerically singular at degenerate spectra,
exactly where extremal operators tend to lie; (v) extremal
values sit on sets of tiny measure that random sampling does not reach in higher dimensions
(Section~\ref{sec:optimisation}); (vi) batched workloads easily exceed the time or memory budget of a device;
(vii) the number field (real or complex) and the precision change both cost and results
(Section~\ref{sec:numerics}).

\paragraph{Objective and scope.} \texttt{torch\_vn\_algebra} is a library for numerical experiments \emph{in}
the algebra~\eqref{eq:M}: sampling and optimising operators, states and channels, evaluating trace and spectral
functionals, and formulating physical models (spin chains, open systems, noisy symmetric quantum circuits) whose
symmetry is encoded in the sector structure. It is not a circuit simulator and not a tensor-network code, and for
single large systems the specialised packages above remain the better choice; Section~\ref{sec:baselines}
quantifies where the library gains and where it does not.

\paragraph{Contributions and outline.} Sections~\ref{sec:math}--\ref{sec:symmetries} describe the library:
notation and conventions (Section~\ref{sec:math}); the software stack, the workflow, how the challenges (i)--(vii) are
met, and the number field and precision (Section~\ref{sec:design}); an overview of short complete programs in the
repository, including a noisy quantum circuit (Section~\ref{sec:usage}); validation against exact results (Section~\ref{sec:validation}); constrained
optimisation instead of Monte Carlo sampling, which proves the envelopes of the trace inequalities of~\cite{v1} for
sectors of every size and locates the data-processing thresholds of the Rényi divergences (Sections~\ref{sec:optimisation}
and~\ref{sec:renyi}); charge-conserving tensor products and spin chains with momentum sectors, sparse Krylov solvers
and particle loss (Section~\ref{sec:chains}); and inclusions of algebras with conditional expectations and indices, the
Wedderburn decomposition of the algebra generated by a symmetry, and the SU(2) symmetry of spin chains
(Section~\ref{sec:symmetries}). Section~\ref{sec:usecases} uses the library on four research questions: how strongly
characterisations of the trace fail for non-tracial functionals (Section~\ref{sec:tc}), the Connes--St{\o}rmer entropy
of two maximal abelian subalgebras of $M_n$, in closed form for $n=2$ (Section~\ref{sec:cs}), the lifetimes of a
boundary time crystal under local dissipation (Section~\ref{sec:btc}), and symmetry-resolved entanglement in the XXZ
chain (Section~\ref{sec:sre}). Section~\ref{sec:baselines} compares with NumPy/SciPy, Section~\ref{sec:performance}
reports performance, and Section~\ref{sec:limitations} discusses limitations.

\section{Mathematical setting}\label{sec:math}

\paragraph{Algebra and traces.} Let $\M=\bigoplus_{c=1}^C M_{k_c}(\C)$ act on
$\mathcal H=\bigoplus_c\C^{k_c}$; $A\in\M$ is block diagonal, $A=\bigoplus_c A_c$. Every trace on $\M$ is
$\operatorname{tr}_w(A)=\sum_c w_c\Tr A_c$ with weights $w_c\ge0$, and it is faithful iff all $w_c>0$. The library implements
\[
  \Tr_{\mathrm{blunt}}:\ w_c=1,\qquad \Tr_{\mathrm{norm}}:\ w_c=\tfrac1{k_c},\qquad
  \tau_{\mathrm{vN}}:\ w_c=\tfrac1{Ck_c},
\]
where $\tau_{\mathrm{vN}}$ is the faithful normal tracial state that gives every sector the same weight
(the library allows empty sectors, $k_c=0$; they carry no weight, and $C$ in $\tau_{\mathrm{vN}}$ counts the non-empty ones).
For $A\in\M$, $|A|=(A^*A)^{1/2}$; for $A\ge 0$ the Michelson contrast is
$\Delta(A)=(\lambda_{\max}-\lambda_{\min})/(\lambda_{\max}+\lambda_{\min})\in[0,1]$, with the extreme
eigenvalues taken over all sectors, for $A\ne0$ (the library returns $\Delta(0)=0$); for $A\ge0$, $A\ne0$, $\Delta(A)=0$ if and only if $A$ is a multiple of the identity.

\paragraph{States.} A state $\omega$ on $\M$ is represented by its density $\rho\ge0$ with respect to
$\Tr_{\mathrm{blunt}}$, $\omega(A)=\Tr(\rho A)$, $\sum_c\Tr\rho_c=1$. The restriction of $\omega$ to the
centre is the classical distribution $p_c=\Tr\rho_c$ over sectors. The density of the same state with
respect to $\operatorname{tr}_w$ is $\rho^{(w)}_c=\rho_c/w_c$; the library converts between the
conventions explicitly, because statements such as ``trace preserving'' or ``the dual map'' depend on
them as soon as the sectors change.

\paragraph{Maps between sectors.} Let $\M_{\mathrm{in}}=\bigoplus_c M_{k_c}$ and
$\M_{\mathrm{out}}=\bigoplus_d M_{m_d}$. Every normal completely positive map
$\Phi:\M_{\mathrm{in}}\to\M_{\mathrm{out}}$ (in the Schrödinger picture) is a sum of completely
positive maps between blocks,
\[
  \Phi(\rho)_d=\sum_c\sum_i K^{dc}_i\,\rho_c\,(K^{dc}_i)^*,\qquad K^{dc}_i:\C^{k_c}\to\C^{m_d}.
\]
It preserves $\Tr_{\mathrm{blunt}}$ iff $\sum_{d,i}(K_i^{dc})^*K_i^{dc}=1_c$ for all $c$, and
$T_{dc}=\Tr\Phi_{dc}(\rho_c)/\Tr\rho_c$ are the transition probabilities between sectors for the input
state $\rho$ (defined for $\Tr\rho_c>0$) (independent of $\rho$ iff every $\sum_i(K_i^{dc})^*K_i^{dc}$ is a multiple of $1_c$; by default the
library uses the maximally mixed state of each input sector). The dual
with respect to $\operatorname{tr}_w$ on both sides is
$\Phi^{*,w}(A)_c=\sum_{d,i}(w^{\mathrm{out}}_d/w^{\mathrm{in}}_c)(K_i^{dc})^*A_dK_i^{dc}$; for
sector-preserving maps on one algebra the weights cancel and all three traces give the same dual.

\paragraph{Tensor products and fusion of sectors.} $\bigl(\bigoplus_cM_{k_c}\bigr)\otimes\bigl(\bigoplus_dM_{m_d}\bigr)
=\bigoplus_{(c,d)}M_{k_cm_d}$ is again of Type~I with sectors labelled by pairs. If the sectors carry
conserved charges $q_c$, $q'_d$ (particle numbers, magnetisation), the operators of the composite system
that conserve only the \emph{total} charge form the larger algebra
\[
  \M_1\mathbin{\hat\otimes}\M_2=\bigoplus_q M_{n_q},\qquad n_q=\sum_{q_c+q'_d=q}k_cm_d ,
\]
in which all pairs with the same total charge are merged into one sector. It contains $\M_1\otimes\M_2$
(block-diagonally in every merged sector), but also operators such as $S^+_iS^-_j$ that move charge
between the factors. For additive charges, and distinct charges on distinct sectors of $\M_1$, the partial trace
$(\Tr_2X)_c=\sum_d\Tr_{m_d}X_{(c,d),(c,d)}$ still lands in $\M_1$: a block between pairs $(c,d)\ne(c',d)$ with equal
total charge would need $q_c=q_{c'}$, hence $c=c'$.
In both cases the partial trace is a map between sectors in the sense above, and its dual with respect
to $\Tr_{\mathrm{blunt}}$ is the embedding $A\mapsto A\otimes1$.

\paragraph{Spin chains.} For $L$ spins $1/2$ with conserved $S^z_{\mathrm{tot}}$, iterating the fused
product of single-site algebras $M_1\oplus M_1$ (charges $0,1$) gives
$\M=\bigoplus_{N=0}^LM_{\binom LN}(\C)$, $N$ the number of up spins, which contains every
$S^z_{\mathrm{tot}}$-conserving Hamiltonian, e.g.\ the XXZ chain
$H=\sum_{\langle ij\rangle}\bigl[\tfrac J2(S^+_iS^-_j+S^-_iS^+_j)+J\Delta S^z_iS^z_j\bigr]+\sum_ih_iS^z_i$ (the
anisotropy $\Delta$ is unrelated to the Michelson contrast $\Delta(\cdot)$).

\paragraph{Dynamics.} For $H=H^*\in\M$ and jump operators $L_j\in\M$ the
Gorini--Kossakowski--Sudarshan--Lindblad equation
$\dot\rho=-i[H,\rho]+\sum_j\gamma_j(L_j\rho L_j^*-\tfrac12\{L_j^*L_j,\rho\})$ preserves every sector,
so the sector probabilities $p_c$ are conserved (superselection). Jumps between sectors, such as particle loss,
are described by maps between sectors as above (Section~\ref{sec:chains}).

\section{Design and implementation}\label{sec:design}

\paragraph{Software stack.} The library is pure Python on top of PyTorch~\cite{paszke2019} ($\ge1.13$) and NumPy~\cite{harris2020}. All dense linear algebra goes through \texttt{torch.linalg}, i.e.\ LAPACK
(MKL or OpenBLAS) on CPUs and cuSOLVER, cuBLAS and MAGMA on NVIDIA GPUs; sparse products use PyTorch's CSR
tensors. Automatic differentiation is PyTorch's reverse mode. The validation and paper scripts additionally use
SciPy, pandas, Matplotlib and mpmath (for the high-precision certification of Section~\ref{sec:renyi}). Packaging follows \texttt{pyproject.toml}; the test suite runs in continuous
integration on Python~3.8--3.13 (CPU) and with the oldest supported versions (Python~3.8, PyTorch~1.13, NumPy~1.21), and the paper scripts state the device and mode in the table captions and record the matmul precision, the
PyTorch, CUDA and cuDNN versions and the compute capability in a JSON file next to every table.

\paragraph{Workflow.} A computation proceeds in four steps. (1) An algebra is declared by its sector sizes,
optionally sector charges, number field, precision and device (\texttt{TypeIAlgebra}, or a model such as
\texttt{SpinChain} that creates it). (2) Objects are created in it: random operators with prescribed spectra,
states, channels, Hamiltonians, or differentiable parametrisations. (3) They are combined lazily (products,
functional calculus, channel application, tensor products, time evolution). (4) Functionals are evaluated for the
whole batch at once (traces, spectra, contrasts, entropies, sector probabilities) and reduced to statistics, or
fed into an optimiser that updates the parameters of step~(2). Every step keeps the batch dimension, so the same
code serves one operator and $10^5$ samples.

\begin{table}[h]
\centering\footnotesize
\begin{tabular}{p{0.3\textwidth}p{0.62\textwidth}}
\toprule
challenge & solution in the library\\
\midrule
(i) blocks of different size on batched kernels & padded $(B,C,k_{\max},k_{\max})$ layout; products and traces in one batched call, decompositions in one batched call per distinct block size; single-sector algebras and sparse sector Hamiltonians where padding would dominate\\
(ii) conventions between sectors & states stored as $\Tr_{\mathrm{blunt}}$-densities with explicit conversions;
duals of maps between sectors computed with the weights of the chosen trace\\
(iii) exact random ensembles & Mezzadri's QR algorithm with phase correction, $\det^{-1/n}$ for $SU(n)$, $W^TW$ and
$W^RW$ for COE/CSE, validated against Weingarten moments and spacing statistics\\
(iv) singular gradients at degenerate spectra & functional calculus differentiated by the Daleckii--Krein formula
(divided differences, finite at degenerate spectra); parametrisations that build operators from their spectra in the
optimisation, exact $|X|$ from the parametrisation\\
(v) extremes missed by sampling & batched multistart optimisation with exact contrast constraints\\
(vi) time and memory budgets & probing of batched kernels with an ETA warning, memory estimates checked against the
free device memory, ETA for time stepping\\
(vii) field and precision & two switches per algebra, conversions between them, measured trade-offs\\
\bottomrule
\end{tabular}
\caption{Implementation challenges (Section~\ref{sec:intro}) and how they are addressed.}
\label{tab:challenges}
\end{table}

\paragraph{Data layout.} An operator is a tensor of shape $(B,C,k_{\max},k_{\max})$: $B$ samples, $C$
sectors, block $c$ in the top-left $k_c\times k_c$ corner and zeros elsewhere. Linear algebra is applied
to the batch of active blocks of each sector, so one call to a batched LAPACK/cuSOLVER routine serves all
samples. The padding costs memory when the sector sizes differ strongly (for a chain of 14 spins,
77\,\% of a padded operator is zeros); the library warns in that case, and \texttt{SpinChain} lets one work
in single sectors $M_{\binom LN}$ without padding, which is what dense diagonalisation needs anyway.

\paragraph{Objects.} \texttt{Operator} (with \texttt{DensityMatrix} as a subclass for states),
\texttt{Channel} (sector preserving) and \texttt{InterSectorChannel} (between sectors or algebras) are
thin wrappers around such tensors; Table~\ref{tab:modules} lists the modules. Operations on operators
are lazy: they record how to compute the result and do so on first use; the result is cached and the
recipe dropped, so intermediates of long expressions can be freed. Laziness defers and skips work but
does not fuse operations; where this matters the library uses direct formulas, e.g.\
$\Tr(\rho A)=\sum_{ij}\rho_{ij}A_{ji}$ in $O(k^2)$ instead of forming $\rho A$.

\begin{table}[h]
\centering\small
\begin{tabular}{lp{0.78\textwidth}}
\toprule
module & content \\
\midrule
\texttt{algebra} & \texttt{TypeIAlgebra}, \texttt{Operator}: random operators, unitary ensembles, traces,
functional calculus\\
\texttt{states} & \texttt{DensityMatrix}, random states, Gibbs states, Lüders rule, entropies, fidelity\\
\texttt{channels} & \texttt{Channel}, \texttt{InterSectorChannel}, Choi matrices, standard and random channels\\
\texttt{composite} & plain and fused (charge-conserving) tensor products, \texttt{kron}, partial traces\\
\texttt{chains} & \texttt{SpinChain}: $U(1)$ and momentum sectors, XXZ Hamiltonians with batched disorder,
sparse form, reduced states, jump operators between sectors\\
\texttt{krylov} & sparse sector Hamiltonians, batched Lanczos ground states, Krylov $e^{-itH}\psi$\\
\texttt{dynamics} & exact Schrödinger/von Neumann evolution, RK4, Lindblad equation, $e^{t\mathcal L}$\\
\texttt{functional} & functional calculus: self-adjoint, normal, singular values, holomorphic (contour); spectral cache\\
\texttt{optimize} & parametrisations with prescribed contrast, batched multistart \texttt{extremize}\\
\texttt{cost} & run-time and memory estimates, warnings and errors\\
\texttt{hilbert\_space} & the padded block layout of $\bigoplus_c\C^{k_c}$ and its random subspace bases\\
\texttt{inclusions} & \texttt{Inclusion}: inclusions $\mathcal N\subset\M$ by inclusion matrices, conditional expectations,
Markov trace, Jones, Watatani and Pimsner--Popa indices, basic construction, spin-chain tower\\
\texttt{decompose} & numerical Wedderburn decomposition of the algebra generated by matrices or a finite group; commutants;
real, complex and quaternionic types\\
\texttt{su2} & \texttt{SU2Chain}: SU(2) multiplets of spin-$\frac12$ chains, highest-weight bases, Clebsch--Gordan
coefficients, Wigner--Eckart theorem\\
\texttt{divergences} & Petz and sandwiched Rényi divergences and their limits, Petz and universal recovery maps,
$L^p$ and Kosaki norms\\
\texttt{jordan} & Jordan product, quadratic representation and triple product; Jordan algebras generated by
Hermitian matrices; spin factors\\
\bottomrule
\end{tabular}
\caption{Modules of \texttt{torch\_vn\_algebra}.}
\label{tab:modules}
\end{table}

\paragraph{Random unitaries.} Haar unitaries on $U(n)$ and $O(n)$ are obtained from the QR decomposition
of a Ginibre matrix with i.i.d.\ complex (respectively real) Gaussian entries after fixing the phases (signs) of
$\operatorname{diag}R$~\cite{mezzadri2007}; $SU(n)$ and $SO(n)$ by multiplying with $\det^{-1/n}$ or by
a sign flip of one column; the circular orthogonal and symplectic ensembles as $W^TW$ and $W^RW$ with
$W$ Haar on $U(n)$ ($n$ even for $W^RW$, $W^R$ the quaternion dual)~\cite{mehta2004}. Section~\ref{sec:validation} checks all of them.

\paragraph{Functional calculus.} The module \texttt{functional} provides $f(A)$ at four levels of generality,
all built on batched kernels, one call per distinct block size~\cite{higham2008}. (1) For self-adjoint $A$,
$f(A)=Vf(\Lambda)V^*$ from the Hermitian eigensolver, for any function on the real spectrum. (2) For normal
$A=H+iK$ (unitaries), $H$ and $K$ commute, and the eigenvectors of the Hermitian matrix $H+\alpha K$ diagonalise $A$ for
all but finitely many $\alpha$; the library tries a few irrational values, refines the eigenvectors of nearly tied values
with a second coefficient and corrects the remaining residual to first order, so normal operators also need only
the Hermitian eigensolver (in single precision this decomposition is computed in double precision). (3) Singular values, $|A|$, the pseudo-inverse and the polar decomposition
come from the Hermitian eigensolver of the dilation $\left(\begin{smallmatrix}0&A\\A^*&0\end{smallmatrix}\right)$, whose
eigenvalues are $\pm s_i$: batched SVD on GPUs is limited to blocks of size at most $32$, while the batched Hermitian
eigensolver is not. (4) For arbitrary $A$ and $f$ holomorphic near the spectrum,
$f(A)=(2\pi i)^{-1}\oint f(z)(z-A)^{-1}dz$ is evaluated by the trapezoidal rule on a circle~\cite{trefethen2014},
i.e.\ by batched LU solves, on a circle just around the spectrum, with the number of nodes doubled until two
results agree; the error estimate includes the rounding error of cancelling terms, and an unconverged result raises
an error instead of returning a wrong matrix. An opt-in cache keeps the
eigendecomposition of an operator, so that many functions of it (for example $e^{-itH}$ for many $t$) cost matrix
products only.

\paragraph{Differentiation.} All operations are differentiable. Gradients propagated naively through
eigenvectors contain factors $1/(\lambda_i-\lambda_j)$ and fail at degenerate spectra, which are typical here
(projections, symmetric operators, extremal operators). The functional calculus $f(A)=Vf(\Lambda)V^*$ therefore
has its own backward pass: by the Daleckii--Krein formula~\cite{daleckii1965,bhatia1997}, the derivative of $f$ at
a self-adjoint $A$ is $dF=V\,(L\circ V^*dA\,V)\,V^*$ with the Loewner matrix of divided differences
$L_{ij}=(f(\lambda_i)-f(\lambda_j))/(\lambda_i-\lambda_j)$, replaced by $f'$ for coinciding eigenvalues, which is
finite whenever $f$ is differentiable on the spectrum (for $f\in C^1$ it is the Fr\'echet derivative). Backward
passes through the eigenvectors themselves stay finite at degenerate spectra only for perturbations or losses that
do not mix a degenerate eigenspace~\cite{kasim2020}; they then miss the entries $f'(\lambda)$ of $L$ inside the
degenerate blocks, which the Daleckii--Krein formula includes. Gradients with respect to
parameters of $f$ (a time in $e^{-itA}$) are propagated as usual. The same applies to $\sqrt A$, $\log A$, $A^p$,
and to $|A|$ and the inverse of self-adjoint $A$, which use the Hermitian eigensolver instead of the SVD. For normal $A$ the same formula with complex divided differences is the derivative of the
holomorphic calculus, for holomorphic $f$. $|A|$ of a general $A$ is differentiated as the square root of $A^*A$,
with $L_{ij}=1/(s_i+s_j)$, and the pseudo-inverse by the constant-rank formula; the contour integral is
differentiated through the solves. Only gradients through the eigenvectors returned by \texttt{eigh} keep the
factors $1/(\lambda_i-\lambda_j)$; the parametrisations of Section~\ref{sec:optimisation} avoid them. \texttt{Operator.eigenvalues()} computes spectra without forming eigenvectors, which also saves time and
memory.

\paragraph{GPUs, tensor cores and TF32.} The library computes in FP32 or FP64 (real or complex); it does not use
half precision. GPU generations differ strongly in what this costs. Tensor cores, absent on the P100 used for the
first version, accelerate matrix products on Volta and later GPUs; from Ampere (A100) on, FP32 products may be
executed in the TF32 format with a 10-bit mantissa, which is fast but accurate only to about $10^{-3}$ relative.
PyTorch (since 1.12) keeps true FP32 for matrix products by default, and the library does not change this setting;
the paper scripts expose a \texttt{--tf32} switch to measure the effect and record the setting in every table.
TF32 is acceptable where the statistical error dominates (moments and histograms of many samples) and not where
small differences matter (eigenvalues near degeneracies, the validation tables, trace inequalities near equality). Eigen- and singular-value solvers do not use
tensor cores. FP64 throughput is half of FP32 on data-centre GPUs (P100, V100, A100, H100), the latter two with FP64
tensor cores, but only 1/32--1/64 on consumer GPUs. We therefore recommend a data-centre GPU for double-precision
work. The full-size runs of this paper were made on a consumer GPU (the tables state the hardware); the
double-precision timings there are therefore pessimistic for the GPU.

\subsection*{Number field and precision}\label{sec:numerics}

Every algebra is created with two independent switches, \texttt{complex\_valued} (real or complex
matrices) and \texttt{precision} (\texttt{'single'} or \texttt{'double'}), giving the storage types
\texttt{float32}, \texttt{float64}, \texttt{complex64} and \texttt{complex128}; everything built on the
algebra (random operators, states, channels, chains, parametrisations) inherits them.
\texttt{alg.like(complex\_valued=..., precision=...)} creates the same sector structure with other
numerics and \texttt{op.cast(alg2)} moves an operator there, e.g.\ to explore in single precision and
re-evaluate the interesting cases in double precision.

\paragraph{The field is part of the model.} Choosing real matrices is not only a matter of speed: the
symmetry group and the random ensembles change. Haar matrices of a real algebra are orthogonal, so
moments, level statistics and extremal values refer to $O(n)$ instead of $U(n)$
(e.g.\ $\mathbb E|U_{11}|^4=3/(n(n+2))$ instead of $2/(n(n+1))$, Table~\ref{tab:validation}); random
real symmetric and complex Hermitian matrices have different level statistics, with mean spacing ratios
$0.5307$ and $0.5996$ for GOE and GUE (large-$n$ numerical values of~\cite[Table~I]{atas2013}; the $3\times3$
surmise gives $4-2\sqrt3\approx0.5359$ and $2\sqrt3/\pi-\frac12\approx0.6027$)
(our samples: \IfFileExists{generated/numerics.tex}{\NumRGOE\ and \NumRGUE}{not generated}). Real arithmetic is \emph{exact} (loses nothing) when the problem is
real: spectra and eigenvectors of real symmetric matrices, ground states and their entanglement for
Hamiltonians without complex hoppings, statistics of time-reversal-invariant systems, and also some
extremal problems — for real $X,Y\ge0$ the supremum of $|\Tr(XUY)|$ over unitaries is attained by a real
orthogonal matrix (the polar decomposition of the real matrix $YX$ is real). Complex matrices are
\emph{required} for unitary time evolution $e^{-iHt}$ and Lindblad dynamics with a Hamiltonian (the library
refuses to run them in a real algebra), for Haar unitaries, CUE/CSE and $SU(n)$, for random pure states
and quantities such as Page's entropy, and for systems without an antiunitary symmetry squaring to $+1$ (orbital fluxes, complex
couplings).

\paragraph{Precision.} Single precision resolves about $10^{-7}$ relative to the largest entry. It is
sufficient for Monte Carlo statistics whose statistical error is much larger (moments, histograms,
spacing ratios with $10^3$--$10^5$ samples) and for the first, exploratory pass of an optimisation.
Double precision is needed when the quantity of interest is small compared with the operators: an
inequality checked near equality, small spectral gaps, entropies and logarithms of nearly singular
states, contrasts close to $1$, comparisons with exact results beyond six digits, long RK4 integrations
and optimisation to high relative accuracy. Table~\ref{tab:numerics-acc} shows how drastic the
difference can be: the von Neumann entropy of an almost pure state with eigenvalues of order $10^{-9}$ is
wrong by more than $100\,\%$ in single precision, and the round-off floor of a trace inequality that
holds with equality is several orders of magnitude higher in single than in double precision (Table~\ref{tab:numerics-acc}).

\paragraph{Cost.} An entry takes 4, 8, 8 and 16 bytes for \texttt{float32}, \texttt{float64},
\texttt{complex64} and \texttt{complex128}, so memory, and the batch that fits, changes by a factor up to four
(\texttt{alg.bytes\_per\_operator(B)} reports it). A complex multiply--add costs about four real ones, which bounds the
slowdown of complex matrix products; for eigen- and singular-value solvers, and for double against single precision
on CPUs, the measured ratios are smaller (Table~\ref{tab:numerics-time}). On GPUs double precision depends on the
hardware: data-centre GPUs execute
FP64 at half the FP32 rate, consumer GPUs at 1/32--1/64 of it, so that double precision there can be an
order of magnitude slower. Table~\ref{tab:numerics-time} gives the measured ratios for the current run;
the recommendation that follows is to use the smallest type that is exact for the model (real if the
problem is real) and the precision that the size of the effect requires.

\begin{table}[h]
\centering\small
\begin{tabular}{lrrrrrrrr}
\toprule
type & B/entry & MB/batch & \NumTimeHeader\\
\midrule
\NumTimeRows
\bottomrule
\end{tabular}
\caption[]{Bytes per entry, memory of one batch [MB], and time relative to \texttt{float32} for $C=\IfFileExists{generated/numerics.tex}{\NumC}{8}$
sectors of size $k=\IfFileExists{generated/numerics.tex}{\NumK}{16}$ and a batch of
\IfFileExists{generated/numerics.tex}{\NumB}{B} operators\IfFileExists{generated/numerics.tex}{ (\NumMode\ run
on \NumDevice; float32 times per sample: \NumRefTimes)}{}.}
\label{tab:numerics-time}
\end{table}

\begin{table}[h]
\centering\small
\begin{tabular}{lcccc}
\toprule
type & $\max\|UU^*-1\|$ & $\|\sqrt A^2-A\|/\|A\|$ & rel.\ error $S(\rho)$, $\rho\approx$ pure & $|z|$ at equality\\
\midrule
\NumAccRows
\bottomrule
\end{tabular}
\caption[]{Accuracy of the same computations in the four storage types: Haar unitaries ($n=64$), SVD square
root, entropy of $\rho=(1-\varepsilon)|0\rangle\langle0|+\varepsilon\,1/n$ with $\varepsilon=10^{-7}$, and the
round-off floor of $z=\Tr(Y|X|Y)-\Tr|YXY|$ for commuting $X,Y$ (exactly $0$)\IfFileExists{generated/numerics.tex}{; \NumMode\ run on \NumDevice}{}.}
\label{tab:numerics-acc}
\end{table}

\paragraph{Safeguards.} Batched computations make it easy to request hours of work or more memory than
the device has. The \texttt{cost} module therefore (a) times the first few items of every large
batched eigen-, singular-value or QR decomposition and matrix exponential and issues a warning with the expected
total time when it exceeds a threshold (10\,s by default), (b) compares the estimated peak memory of
large allocations with the free memory of the device, warning above 50\,\% and raising an
\texttt{InsufficientMemoryError} above 90\,\%, and (c) after the second step of a time-stepping
or optimisation loop, warns with the expected duration if it exceeds the same threshold; a progress bar is optional. All thresholds can
be changed, and the checks can be switched off locally.

\section{Usage by example}\label{sec:usage}

The repository contains short complete programs in \texttt{man/article/listings}; the test suite runs them, so they
stay in sync with the code. Each runs on a CPU in under half a minute and uses a GPU when available.
Table~\ref{tab:examples} lists the nine introductory examples; the programs of later sections are described where
they are used.

\begin{table}[h]
\centering\scriptsize
\begin{tabular}{>{\raggedright\arraybackslash}p{0.25\textwidth}>{\raggedright\arraybackslash}p{0.39\textwidth}%
>{\raggedright\arraybackslash}p{0.28\textwidth}}
\toprule
program & what it shows & key API\\
\midrule
\repofile{ex1_random_operators.py} & $10^4$ positive operators in $M_2\oplus M_3\oplus M_4$ with contrast $0.6$ and
Haar-random eigenvectors in every sector; the three traces and spectral functions &
\texttt{TypeIAlgebra}, \texttt{operator\_from\_eigenvalues}, \texttt{michelson\_contrast}, \texttt{Tr\_blunt},
\texttt{tau\_vN}, \texttt{apply\_function}\\
\repofile{ex2_states_measurement.py} & Lüders conditioning on a projective measurement moves all probability into one
sector; expectation values do not depend on the trace that represents the state &
\texttt{states.random\_density\_matrix}, \texttt{condition\_on}, \texttt{sector\_probabilities}, \texttt{density}\\
\repofile{ex3_channels.py} & particle loss of two fermionic modes as a map between the sectors $N=0,1,2$; unital
and non-unital channels and the contrast & \texttt{InterSectorChannel.from\_blocks},
\texttt{transition\_matrix}, \texttt{random\_mixed\_unitary\_channel}\\
\repofile{ex4_dynamics.py} & spontaneous emission of a two-level atom by RK4 and by the exponential of the
generator, against $P_e(0)e^{-\gamma t}$ and $|\rho_{eg}(0)|e^{-\gamma t/2}$ & \texttt{dynamics.lindblad\_evolve},
\texttt{dynamics.lindblad\_channel}\\
\repofile{ex5_page_entropy.py} & mean entanglement entropy of random pure states on $\C^m\otimes\C^n$ against
Page's formula~\cite{page1993} & \texttt{tensor\_product}, \texttt{partial\_trace},
\texttt{states.von\_neumann\_entropy}\\
\repofile{ex6_inequality_search.py} & Haar sampling against batched gradient ascent for
$\sup_U(|\Tr(XUY)|-\Tr(XY))=\|YX\|_1-\Tr(XY)$ (Section~\ref{sec:optimisation}) &
\texttt{random\_unitary\_operator}, \texttt{expm}, \texttt{trace\_norm}\\
\repofile{ex7_spin_chain.py} & Heisenberg ring of ten spins: the sector of the ground state in a uniform field,
the mean level-spacing ratio of 100 disorder realisations in one batch, entanglement entropies &
\texttt{SpinChain}, \texttt{random\_field\_heisenberg}, \texttt{level\_spacing\_ratio},
\texttt{entanglement\_entropy}\\
\repofile{ex8_large_and_open_chains.py} & Lanczos ground state of a sector of dimension $12870$, Krylov evolution
of a domain wall, momentum sectors, particle loss through jumps between sectors & \texttt{xxz\_sparse},
\texttt{krylov.ground\_state}, \texttt{krylov.evolve}, \texttt{xxz\_momentum}\\
\repofile{ex9_qaoa_xy_mixer.py} & XY mixer on a ring of six qubits with amplitude damping, dephasing and
post-selection on the Hamming weight & \texttt{SpinChain.xxz}, \texttt{lowering}, \texttt{algebra.central},
\texttt{condition\_on}\\
\bottomrule
\end{tabular}
\caption{The introductory programs in \texttt{man/article/listings} of the repository.}
\label{tab:examples}
\end{table}

Some of the examples check exact statements. In \texttt{ex3}, the Kraus operators $\sqrt{1-gN}$,
$\sqrt g\,a_1$, $\sqrt g\,a_2$ ($0\le g\le\frac12$) map sector $N$ to $N$ and $N-1$, and the classical part is the
transition matrix of a pure death process (columns: from $N=0,1,2$; rows: to $N=0,1,2$),
$T=\bigl(\begin{smallmatrix}1&g&0\\0&1-g&2g\\0&0&1-2g\end{smallmatrix}\bigr)$; a unital positive map cannot increase
the Michelson contrast, since $\lambda_{\min}1\le A\le\lambda_{\max}1$ implies the same bounds for $\Phi(A)$, while the
non-unital loss channel increases it for most samples. In \texttt{ex7} the ground state moves through the sectors until
the chain is fully polarised above the saturation field $h_s=2J$ (for our sign $+\sum_ih_iS^z_i$ the all-down state
$N=0$; the last crossing is with the one-magnon level at $k=\pi$). In \texttt{ex9} the mixer
$H_M=\frac14\sum_{\langle ij\rangle}(X_iX_j+Y_iY_j)$ keeps the Hamming weight~\cite{hadfield2019}, so the register
carries the algebra $\bigoplus_{N=0}^nM_{\binom nN}$; amplitude damping ($S^-_i$, rate $\gamma_1$) is a jump between
sectors, dephasing ($S^z_i$, rate $\gamma_2$) a jump inside them, and reading out the weight measures a central
projection. With damping only, the post-selected fidelity is exactly one and the acceptance probability is
$e^{-w\gamma_1T}$, because the no-jump evolution is generated by $H-\frac i2\gamma_1N$ and $N$ is central; for the same
reason the post-selected state does not depend on $\gamma_1$ even with dephasing. For $|111000\rangle$, $T=2$ and
$\gamma_1=\gamma_2=0.02$ the program gives the acceptance probability $0.8869$, the fidelity $0.8629$ and $0.9730$ after
post-selection.

\section{Validation against exact results}\label{sec:validation}

Table~\ref{tab:validation} compares Monte Carlo estimates produced with the library (double precision)
with exact values: fourth moments of Haar unitaries and orthogonal matrices; the Weingarten-calculus
value $\mathbb E\,\tau\bigl((AUBU^*)^2\bigr)=-1/(n^2-1)$ for traceless involutions $A=B$ and Haar $U$ in
$M_n(\C)$~\cite{collins2006}; Page's formula for $\langle S(\rho_A)\rangle$, conjectured in~\cite{page1993} and proved
in~\cite{foong1994,sanchezruiz1995}; and the mean purity $(k+r)/(kr+1)$ of random states of $M_k$ induced by an
ancilla of dimension $r\le k$ (so of rank $r$), due to Lubkin~\cite{lubkin1978} (see
also~\cite[Eq.~(3.20)]{zyczkowski2001}). The last column is the deviation in units of the standard
error; for a correct sampler these should look like standard normal deviates. For the circular ensembles the Wigner surmise is only an approximation to
the exact spacing distribution, so no $z$-score is given; Figure~\ref{fig:spacings} shows the
distributions, with $L^1$ distances
\IfFileExists{generated/validation.tex}{\ValLoneOne, \ValLoneTwo\ and \ValLoneFour\ for
$\beta=1,2,4$}{(not generated)} from the surmise; level repulsion $P(s)\sim s^\beta$ is clearly
visible. The spontaneous-emission example agrees with the analytic solution to
\IfFileExists{generated/validation.tex}{$\ValEmissionExp$ (exact exponential) and $\ValEmissionRK$ (RK4)}{(not generated)}.

\begin{table}[h]
\centering\small
\begin{tabular}{llrrrr}
\toprule
quantity & parameters & exact & Monte Carlo & std.\ error & $z$\\
\midrule
\ValidationRows
\bottomrule
\end{tabular}
\caption[]{Monte Carlo estimates versus exact values (for $\langle s^2\rangle$ the Wigner surmise, an approximation,
hence no $z$)\IfFileExists{generated/validation.tex}{; samples: \ValSHaar\ per Haar moment, \ValSW\ per Weingarten
row, \ValSce\ matrices per circular ensemble, \ValSPage\ states per Page row, \ValSPur\ per purity row
(\ValMode\ run on \ValDevice, PyTorch \ValTorch)}{}.}
\label{tab:validation}
\end{table}

\begin{figure}[h]
\centering
\fig{spacings.pdf}{\textwidth}\\[2mm]
\fig{page.pdf}{0.45\textwidth}
\caption[]{Top: nearest-neighbour spacings of COE, CUE and CSE eigenphases generated by the library\IfFileExists{generated/validation.tex}{ ($n=\ValNce$, \ValSce\ matrices per ensemble)}{},
against the Wigner surmise. Bottom: mean entanglement entropy of random pure states against Page's
formula\IfFileExists{generated/validation.tex}{ (\ValSPage\ states per point)}{} (error bars: two standard errors).}
\label{fig:spacings}
\end{figure}

\paragraph{Trace functionals.} The three traces are checked on random non-Hermitian operators in
$M_2\oplus M_3\oplus M_5$ (Table~\ref{tab:traces}): normalisation on the identity ($\sum_ck_c$, $C$ and $1$),
the tracial property $\operatorname{tr}(AB)=\operatorname{tr}(BA)$, invariance under unitaries of $\M$, and
positivity $\operatorname{tr}(A^*A)>0$ (faithfulness); the last column is the smallest ratio
$\operatorname{tr}(A^*A)/\|A\|_F^2$ over the samples, which is bounded below by $\min_cw_c$. The relation $\Tr_{\mathrm{norm}}=C\,\tau_{\mathrm{vN}}$ holds to
\IfFileExists{generated/validation.tex}{$\ValTraceNormTau$}{round-off}. All three return values
in the number field of the algebra (complex traces of non-Hermitian operators).

\begin{table}[h]
\centering\small
\begin{tabular}{lcccc}
\toprule
trace & error of $\operatorname{tr}(1)$ & $\max|\operatorname{tr}(AB)-\operatorname{tr}(BA)|$ &
$\max|\operatorname{tr}(UAU^*)-\operatorname{tr}(A)|$ & $\min\operatorname{tr}(A^*A)/\|A\|_F^2$\\
\midrule
\ValTraceRows
\bottomrule
\end{tabular}
\caption[]{Identities of the three trace functionals, 2000 random operators in double precision\IfFileExists{generated/validation.tex}{; \ValMode\ run on \ValDevice}{}.}
\label{tab:traces}
\end{table}

\section{From sampling to optimisation}\label{sec:optimisation}

The first version of this work~\cite{v1}, and before it~\cite[Sect.~5, Figs.~1--3]{abed2024} for a single block
$M_n(\mathbb R)$ with $n\le5$, studied three trace functionals by Monte Carlo:
\begin{align*}
  \text{Exp.~1:}\quad & z=|\Tr(XUY)|-\Tr(XY), && X,Y\ge0,\ U\in\M\ \text{unitary},\\
  \text{Exp.~2:}\quad & z=\Tr|XY|-\Tr(|X|\,|Y|), && X=UX_0,\ Y=VY_0,\ X_0,Y_0\ge0,\ U,V\in\M\ \text{unitary},\\
  \text{Exp.~3:}\quad & z=\Tr(Y|X|Y)-\Tr|YXY|, && X=X^*,\ Y\ge0,
\end{align*}
with $\Tr=\Tr_{\mathrm{blunt}}$, operators normalised to $\lambda_{\max}=1$, and the contrasts of the
operators as parameters (for the self-adjoint $X$ of Exp.~3 and the factors of Exp.~2, the contrasts of $|X|$, $|Y|$). Two observations motivate a different approach.

\paragraph{Sampling misses the extremes.} For fixed $X,Y$ the supremum in Exp.~1 is known exactly,
$\sup_U z=\|YX\|_1-\Tr(XY)>0$. Table~\ref{tab:search} and Figure~\ref{fig:search} show that the best of
\IfFileExists{generated/search.tex}{\SearchSamples}{many} Haar samples falls short of the supremum for every $k$, misses
the violation for all pairs from $k=4$ on, and moves further away from it as $k$ grows, while batched gradient ascent
(\IfFileExists{generated/search.tex}{\SearchStarts\ starts, \SearchSteps\ steps}{multistart}) finds it
to a relative error of at most \IfFileExists{generated/search.tex}{$\SearchMaxErr$}{about $10^{-5}$} (Table~\ref{tab:search}). This is concentration of measure: the unitaries close to
the maximiser form a set whose Haar measure decays rapidly with the dimension. This search is done in the complex algebra $M_k(\C)$ ($C=1$, eigenvalues of $X,Y$ uniform in $[0.2,1.2]$, double
precision); the same concentration effect explains the absence of violations in the (real) Monte Carlo samples
for $k=16$ reported in~\cite{v1}: it is a property of the sampling, not of the inequality.

\begin{table}[h]
\centering\small
\begin{tabular}{rcccrccrr}
\toprule
$k$ & \multicolumn{2}{c}{violation found} & $\sup z/\Tr XY$ & best sample & \multicolumn{2}{c}{ascent rel.\ error}
& \multicolumn{2}{c}{time [s]}\\
& sampling & ascent & & $z/\Tr XY$ & median & max & sampl. & ascent\\
\midrule
\SearchRows
\bottomrule
\end{tabular}
\caption[]{Exp.~1 for fixed random $X,Y$ in $M_k(\C)$ (eigenvalues uniform in $[0.2,1.2]$, double precision): Haar
sampling versus gradient ascent (medians over
\IfFileExists{generated/search.tex}{\SearchPairs}{several} pairs\IfFileExists{generated/search.tex}{;
\SearchMode\ run on \SearchDevice}{}).}
\label{tab:search}
\end{table}

\begin{figure}[h]
\centering\fig{search.pdf}{0.55\textwidth}
\caption{Best $z$ found by sampling and by gradient ascent, relative to $\Tr(XY)$, against the exact
supremum (symmetric log scale).}
\label{fig:search}
\end{figure}

\paragraph{Constrained optimisation.} The natural objects are the envelopes
\[
  \overline z(\delta)=\sup\{z:\Delta(X)=\delta\},\qquad \underline z(\delta)=\inf\{z:\Delta(X)=\delta\},
\]
(and the same with $\Delta(Y)$, or with both contrasts prescribed), the supremum being taken over all
remaining variables. The module \texttt{optimize} parametrises the constraint set exactly: the spectrum
is $\lambda_{\max}=1$, $\lambda_{\min}=(1-\delta)/(1+\delta)$ and the other eigenvalues in between, as
$\lambda_{\min}+(1-\lambda_{\min})(1+\sin\theta)/2$ (unlike a sigmoid, this reaches both ends at finite $\theta$; extremal
spectra often have all eigenvalues at the ends); the two eigenvalue positions that carry $\lambda_{\max}$ and $\lambda_{\min}$ are drawn at
random per start, in the same sector or in different ones, so that the starts together cover a dense subset of
the constraint set, the eigenvectors are $U=U_0e^{G-G^*}$ with a Haar-random $U_0$ per start and a chart
that is re-centred after every round. For Exp.~3 the signs of the eigenvalues of $X$ are discrete
variables: the number of negative eigenvalues is drawn per sector (for half of the starts as one fraction common to
all sectors), and between the gradient rounds single signs are flipped where this improves $z$. Every iterate is an
admissible point, so the best value seen per start is kept. Every grid point and every start is one element of the batch, so each envelope (supremum
or infimum along one contrast) is computed in one batched optimisation. All envelopes are computed in real algebras
(real symmetric $X,Y$ and orthogonal $U$), the setting of the Monte Carlo study of~\cite{v1}.

Figures~\ref{fig:bounds} and~\ref{fig:bounds2d} and Table~\ref{tab:bounds} show the result. The search is local,
so the computed envelopes are checked in three independent ways. (i)~Where an
envelope is known in closed form (below), it is compared with the formula. (ii)~Since $z$ is additive over
orthogonal blocks (in Exp.~1 the trace $|\Tr(XUY)|$ is only subadditive, but replacing $U$ by $-U$ on a block aligns
the phases of the block traces, so the bound below holds) and the contrast of a direct sum is the largest contrast of
its summands, a direct sum of
$kC/2$ copies of a solution for $k=2$, $C=1$ is an admissible configuration: with $s(\delta)$ the upper envelope
for $k=2$, $C=1$,
\[
  \overline z_{k,C}(\delta)\ \ge\ s(\delta)+\bigl(\tfrac{kC}2-1\bigr)\max_{\delta'\le\delta}s(\delta')\qquad(k\ \text{even}),
\]
and the same with $\min$ for the lower envelopes; the computed envelopes must not lie below this value (above, for
the infima). (iii)~Each Monte Carlo sample of~\cite{v1} and each of
\IfFileExists{generated/bounds.tex}{\BoundsStChecked}{several thousand} \emph{structured samples} is checked against
both envelopes, along $\Delta(X)$ and along $\Delta(Y)$; samples that fall outside the linear interpolation between
grid points are re-optimised at their exact contrast. The structured samples are computed with plain NumPy: all
spectra take only the two values $1$ and $\lambda_{\min}$, and the eigenbases and $U$ are random signed permutations
composed with direct sums of $2\times2$ rotations. Unlike Haar samples they come close to the extremal
configurations in every dimension, so they test the optimiser where Monte Carlo cannot.
\IfFileExists{generated/bounds.tex}{Of the \BoundsMcChecked\ Monte Carlo samples with both contrasts in the grid
(\BoundsMcDropped\ have a contrast above $0.999$), \BoundsMcOutside\ remain outside after re-optimisation, and
\BoundsStOutside\ of the structured samples; the envelopes fall short of the direct sums~(ii) by at most
\BoundsShortfall\% of the largest value of the curve.}{} For a single sector of size $k=2$ the Haar samples come close
to the envelopes; with two sectors they already miss a substantial part of the range (Table~\ref{tab:bounds}), and
for larger $k$ or $C$ they stay in a narrow band far from both envelopes (Figure~\ref{fig:bounds}, bottom row; Table~\ref{tab:search}):
the extremes are attained at degenerate spectra with aligned eigenbases, a set that Haar sampling essentially never
reaches.

Several envelopes can be checked analytically. In Exp.~3, $z\ge0$ always: the inequality $\Tr|YXY|\le\Tr(Y|X|Y)$ is
standard (it follows from $-Y|X|Y\le YXY\le Y|X|Y$, since $-K\le A\le K$ implies $\|A\|_1\le\Tr K$), and what we
need in addition is the case of equality:

\begin{proposition}
For $X=X^*$ and $Y\ge0$, equality holds in $\Tr|YXY|\le\Tr(Y|X|Y)$ if and only if $YX_+Y^2X_-Y=0$, in particular
if $Y$ commutes with $X$.
\end{proposition}
\begin{proof}
With $X=X_+-X_-$, $|X|=X_++X_-$, the triangle inequality for the trace norm gives the inequality,
$\|YX_+Y-YX_-Y\|_1\le\Tr(YX_+Y)+\Tr(YX_-Y)$; for positive $A,B$, $\|A-B\|_1=\Tr A+\Tr B$ if and only if
$AB=0$. Indeed, with $P$ the projection onto the positive part of $A-B$,
$\|A-B\|_1=\Tr\bigl((2P-1)(A-B)\bigr)=\Tr A+\Tr B-2\Tr\bigl((1-P)A\bigr)-2\Tr(PB)$, and both traces subtracted are
non-negative; equality forces $(1-P)A=0=PB$, so $A=AP$, $B=(1-P)B$ and $AB=AP(1-P)B=0$. Conversely, if $AB=0$ the
supports of $A$ and $B$ are orthogonal and $|A-B|=A+B$. Apply this to $A=YX_+Y$, $B=YX_-Y$. If $XY=YX$, then $YX_+Y^2X_-Y=Y^4X_+X_-=0$.
\end{proof}
\noindent The computed lower envelope of Exp.~3 is zero to within \IfFileExists{generated/bounds.tex}{$\BoundsExpThreeInf$}{round-off}. In Exp.~1, the lower envelope
for prescribed $\Delta(X)=\delta$ equals $-(kC-2\delta/(1+\delta))$ when all sectors have size $k\ge2$:
$\Tr(XY)\le\Tr X\le kC-1+\lambda_{\min}(X)$ with equality for $Y=1$, while $|\Tr(UYX)|$ can be made zero
by a cyclic permutation of the eigenbasis of $X_c$ in every sector. The optimiser reproduces this curve to a relative deviation of
\IfFileExists{generated/bounds.tex}{\BoundsExactDev}{(not generated)} (along either contrast: in a real algebra
$\Tr(XUY)=\Tr(YU^TX)$, so $X$ and $Y$ exchange roles). In Exp.~2, $|X|=X_0$ and
$|Y|=Y_0$ and the unitary $U$ drops out, so $z=\|X_0VY_0\|_1-\Tr(X_0Y_0)$. Since $V$ and the eigenbases act
sector by sector, von Neumann's trace inequality $\|AB\|_1\le\sum_is_i(A)s_i(B)$, its lower counterpart
$\|AB\|_1\ge\sum_is_i(A)s_{k+1-i}(B)$ and rearrangement give, for given spectra,
\[
  \sup z=-\inf z=\sum_c\sum_i x_{c,i}^\downarrow\bigl(y_{c,i}^\downarrow-y_{c,i}^\uparrow\bigr),
\]
with $x_{c,i}^\downarrow$ the eigenvalues of $X_c$ in decreasing order and $y_{c,i}^\downarrow$, $y_{c,i}^\uparrow$ those of
$Y_c$ in decreasing and increasing order. Pairing $i$ with $k+1-i$,
$\sum_ix_{c,i}^\downarrow(y_{c,i}^\downarrow-y_{c,i}^\uparrow)=\sum_{i\le k/2}(x_{c,i}^\downarrow-x_{c,k+1-i}^\downarrow)
(y_{c,i}^\downarrow-y_{c,k+1-i}^\downarrow)\le\lfloor k/2\rfloor(1-\lambda_{\min}(X))(1-\lambda_{\min}(Y))$, with equality
for two-level spectra. Hence, for sectors of sizes $k_c$, $\sup z=-\inf z=\sum_c\lfloor k_c/2\rfloor(1-\lambda_{\min}(X))(1-\lambda_{\min}(Y))$, and for
$C$ sectors of size $k$, $\Delta(|X|)=\delta$ prescribed and $\Delta(|Y|)\le0.999$
free,
\[
  \sup z=-\inf z=C\Bigl\lfloor\frac k2\Bigr\rfloor\,\frac{2\delta}{1+\delta}\Bigl(1-\frac{0.001}{1.999}\Bigr),
\]
which tends to $C\lfloor k/2\rfloor\frac{2\delta}{1+\delta}$ as $\Delta(|Y|)\to1$. The optimiser reproduces the upper and
lower envelopes of Exp.~2 for all eight configurations, along either contrast, to a relative deviation of
\IfFileExists{generated/bounds.tex}{\BoundsExactDevTwo}{(not generated)}.
For sectors of size $k=2$ the upper envelopes of Exp.~1 and Exp.~3 can be computed exactly. In $M_2(\C)$ the
eigenbases $(e_i)$ of $X$ and $(f_j)$ of $Y$ enter only through $\gamma=|\langle e_1,f_1\rangle|^2\in[0,1]$, since
$(|\langle e_i,f_j\rangle|^2)$ is doubly stochastic; phases play no role, and every $\gamma$ is reached by a real rotation, so
real and complex algebras have the same envelopes.

\begin{proposition}\label{prop:k2}
(a) Let $X,Y\ge0$ in $M_2(\C)$ have eigenvalues $x_1\ge x_2\ge0$ and $y_1\ge y_2\ge0$, $x_1,y_1>0$. For every unitary $U$,
\[
  |\Tr(XUY)|-\Tr(XY)\ \le\ \frac{(x_1-x_2)^2(y_1-y_2)^2}{4(x_1+x_2)(y_1+y_2)} .
\]
(b) Let $X=X^*$ in $M_2(\C)$, let $|X|$ have eigenvalues $x_1\ge x_2\ge0$, $x_1>0$, and let $Y$ be as in (a). Then
\[
  \Tr(Y|X|Y)-\Tr|YXY|\ \le\ \frac{2x_1x_2}{x_1+x_2}\,(y_1-y_2)^2 .
\]
Let $S=x_1y_2+x_2y_1$, $q=(x_1+x_2)(y_1+y_2)$, and let $e_1$ be an eigenvector of $X$ for its eigenvalue of modulus $x_1$,
$f_1$ one of $Y$ for $y_1$. The bound (a) is attained iff $|\Tr(XUY)|=\|YX\|_1$ and, if $x_1>x_2$ and $y_1>y_2$,
$\gamma=\frac14+\frac S{2q}$ (otherwise the bound is $0$). If $x_2>0$ and $y_1>y_2$, the bound (b) is attained iff $X$ is
indefinite and $\gamma=\frac Sq$ (otherwise it is $0$, and so is $z$). Both bounds are attained by real $X$, $Y$, $U$.\\
(c) Let $\M=\bigoplus_{c=1}^CM_2(\C)$, $X$ and $Y$ (in Exp.~3: $|X|$ and $Y$) normalised to $\lambda_{\max}=1$ with contrasts
$\delta<1$ and $\delta'\le1$, and $\ell=\frac{1-\delta}{1+\delta}$, $\ell'=\frac{1-\delta'}{1+\delta'}$. Then
\begin{align*}
  \text{Exp.~1:}\quad & \sup z=\frac{C\,\delta^2\delta'^2}{(1+\delta)(1+\delta')}, && \inf z=-(2C-1+\ell\ell'),\\
  \text{Exp.~3:}\quad & \sup z=(C-\delta)\Bigl(\frac{2\delta'}{1+\delta'}\Bigr)^2, && \inf z=0,
\end{align*}
and all four values are attained.
\end{proposition}
\begin{proof}
Let $p=(x_1-x_2)(y_1-y_2)$, so that $q-2S=p$. (a) $|\Tr(XUY)|=|\Tr(YXU)|\le\|YX\|_1$, with equality for $U=W^*$,
$YX=W|YX|$. Expanding in the eigenbases, $\Tr(XY)=\sum_{ij}x_iy_j|\langle e_i,f_j\rangle|^2=S+\gamma p$; the singular values of
$YX$ satisfy $s_1^2+s_2^2=\Tr(X^2Y^2)$ and $s_1s_2=|\det YX|=x_1x_2y_1y_2$, so $\|YX\|_1^2=\Tr(X^2Y^2)+2x_1x_2y_1y_2=S^2+\gamma pq$.
With $\sqrt r\le r/q+q/4$ (equality iff $\sqrt r=q/2$),
$z\le\|YX\|_1-\Tr(XY)\le(S^2+\gamma pq)/q+q/4-S-\gamma p=(q/2-S)^2/q=p^2/(4q)$, with equality iff in addition
$\gamma pq=q^2/4-S^2=\frac p2(\frac q2+S)$.
(b) $z$ is invariant under $X\mapsto-X$ and vanishes for semidefinite $X$. Otherwise let $Xe_1=x_1e_1$, $Xe_2=-x_2e_2$,
$u=\langle e_1,Y^2e_1\rangle=\gamma y_1^2+(1-\gamma)y_2^2$ and $s=y_1^2+y_2^2$. Then $\Tr(Y|X|Y)=\Tr(|X|Y^2)=A:=(x_1-x_2)u+x_2s$, and
since $\det(YXY)=-x_1x_2y_1^2y_2^2\le0$, $\Tr|YXY|=\sqrt{\Tr(YXY)^2-4\det(YXY)}=\sqrt{L^2+K^2}$ with $L=\Tr(XY^2)=(x_1+x_2)u-x_2s$
and $K=2\sqrt{x_1x_2}\,y_1y_2$. For the unit vector $(\alpha,\beta)=(x_1-x_2,\,2\sqrt{x_1x_2})/(x_1+x_2)$, the Cauchy--Schwarz
inequality gives $\sqrt{L^2+K^2}\ge\alpha L+\beta K$, and $A-\alpha L-\beta K=\frac{2x_1x_2}{x_1+x_2}(y_1-y_2)^2$ for every $u$;
equality holds iff $\beta L=\alpha K$, i.e.\ $u-y_2^2=(y_1-y_2)S/(x_1+x_2)$, i.e.\ $\gamma=S/q$. Real maximisers: diagonal $X$, $Y$ rotated by
$\arccos\sqrt\gamma$, and $U$ from a real singular value decomposition of $YX$.
(c) In Exp.~1, $|\Tr(XUY)|\le\sum_c\|Y_cX_c\|_1$, and by (a) block $c$ contributes at most
$\frac{x_{c,1}y_{c,1}}4\,g(\frac{x_{c,2}}{x_{c,1}})\,g(\frac{y_{c,2}}{y_{c,1}})$ with $g(r)=\frac{(1-r)^2}{1+r}$ decreasing; as
$x_{c,1},y_{c,1}\le1$, $x_{c,2}/x_{c,1}\ge\ell$, $y_{c,2}/y_{c,1}\ge\ell'$ and $g(\ell)=2\delta^2/(1+\delta)$, every block contributes at
most $\delta^2\delta'^2/((1+\delta)(1+\delta'))$; equality for $X_c=\mathrm{diag}(1,\ell)$, $Y_c$ with spectrum $\{1,\ell'\}$ at the angle
of (a), and $U_c$ the polar unitaries, which make all block traces positive. For the infimum, $U_c$ with
$\Tr(X_cU_cY_c)=0$ (a cyclic shift in the singular bases) gives $z=-\Tr(XY)$, and $\Tr(XY)\le\sum_c(x_{c,1}y_{c,1}+x_{c,2}y_{c,2})\le
2C-1+\ell\ell'$ (if $\ell$ and $\ell'$ lie in different blocks the bound is $2C-2+\ell+\ell'$, smaller by $(1-\ell)(1-\ell')$), with
equality for $X=\mathrm{diag}(1,\ell)\oplus1$, $Y=\mathrm{diag}(1,\ell')\oplus1$. In Exp.~3, $z=\sum_cz_c$ and by (b)
$z_c\le h_c(1-\ell')^2$ with $h_c=\frac{2x_{c,1}x_{c,2}}{x_{c,1}+x_{c,2}}\le x_{c,1}\le1$; in the block that holds the eigenvalue $\ell$
of $|X|$, $h_c=2\ell/(1+\ell/x_{c,1})\le2\ell/(1+\ell)=1-\delta$. Equality for $X=\mathrm{diag}(1,-\ell)\oplus\mathrm{diag}(1,-1)^{\oplus(C-1)}$
and $Y_c$ with spectrum $\{1,\ell'\}$ at the angle of (b); $\inf z=0$ is Proposition~1.
\end{proof}

\noindent Both suprema increase with $\delta'$. With the free contrast capped at $D$, the upper envelopes along
$\Delta(X)=\delta$ are therefore $C\delta^2D^2/((1+\delta)(1+D))$ and $(C-\delta)\,(2D/(1+D))^2$; for $C=1$ and $D=0.999$ these
are the closed forms that we first found numerically. As $D\to1$ they increase
to $C\delta^2/(2(1+\delta))$ and $C-\delta$, the values at $\Delta(Y)=1$ (a singular $Y$), which no $Y$ with $\Delta(Y)<1$
attains (for $\delta>0$ in Exp.~1). Along $\Delta(Y)=\delta'$ the upper envelopes are $C\delta'^2D^2/((1+\delta')(1+D))$ and,
attained at $|X|=1$, $C\,(2\delta'/(1+\delta'))^2$ for every $D$. For $k=2$ the direct sums~(ii) are thus optimal for every
$C$.

\paragraph{Sectors of any size.} The upper envelopes of Exp.~1 and Exp.~3 are known for every sector size. Let
$N=\sum_c\lfloor k_c/2\rfloor$ be the largest number of mutually orthogonal two-dimensional subspaces that fit into the
sectors.

\begin{theorem}\label{thm:allk}
Let $\M=\bigoplus_{c=1}^CM_{k_c}(\C)$ or $\bigoplus_{c=1}^CM_{k_c}(\mathbb R)$ with $N\ge1$, and let $X$ and $Y$ (in Exp.~3:
$|X|$ and $Y$) be normalised to $\lambda_{\max}=1$, with prescribed contrasts $\delta\in[0,1)$ and $\delta'\in[0,1]$. Then
\[
  \text{Exp.~1:}\ \ \sup z=\frac{N\,\delta^2\delta'^2}{(1+\delta)(1+\delta')},\qquad
  \text{Exp.~3:}\ \ \sup z=\begin{cases}(N-\delta)\bigl(\frac{2\delta'}{1+\delta'}\bigr)^2&\text{if every }k_c\text{ is even},\\[2pt]
  N\bigl(\frac{2\delta'}{1+\delta'}\bigr)^2&\text{if some }k_c\text{ is odd},\end{cases}
\]
and both suprema are attained by real direct sums of the $2\times2$ maximisers of Proposition~\ref{prop:k2} and
$1\times1$ blocks.
\end{theorem}

\noindent For $C$ sectors of size $k$ and the other contrast free in $[0,D]$, the upper envelopes are therefore
$C\lfloor k/2\rfloor\delta^2D^2/((1+\delta)(1+D))$ in Exp.~1, along either contrast; in Exp.~3 along $\Delta(|X|)=\delta$
\[
  \overline z(\delta)=\Bigl(\frac{Ck}2-\delta\Bigr)\Bigl(\frac{2D}{1+D}\Bigr)^2\quad(k\text{ even}),\qquad
  \overline z(\delta)=C\Bigl\lfloor\frac k2\Bigr\rfloor\Bigl(\frac{2D}{1+D}\Bigr)^2\quad(k\text{ odd}),
\]
and along $\Delta(Y)=\delta'$, attained at $|X|=1$, $C\lfloor k/2\rfloor(2\delta'/(1+\delta'))^2$. For even $k$ these are
the values of the direct sums~(ii): the direct sums are optimal for even $k$, in both functionals. For odd $k$ the
direct sums with a one-dimensional summand are optimal in Exp.~1 but not in Exp.~3: there the eigenvalue
$\lambda_{\min}$ of $|X|$ can be placed on a one-dimensional summand, where it costs nothing, and the envelope along
$\Delta(|X|)$ does not depend on $\delta$ (for $k=3$, $C=1$: $(2D/(1+D))^2>(1-\delta)(2D/(1+D))^2$). The
proof (Appendix~\ref{app:envelopes}) reduces both functionals to two-level spectra and, with the two-projection lemma~\cite{halmos1969}, to
$2\times2$ blocks, which Proposition~\ref{prop:k2} handles, except in one case: in Exp.~3 the eigenvector of $|X|$ for
$\lambda_{\min}$ may be spread over all two-dimensional blocks of its sector, so that $X$ is not block-diagonal. This case
is settled by the following bound, which is sharp; its proof linearises the problem.

\begin{proposition}\label{prop:mblocks}
Let $X=X_+\oplus(-1)$ on $\C^m\oplus\C^m$, where $0<X_+\le1$ has eigenvalues $x_1,\dots,x_m$; let
$c_\theta=\mathrm{diag}(\cos\theta_j)$, $s_\theta=\mathrm{diag}(\sin\theta_j)$, $Q$ the projection onto the range of
$\binom{c_\theta}{s_\theta}:\C^m\to\C^{2m}$, $b\in[0,1)$ and $Y=b+(1-b)Q$. Then
\[
  \Tr(Y|X|Y)-\Tr|YXY|\ \le\ (1-b)^2\sum_{i=1}^m\frac{2x_i}{1+x_i},
\]
with equality for $X_+=\mathrm{diag}(x_1,\dots,x_m)$ and $\cos^2\theta_i=(1+bx_i)/((1+x_i)(1+b))$.
\end{proposition}
\begin{proof}
Let $\gamma=1-b^2$, $K=1+X_+$, $A=c_\theta^2$ (so $s_\theta^2=1-A$). In the decomposition $\C^m\oplus\C^m$,
$Y^2=b^2+\gamma Q$ has the blocks $b^2+\gamma A$, $\gamma c_\theta s_\theta$, $\gamma c_\theta s_\theta$, $b^2+\gamma s_\theta^2$, all diagonal.
For $\lambda>0$ the lower right block of $\lambda-XY^2$ is the invertible diagonal $E=\lambda+1-\gamma A$, and as
$(b^2+\gamma A)E-\gamma^2As_\theta^2=\lambda(b^2+\gamma A)+b^2$, the Schur complement formula gives
\[
  \det(\lambda-XY^2)=\det\bigl(\lambda E-X_+(\lambda(b^2+\gamma A)+b^2)\bigr)=\det\bigl((\lambda+1)(\lambda-b^2X_+)-\lambda\gamma KA\bigr)
  =p(\lambda),
\]
$p(\lambda)=\det\bigl((\lambda+1)(\lambda-b^2X_+)-\lambda\gamma K^{1/2}AK^{1/2}\bigr)$. The Schur complement of the lower right
block $\lambda$ gives $\det(\lambda-\mathcal L)=p(\lambda)$ as well, for the self-adjoint matrix
\[
  \mathcal L=\begin{pmatrix}\gamma K^{1/2}AK^{1/2}+b^2X_+-1 & bX_+^{1/2}\\ bX_+^{1/2} & 0\end{pmatrix},
\]
which is affine in $A$. Since $\det(\lambda-YXY)=\det(\lambda-XY^2)$, the polynomials agree and $YXY$ and $\mathcal L$ have
the same eigenvalues; hence $\Tr|YXY|=\|\mathcal L\|_1\ge\Tr(T\mathcal L)$ for every self-adjoint $T$ with $\|T\|\le1$.
Take
\[
  T=\begin{pmatrix}-\Delta & G\\ G & \Delta\end{pmatrix},\qquad \Delta=\frac{1-X_+}{1+X_+},\quad G=\frac{2X_+^{1/2}}{1+X_+},
\]
so that $T^2=1$ ($\Delta$ and $G$ commute and $\Delta^2+G^2=1$). With $\Tr(Y|X|Y)=\Tr(|X|Y^2)=b^2\Tr X_++\gamma\Tr(X_+A)+mb^2+\gamma\Tr(1-A)$
and $K^{1/2}\Delta K^{1/2}=1-X_+$, the terms containing $A$ cancel in $\Tr(|X|Y^2)-\Tr(T\mathcal L)$, and
\[
  \Tr(Y|X|Y)-\Tr|YXY|\le\sum_{i=1}^m\Bigl(b^2x_i+1+\frac{1-x_i}{1+x_i}(b^2x_i-1)-\frac{4bx_i}{1+x_i}\Bigr)=(1-b)^2\sum_{i=1}^m\frac{2x_i}{1+x_i}.
\]
For diagonal $X_+$ both operators are direct sums of $2\times2$ blocks, on which $X=\mathrm{diag}(x_i,-1)$, and
Proposition~\ref{prop:k2}(b) gives the value $(1-b)^2\,2x_i/(1+x_i)$ at the stated angle (there $\gamma=\sin^2\theta_i$).
\end{proof}

\noindent Since $\mathcal L$ is affine in $A=PQP$ ($P$ the projection onto the first summand), $z$ is a concave
function of $A$, and $T$ is the sign of $\mathcal L$ at the maximiser. For fixed angles, however, the maximum over the
direction of the eigenvector for $\lambda_{\min}(|X|)$ is in general not attained in a single block; only the joint
maximum is. Before this proof we had certified the cases $k=4$ and $k=6$ ($C=1$) numerically, by a rigorous branch and
bound over the angles of the reduced model (\texttt{scripts/certify\_envelopes.py}; gaps $10^{-7}$ and $10^{-5}$,
including floating-point rounding); it agrees with Theorem~\ref{thm:allk} and serves as an independent check. All
upper envelopes of the three functionals studied here, along either contrast, are now proved, and so are the lower
envelopes with the other contrast free; we have not studied the lower envelope of Exp.~1 with both contrasts
prescribed for $k\ge3$. The optimiser reproduces the upper envelopes of Theorem~\ref{thm:allk} for all eight
configurations, along either contrast, to a relative deviation of \BoundsSupDevOne\ (Exp.~1) and \BoundsSupDevThree\
(Exp.~3). For this, half of the starts begin at the structure of the maximisers
(\texttt{init} and \texttt{spectrum\_init} of \texttt{optimize}): $X$ diagonal, $Y$ rotated in pairs of eigenvectors of
$X$, two-level spectra (in Exp.~3 one positive and one negative eigenvalue of $X$ per pair), and $U$ rotating the same
pairs. The last point matters in a real algebra: a chart $U_0e^{G-G^*}$ never leaves the connected component of
$U_0$, and the maximiser of Exp.~1 has determinant $+1$ in every sector, which a Haar-random $U_0$ has with probability
$2^{-C}$. With Haar-random starts only, the run of the first version of this section fell short of
Theorem~\ref{thm:allk} by up to $35\%$ (Exp.~3) and $13\%$ (Exp.~1) for $k=2$, $C=32$; with the structured starts, a CPU
run with 12 starts per grid point and 240 steps reaches it to a relative $10^{-3}$ (Exp.~3) and $10^{-5}$ (Exp.~1)
for $k=2$, $C=32$ and for $k=16$, $C=1,2,16$ (the remaining gap in Exp.~3 comes from the kink of $z$ at the maximiser,
which is smoothed only on the scale $\lambda_{\min}(Y)=0.001/1.999$; it decreases with the step size).

\paragraph{Relation to Kantorovich-type inequalities.} Exp.~1 compares $|\Tr(XUY)|$, whose supremum over $U$ is
$\|YX\|_1$, with $\Tr(XY)$ for positive $X,Y$ whose spectra are confined by their contrasts. It thus belongs to the
family of Kantorovich-type inequalities (reverse Cauchy--Schwarz and condition-number inequalities), in which the ratio
$\lambda_{\max}/\lambda_{\min}$ of an operator, here $(1+\delta)/(1-\delta)$ for the contrast $\delta$, bounds how
far an inequality can be reversed~\cite{greub1959,furuta2005}. The sharp envelopes with prescribed contrasts of
Proposition~\ref{prop:k2} and Theorem~\ref{thm:allk}, and their additive form over the sectors (the factor $N$), are
to our knowledge new.

\begin{table}[h]
\centering\small
\begin{tabular}{rrrrrcr}
\toprule
Exp. & $k$ & $C$ & $\inf z$ & $\sup z$ & Monte Carlo range of $z$ & outside\\
\midrule
\BoundsRows
\bottomrule
\end{tabular}
\caption[]{Global extremes of $z$ over both envelopes and contrasts in $[0,0.999]$ found by constrained optimisation (real algebras,
orthogonal $U$), the range of the Monte Carlo samples with both contrasts in that interval, and the number of
samples outside the envelopes (along either contrast) after re-optimisation
\IfFileExists{generated/bounds.tex}{(\BoundsStarts\ starts per grid point, \BoundsSteps\ steps; \BoundsMode\
run on \BoundsDevice)}{}; a sample counts as outside if it exceeds the envelope by more than $10^{-6}+10^{-3}\max|z|$;
``--'': no Monte Carlo data of~\cite{v1} for this configuration. Monte Carlo values of $z$ slightly below
$0$ in Exp.~3 (where $z\ge0$ is proved) are single-precision round-off of the samples of~\cite{v1}.}
\label{tab:bounds}
\end{table}

\begin{figure}[p]
\centering
\fig{bounds_exp1_k2_C1.pdf}{0.64\textwidth}\\
\fig{bounds_exp2_k2_C1.pdf}{0.64\textwidth}\\
\fig{bounds_exp3_k2_C1.pdf}{0.64\textwidth}
\IfFileExists{figures/bounds_exp2_k16_C1.pdf}{\\\fig{bounds_exp2_k16_C1.pdf}{0.64\textwidth}}{}
\caption{Envelopes from constrained optimisation versus sampling for the three trace functionals (rows 1--3: $k=2$,
$C=1$; row 4: Exp.~2 for $k=16$, $C=1$).
Horizontal axis: Michelson contrast of the first operator (left) or of the second operator (right); vertical axis:
$z$ for operators normalised to $\lambda_{\max}=1$. Grey: the Monte Carlo samples of~\cite{v1} (Haar eigenbases,
uniform spectra); orange: structured samples (two-level spectra, permutation and rotation bases; computed with NumPy,
independently of the library); red and blue: the upper and lower envelopes found by constrained optimisation, the
other contrast being free; dashed: envelopes known in closed form; dotted: direct sums of the $k=2$, $C=1$ solution
(where they apply). No sample lies outside the envelopes beyond the counts in Table~\ref{tab:bounds}. In the
smallest case the Haar samples fill the region between the envelopes; for $k=16$ (row 4) they stay in a narrow band
around $z=0$, while the structured samples and the exact envelopes $\pm8\cdot\frac{2\delta}{1+\delta}(1-\frac{0.001}{1.999})$
show how far the range really extends. In Exp.~1 the region $z>0$ (violation of $|\Tr(XUY)|\le\Tr(XY)$)
extends to the value of $\sup z$ in Table~\ref{tab:bounds}; in Exp.~3 the lower envelope is the proven bound $z\ge0$.}
\label{fig:bounds}
\end{figure}

\begin{figure}[h]
\centering
\fig{bounds2d_exp1_k2_C1.pdf}{0.8\textwidth}
\caption{Exp.~1, $k=2$, $C=1$, both contrasts prescribed: $\sup z$ (left) and $\inf z$ (right) on a grid of
$\Delta(Y)$ (horizontal) and $\Delta(X)$ (vertical). The violation $\sup z>0$ needs both operators to be far from
scalar and grows with both contrasts.}
\label{fig:bounds2d}
\end{figure}

\subsection{Rediscovering the data-processing thresholds of the Rényi divergences}\label{sec:renyi}

The closed forms above were first found numerically and proved afterwards. A further test of the method, on an
inequality whose range of validity is known exactly, is to ask it where that range ends, without telling it. For states $\rho,\sigma$ on $\M$ (densities with
respect to $\Tr_{\mathrm{blunt}}$) and $\alpha\in(0,1)\cup(1,\infty)$ the Petz~\cite{petz1986} and the
sandwiched~\cite{mullerlennert2013,wilde2014} Rényi divergences are
\[
  D_\alpha(\rho\|\sigma)=\frac{1}{\alpha-1}\log\Tr\bigl(\rho^\alpha\sigma^{1-\alpha}\bigr),\qquad
  \widetilde D_\alpha(\rho\|\sigma)=\frac{1}{\alpha-1}\log\Tr\bigl(\sigma^{\frac{1-\alpha}{2\alpha}}\,\rho\,
  \sigma^{\frac{1-\alpha}{2\alpha}}\bigr)^{\alpha}.
\]
Both tend to the Umegaki relative entropy $D(\rho\|\sigma)=\Tr\rho(\log\rho-\log\sigma)$ as $\alpha\to1$;
$\widetilde D_\alpha\le D_\alpha$, $\widetilde D_{1/2}=-\log F(\rho,\sigma)$ with the fidelity
$F=\|\sqrt\rho\sqrt\sigma\|_1^2$, and $\widetilde D_\alpha\to D_{\max}$ as $\alpha\to\infty$. The module
\texttt{divergences} evaluates them, batched and differentiable (Daleckii--Krein gradients), with exact support
conventions (powers are taken on the numerical support, and the divergence is $+\infty$ for $\alpha>1$ when
$\operatorname{supp}\rho\not\le\operatorname{supp}\sigma$), as logarithms of sums over sectors, so that large
$\alpha$ neither overflows nor underflows. On $\M$ they split into a classical part and the conditional divergences
of the sectors, and they do not depend on the trace used to represent the states (the central weights cancel); the
tests compare them with dense references and check these properties.

The data-processing inequality $D(\Phi(\rho)\|\Phi(\sigma))\le D(\rho\|\sigma)$ for all channels $\Phi$ holds for the
Petz divergences for $\alpha\in(0,1)\cup(1,2]$~\cite{petz1986} and for the sandwiched ones for
$\alpha\ge\frac12$; for the latter it was proved for $\alpha\in(1,2]$ in~\cite{mullerlennert2013,wilde2014}, for all
$\alpha>1$ in~\cite{beigi2013}, and for all $\alpha\ge\frac12$ in~\cite{frank2013}. Outside these ranges it fails in
general: for the Petz divergences with $\alpha>2$ the trace functional is neither convex nor concave in
$\sigma$~\cite{mullerlennert2013}, and for the sandwiched ones with $\alpha<\frac12$ numerical counterexamples were
reported in~\cite{mullerlennert2013}, and by~\cite{berta2017} it fails, for a suitable measurement, for every
pair of non-commuting states. Both ranges are the cases $z=1$ and $z=\alpha$ of the data-processing range of the
$\alpha$-$z$ R\'enyi divergences, which Zhang determined completely~\cite{zhang2020}.
For the relative entropy, equality characterises the channels that can be reversed on $\rho$ and $\sigma$, and the
reversal is the Petz recovery map (for subalgebras and conditional expectations~\cite{petz1986suff}, for channels~\cite{petz1988chan})
$R_{\sigma,\Phi}(X)=\sigma^{1/2}\Phi^*\bigl(\Phi(\sigma)^{-1/2}X\Phi(\sigma)^{-1/2}\bigr)\sigma^{1/2}$, which satisfies
$R_{\sigma,\Phi}(\Phi(\sigma))=\sigma$. The module provides it for channels within and between sectors, together with
its rotated versions and the universal recovery map $R_{\mathrm{univ}}$ of~\cite{junge2018}, for which
$D(\rho\|\sigma)-D(\Phi(\rho)\|\Phi(\sigma))\ge-\log F(\rho,R_{\mathrm{univ}}(\Phi(\rho)))$ whenever
$\operatorname{supp}\rho\le\operatorname{supp}\sigma$ (\cite[Remark~2.2]{junge2018}, written there with the root
fidelity $\sqrt F$ and the factor $-2\log$); the tests check the
inequalities inside their ranges on random states and channels, Petz's equality condition and the recovery bound.

\texttt{scripts/renyi\_threshold.py} maximises the violation
\[
  v(\alpha)=D(\Phi(\rho)\|\Phi(\sigma))-D(\rho\|\sigma)
\]
for both families on a grid of \IfFileExists{generated/renyi_threshold.tex}{\RenyiNumAlphas}{31} values
$\alpha\in[0.2,3]$, over states $\rho,\sigma$ on $M_d$ and over two kinds of channels: the complete dephasing (the
conditional expectation onto the diagonal subalgebra) and all channels with a $d$-dimensional environment, given by
Stinespring isometries $V=Z(Z^*Z)^{-1/2}$ with a free complex $Z$. All values of $\alpha$ and all starts form one
batch ($\alpha$ is given per sample), optimised with the batched multistart ascent of Section~\ref{sec:optimisation}.
The states are restricted to $\rho,\sigma\ge\varepsilon/d$; this keeps the supremum finite and the sandwiched
divergence resolved in double precision for $\alpha\ge0.2$. The bracket between the last violating and the first
non-violating grid point is then refined by searches at interior points, started from random points and from
perturbations of the best violating point. Every reported violation is attained at an explicit point (two states and
the Kraus operators of a channel), and $v$ is recomputed there independently in mpmath with
\IfFileExists{generated/renyi_threshold.tex}{\RenyiDigits}{50} significant digits, after re-orthonormalising the Kraus
operators so that the point is an exactly trace-preserving channel; it counts as a violation only if the certified
value exceeds a tolerance far above the error of the certification. A certified violation therefore proves that the
inequality fails at that $\alpha$; the absence of a violation at the other end of a bracket is evidence only, since a
local search can miss small violations.
\IfFileExists{generated/renyi_threshold.tex}{The sandwiched inequality is violated (certified) at
$\alpha=\RenyiSandLow$ with $v=\RenyiSandLowV$, and no violation above $\RenyiTol$ was found at
$\alpha=\RenyiSandHigh$; the Petz inequality is violated (certified) at $\alpha=\RenyiPetzHigh$ with
$v=\RenyiPetzHighV$, and no violation was found at $\alpha=\RenyiPetzLow$. The thresholds are thus bracketed by
$[\RenyiSandLow,\RenyiSandHigh]$ and $[\RenyiPetzLow,\RenyiPetzHigh]$, in agreement with $\frac12$ and $2$, the
violating end being proved and the other end supported by the search. At the grid points where the theorems guarantee
the inequality, the largest certified $v$ is $\RenyiMaxDPISide$, and the double-precision values of $v$ agree with the
certified ones to $\RenyiCertErr$ ($d\in\{\RenyiDims\}$, $\varepsilon=\RenyiEps$, \RenyiStarts\ starts and
\RenyiSteps\ steps per $\alpha$; \RenyiMode\ run on \RenyiDevice).}{(Results not generated: run
\texttt{scripts/renyi\_threshold.py}.)} Table~\ref{tab:renyi} and Figure~\ref{fig:renyi} show $v(\alpha)$ on the grid.

\begin{table}[h]
\centering\footnotesize
\begin{tabular}{rcccc}
\toprule
 & \multicolumn{2}{c}{sandwiched $\widetilde D_\alpha$} & \multicolumn{2}{c}{Petz $D_\alpha$}\\
$\alpha$ & dephasing & channels & dephasing & channels\\
\midrule
\RenyiRows
\bottomrule
\end{tabular}
\caption[]{Largest certified violation $v(\alpha)$ of the data-processing inequality on $M_2$, for the complete
dephasing and for all channels with a two-dimensional environment (``--'': no violation above the tolerance
\IfFileExists{generated/renyi_threshold.tex}{$\RenyiTol$}{} found)\IfFileExists{generated/renyi_threshold.tex}{;
\RenyiMode\ run on \RenyiDevice}{}. The theorems guarantee $v\le0$ for $\alpha\ge\frac12$ (sandwiched) and
$\alpha\le2$ (Petz).}
\label{tab:renyi}
\end{table}

\begin{figure}[h]
\centering\fig{renyi_threshold.pdf}{\textwidth}
\caption{Largest certified violation $v(\alpha)$ (logarithmic axis) for the sandwiched (left) and the Petz (right)
divergences, for the complete dephasing and for channels with a $d$-dimensional environment (solid: $d=2$, dotted:
$d=3$); open symbols on the bottom line: no violation above the tolerance (dotted horizontal line) found. Dashed:
the known thresholds; shaded: the range in which the data-processing inequality is a theorem.}
\label{fig:renyi}
\end{figure}

\section{Symmetric tensor products and spin chains}\label{sec:chains}

The basis of \texttt{SpinChain} is, by construction, that of the iterated fused product of single sites:
the tests check that $S^z_i$ built by the chain coincides with $S^z_i\otimes1$ formed with
\texttt{kron} in the fused product of the algebra of \texttt{SpinChain(L-1)} with a single site. Beyond that, integrable
chains provide exact results for every part of the module (Table~\ref{tab:ring}, Figure~\ref{fig:chains};
script \texttt{scripts/chains.py}):
\begin{itemize}
\item \emph{Free fermions.} By the Jordan--Wigner transformation the open XX chain ($\Delta=0$) is a free-fermion
model with single-particle energies $\varepsilon_k=J\cos\frac{\pi k}{L+1}$, so the spectrum of sector $N$ consists
of all sums of $N$ distinct $\varepsilon_k$. For $L=\IfFileExists{generated/chains.tex}{\ChainFFL}{10}$ all
$2^L$ levels agree to \IfFileExists{generated/chains.tex}{$\ChainFFErr$}{machine precision}.
\item \emph{Heisenberg ring.} Ground-state energies agree with exact values (Bethe ansatz for $L=4$, independent exact diagonalisation
for the others, rounded to nine digits) to the nine digits of the table; $E_0/L$ approaches the Bethe-ansatz value $\tfrac14-\ln2\approx-0.4431$. Every level of sector
$N<L/2$ reappears in sector $N+1$ (SU(2) multiplets), up to
\IfFileExists{generated/chains.tex}{$\ChainSUErr$}{round-off}; Section~\ref{sec:symmetries} uses this symmetry to
reduce the blocks further.
\item \emph{Dynamics.} The melting of a domain wall $|\!\uparrow\cdots\uparrow\downarrow\cdots\downarrow\rangle$ in
the open XX chain, evolved in the sector $N=L/2$ (dimension \IfFileExists{generated/chains.tex}{\ChainDWDim\ for
$L=\ChainDWL$}{$\binom{L}{L/2}$}), reproduces the single-particle prediction
$\langle n_i(t)\rangle=[e^{iht}C(0)e^{-iht}]_{ii}$, with $h$ the single-particle hopping matrix and
$C_{ij}=\langle c_i^\dagger c_j\rangle$ the correlation matrix, to \IfFileExists{generated/chains.tex}{$\ChainDWErr$}{round-off}.
\item \emph{Entanglement.} The entropy of a block of $n$ sites of the XX ground state, obtained with
\texttt{reduced\_state}, agrees with Peschel's formula $S=-\sum_\nu[\nu\ln\nu+(1-\nu)\ln(1-\nu)]$ over the
eigenvalues $\nu$ of the restricted correlation matrix~\cite{peschel2003}, including the parity oscillations
of open chains, to \IfFileExists{generated/chains.tex}{$\ChainPesErr$}{round-off}.
\item \emph{Disorder.} For the Heisenberg chain with random fields $h_i\in[-W,W]$ the mean ratio of consecutive
level spacings in the middle of the spectrum crosses over from the GOE value $0.5307$ to the Poisson value
$2\ln2-1\approx0.3863$~\cite{oganesyan2007,atas2013} as $W$ grows; at these sizes this is a finite-size crossover,
the standard spectral diagnostic used in studies of many-body localisation.
Here the batch dimension holds the disorder realisations
(\IfFileExists{generated/chains.tex}{\ChainReal}{many} per point).
\item \emph{Momentum.} For periodic chains the basis $|a,k\rangle=N_a^{-1/2}\sum_{j=0}^{L-1}e^{-ikj}T^j|a\rangle$ over
orbit representatives $a$ of the translation $T$ (orbit length $R_a$, $N_a=L^2/R_a$; $a$ contributes to momentum $k$
iff $kR_a\in2\pi\mathbb Z$) splits sector $N$ into blocks of fixed momentum
$k=2\pi m/L$ (for $L=\IfFileExists{generated/chains.tex}{\ChainMomL}{12}$ and $N=L/2$ of sizes
\IfFileExists{generated/chains.tex}{\ChainMomDims}{about $\binom LN/L$}); the union of their spectra is the
spectrum of the sector to \IfFileExists{generated/chains.tex}{$\ChainMomErr$}{round-off}, and the ground state of
the antiferromagnetic Heisenberg ring has $k=0$ for even and $k=\pi$ for odd $L/2$ (Marshall's sign rule). Momentum
blocks with $k\ne0,\pi$ are complex in the plane-wave basis even for real Hamiltonians, one of the cases of
Section~\ref{sec:numerics} where the field is fixed by the problem, at least in this basis: combined with reflection,
time reversal leaves each $k$-block invariant, so for non-integrable translation-invariant Hamiltonians these blocks
show GOE statistics and admit a real basis; the XXZ ring itself is integrable and its blocks are not of GOE type.
\item \emph{Beyond dense matrices.} Lanczos with full re-orthogonalisation and restarts, batched over
Hamiltonians, reproduces the free-fermion ground-state energies of open XX chains and the Heisenberg ring in
sectors up to the sizes of Table~\ref{tab:krylov} (in the full run beyond the reach of dense diagonalisation); Krylov evolution of the domain wall for
$L=\IfFileExists{generated/chains.tex}{\ChainKDWL}{16}$ (sector dimension
\IfFileExists{generated/chains.tex}{\ChainKDWDim}{$\binom L{L/2}$}) agrees with the single-particle solution to
\IfFileExists{generated/chains.tex}{$\ChainKDWErr$}{round-off}.
\item \emph{Open chains.} With jump operators $S^-_i$ on every site (Lindblad terms between sectors $N\to N-1$)
and any $N$-conserving Hamiltonian the particle number performs a pure death process, so the sector
distribution must be binomial, $p_n(t)=\binom{N_0}{n}e^{-n\gamma t}(1-e^{-\gamma t})^{N_0-n}$; the library
reproduces it to \IfFileExists{generated/chains.tex}{$\ChainLossErr$}{the RK4 accuracy}.
\end{itemize}

\begin{table}[h]
\centering\small
\begin{tabular}{lrrrrr}
\toprule
model & sector dim & $E_0$ (Lanczos) & $E_0$ (exact) & difference & time [s]\\
\midrule
\ChainKrylovRows
\bottomrule
\end{tabular}
\caption[]{Batched Lanczos on sparse sector Hamiltonians (time includes building the sparse Hamiltonian):
ground-state energies against free fermions
($\sum$ of the $L/2$ lowest $\cos\frac{\pi k}{L+1}$) and the tabulated Heisenberg value (nine digits)\IfFileExists{generated/chains.tex}{;
\ChainMode\ run on \ChainDevice}{}.}
\label{tab:krylov}
\end{table}

\begin{table}[h]
\centering\small
\begin{tabular}{rrrrrr}
\toprule
$L$ & $\dim$ ($N=L/2$) & $E_0$ (library) & $E_0$ (exact) & difference & $E_0/L$\\
\midrule
\ChainRingRows
\bottomrule
\end{tabular}
\caption[]{Ground-state energy of the Heisenberg ring $H=\sum_i\mathbf S_i\cdot\mathbf S_{i+1}$; exact values to
nine digits\IfFileExists{generated/chains.tex}{ (\ChainMode\ run on \ChainDevice)}{}.}
\label{tab:ring}
\end{table}

\begin{figure}[h]
\centering\fig{chains.pdf}{\textwidth}
\caption{Left: $\langle S^z_i(t)\rangle$ after a domain-wall quench of the XX chain (light cone).
Middle: entanglement entropy of the XX ground state, library vs.\ Peschel's formula. Right: mean
level-spacing ratio of the random-field Heisenberg chain, with the GOE and Poisson values.}
\label{fig:chains}
\end{figure}

\section{Subalgebras, inclusions and symmetries}\label{sec:symmetries}

The algebras of the previous sections are declared by their sector sizes. Three modules relate algebras to each other
and to the symmetries that produce them: \texttt{inclusions} treats a subalgebra $\mathcal N\subset\M$ with its
conditional expectations and indices, \texttt{decompose} computes the algebra generated by given matrices or by a
finite group, and \texttt{su2} uses the full SU(2) symmetry of spin-$\frac12$ chains instead of $S^z_{\mathrm{tot}}$
alone.

\paragraph{Inclusions and conditional expectations.} Let $\mathcal N=\bigoplus_{c=1}^CM_{k_c}$ and
$\M=\bigoplus_{d=1}^DM_{m_d}$. Every unital inclusion $\iota:\mathcal N\to\M$ is, up to an inner automorphism
$\operatorname{Ad}u$ of $\M$, of the form $\iota(x)_d=\bigoplus_cx_c\otimes1_{\lambda_{dc}}$, given by the inclusion
matrix $\Lambda=(\lambda_{dc})\in\mathbb Z_{\ge0}^{D\times C}$ with $m_d=\sum_c\lambda_{dc}k_c$~\cite{ghj1989}:
$\lambda_{dc}$ is the multiplicity of the summand $c$ of $\mathcal N$ in the summand $d$ of $\M$. \texttt{Inclusion}
stores $\Lambda$ and an optional (batched) unitary $u$. With the isometries $V_{dc,i}:\C^{k_c}\to\C^{m_d}$ onto the
$i$-th copy of $\C^{k_c}$, $\iota(x)_d=\sum_{c,i}u_dV_{dc,i}x_cV_{dc,i}^*u_d^*$, so $\iota$ is an
\texttt{InterSectorChannel}, and inclusion matrices of towers multiply. A faithful trace
$\tau=\sum_dw_d\Tr_d$ on $\M$ restricts to $\mathcal N$ with the weights $\nu=\Lambda^Tw$, and
\[
  E(y)_c=\nu_c^{-1}\sum_{d,i}w_d\,V_{dc,i}^*u_d^*\,y_d\,u_dV_{dc,i}
\]
is the unique $\tau$-preserving conditional expectation of $\M$ onto $\mathcal N$: it is unital and completely
positive, $E\circ\iota=\mathrm{id}$, $E(\iota(n_1)y\,\iota(n_2))=n_1E(y)n_2$, and it is the adjoint of $\iota$ with
respect to $\tau$ and $\tau|_{\mathcal N}$. The tests check these properties against dense matrices for the three
traces of Section~\ref{sec:math}, for the Markov trace and for random weights. For a connected inclusion (the bipartite
graph of $\Lambda$ is connected) the Perron--Frobenius vector of $\Lambda\Lambda^T$ gives the Markov trace, and the
index is $[\M:\mathcal N]=\|\Lambda\|^2$~\cite{jones1983,ghj1989}: $(m/k)^2$ for factors $M_k\subset M_m$, $n^2$ for
$\C\subset M_n$ and $n$ for the diagonal subalgebra $\C^n\subset M_n$. The Watatani index
$\operatorname{Ind}E=\sum_ju_ju_j^*$ over a quasi-basis $\{u_j\}$~\cite{watatani1990} is central,
$(\operatorname{Ind}E)_d=(\Lambda\Lambda^Tw)_d/w_d$, and equals $[\M:\mathcal N]\,1$ exactly for the Markov trace. The
Pimsner--Popa constant $\lambda(E)=\max\{\lambda: E(y)\ge\lambda y\ \text{for all}\ y\ge0\}$~\cite{pimsner1986} is
$\bigl(\max_d\sum_c(\nu_c/w_d)\min(k_c,\lambda_{dc})\bigr)^{-1}$~\cite[Theorem~6.1]{pimsner1986}. For every subfactor
of a II$_1$ factor $\lambda(E)=[\M:\mathcal N]^{-1}$~\cite[Theorem~2.2]{pimsner1986}, but for Type~I inclusions the two
differ~\cite[6.5]{pimsner1986}: for $\C\subset M_n$, $\lambda(E)=1/n$, while
the constant of the complete order, $\max\{\lambda:E-\lambda\,\mathrm{id}\ \text{completely positive}\}=
1/\|\operatorname{Ind}E\|$ (\cite[Th\'eor\`eme~3.5]{bdh1988}; see~\cite[proof of Proposition~1.5 and
Example~1.6]{frank1998}), is $1/n^2$. The tests compare both constants with brute-force values
(bisection over rank-one $y$, and Choi matrices). The basic construction $\M_1=\langle\M,e_{\mathcal N}\rangle$ has
the sectors of $\mathcal N$, of sizes $(\Lambda^Tm)_c$, and $\M\subset\M_1$ has the inclusion matrix
$\Lambda^T$~\cite{jones1983,ghj1989}; the library builds the Jones projection $e_{\mathcal N}$ (tested:
$e_{\mathcal N}ye_{\mathcal N}=E(y)e_{\mathcal N}$, $e_{\mathcal N}\in\mathcal N'$, the Markov property
$\tau_1(ye_{\mathcal N})=[\M:\mathcal N]^{-1}\tau(y)$), the Jones tower $\mathcal N\subset\M\subset\M_1\subset\dots$
with constant index, and the relative commutant $\mathcal N'\cap\M=\bigoplus_{\lambda_{dc}>0}M_{\lambda_{dc}}$. All
maps act on batches by index gathers, without host synchronisation, and are differentiable.

\paragraph{The spin-chain tower.} The $U(1)$ algebras $A_L=\bigoplus_{N=0}^LM_{\binom LN}$ of
Section~\ref{sec:chains} form the tower $A_L\subset A_{L+1}$, $x\mapsto x\otimes1_2$, whose Bratteli diagram is
Pascal's triangle: sector $N$ of $A_{L+1}$ receives the sectors $N-1$ (new spin up) and $N$ (new spin down) of $A_L$.
$\Lambda$ is the incidence matrix of the path with $2L+3$ vertices, so
\[
  [A_{L+1}:A_L]=\|\Lambda\|^2=4\cos^2\frac{\pi}{2L+4}\;\nearrow\;4 ,
\]
values $4\cos^2(\pi/n)$, $n=4,6,8,\dots$ ($2,3,2+\sqrt2,\dots$), of the discrete series below $4$ in Jones' index
theorem~\cite{jones1983}. The tests
check the spectrum of the path graph and the index for $L\le8$, that the embedding is $x\otimes1_2$ in the basis of
\texttt{SpinChain}, and that the conditional expectation for the normalised trace of $(\C^2)^{\otimes(L+1)}$ is the
normalised partial trace over the new site. The program \repofile{inclusions.py} shows the interface: for
$M_2\oplus\C\subset M_4\oplus M_3$ with a random unitary it computes the index, the Markov conditional expectation and
its bimodule property, the Watatani index, both Pimsner--Popa constants and the Jones projection of the basic
construction, and the indices of the spin-chain tower (\texttt{inclusions.Inclusion}, \texttt{index},
\texttt{conditional\_expectation}, \texttt{pimsner\_popa\_constant}, \texttt{basic\_construction},
\texttt{jones\_projection}, \texttt{chain\_tower}).

\paragraph{From a symmetry to the algebra: Wedderburn decomposition.} Section~\ref{sec:intro} obtained the
algebra~\eqref{eq:M} from a symmetry through the commutant. \texttt{decompose} carries out this step numerically.
Given matrices $S_1,\dots,S_J\in M_n(\C)$ (for instance the generators of a Lie algebra representation) or all elements
$\pi(g)$ of a unitary representation of a finite group, it returns a unitary $V$ with
\[
  V^*\mathcal AV=\bigoplus_cM_{k_c}(\C)\otimes1_{m_c},\qquad V^*\mathcal A'V=\bigoplus_c1_{k_c}\otimes M_{m_c}(\C),
\]
for the von Neumann algebra $\mathcal A=\{S_j,S_j^*\}''$ generated by the matrices, together with maps between
$n\times n$ matrices and operators of the algebras $\bigoplus_cM_{k_c}$ and $\bigoplus_cM_{m_c}$ of the library. The
algorithm combines two randomised algorithms of Murota and coauthors. As in Maehara and Murota~\cite{maehara2011},
who apply the method of Murota, Kanno, Kojima and Kojima~\cite{murota2010} to the commutant, the
eigenspaces of a random Hermitian element of $\mathcal A'$ are, generically, irreducible $\mathcal A$-submodules, one
per copy of each irreducible representation; the copies are then grouped and aligned as in~\cite[Sec.~4.2]{murota2010},
here with elements of $\mathcal A'$ instead of $\mathcal A$. For a second random element $Y\in\mathcal A'$ and orthonormal bases
$Q_i,Q_j$ of two such submodules, $Q_i^*YQ_j$ intertwines them; by Schur's lemma it vanishes if they are inequivalent
and is a non-zero multiple of a unitary (generically) if they are equivalent, so the non-zero blocks give the
equivalence classes and their polar factors align the bases of the copies. Random elements of $\mathcal A'$ come from
the twirl $|G|^{-1}\sum_g\pi(g)X\pi(g)^*$ for a finite group, or from the null space of $X\mapsto([H_j,X])_j$ with
the Hermitian matrices $H_j$ given by $(S_j+S_j^*)/2$ and $(S_j-S_j^*)/2i$ (as in~\cite{maehara2011}), whose size is
reduced by first diagonalising a random Hermitian element of $\mathcal A$ (the first step of~\cite{murota2010}).
Every result is verified before it is returned: the generators lie in
$\mathcal B=V\bigl(\bigoplus_cM_{k_c}\otimes1_{m_c}\bigr)V^*$, fresh random elements of $\mathcal A'$ lie in
$\mathcal B'$, and $\sum_cm_c^2=\dim\mathcal A'$, which together give $\mathcal A=\mathcal B$ by the double commutant
theorem; a failed check triggers a new random draw. The decomposition is over $\C$ and uses dense linear algebra. For
an algebra closed under complex conjugation (real generators) entrywise conjugation maps every sector onto a sector,
$\overline{V_c}=V_d(J_c\otimes M_c)$ by Schur's lemma, and \texttt{real\_structure} reads off the type of every simple
summand of the real algebra (cf.~\cite{maehara2010} for the numerical decomposition of real matrix $*$-algebras with
general irreducible components): complex for $d\neq c$ (a pair of conjugate sectors), real for $d=c$ and
$J_c\overline{J_c}=1$, quaternionic for $J_c\overline{J_c}=-1$; the sign has a gap of 2, so the decision needs no
tolerance, and \texttt{real\_basis} returns a real orthogonal basis adapted to the types (Takagi factorisation of the
symmetric unitary $J_c$). In a quaternionic sector the multiplicity block of a real symmetric commuting operator has
doubly degenerate eigenvalues (Kramers). The tests recover block structures hidden
by random unitaries, the Fourier basis of the regular representation of $\mathbb Z_n$, and Schur--Weyl
duality~\cite{fulton1991}: on $(\C^2)^{\otimes L}$, $L\le6$, the algebra generated by the permutations of the tensor
factors and the algebra generated by the total spin have the same central projections, with $k_c$ and $m_c$
exchanged (the dimensions of the irreducible representations of $S_L$ of shape $(L-j,j)$ and $2S+1$, $S=\frac L2-j$).
The program \repofile{symmetry.py} shows this for $L=4$ (\texttt{decompose.from\_group} with
\texttt{permutation\_representation}, \texttt{decompose.from\_generators} with \texttt{total\_spin}) and computes the
SU(2) blocks and the lowest energy per total spin of a Heisenberg ring of 14 spins (\texttt{SU2Chain},
\texttt{heisenberg}, \texttt{ground\_state}, \texttt{lowest\_energies}).

\paragraph{SU(2) symmetry of spin chains.} The Heisenberg Hamiltonian $H=\sum J_{ij}\,\mathbf S_i\cdot\mathbf S_j$
commutes with all of SU(2), not only with $S^z_{\mathrm{tot}}$. On $(\C^2)^{\otimes L}=\bigoplus_S\C^{m_S}\otimes
V_S$, with $V_S$ the spin-$S$ representation,
\[
  \mathrm{SU}(2)'=\bigoplus_SM_{m_S}(\C),\qquad m_S=\binom{L}{L/2-S}-\binom{L}{L/2-S-1},\qquad
  \sum_S(2S+1)\,m_S=2^L ,
\]
and an invariant Hamiltonian is $\bigoplus_SH_S\otimes1_{2S+1}$: every multiplet is computed once instead of $2S+1$
times, and the blocks shrink from $\binom L{L/2+S}$ to $m_S$. \texttt{SU2Chain} realises $\C^{m_S}$ as the
highest-weight vectors $\ker S^+\cap\{S^z_{\mathrm{tot}}=S\}$ in the $U(1)$ sector $N=L/2+S$ of \texttt{SpinChain}.
Their orthonormal basis $B_S$ is built site by site with Clebsch--Gordan coefficients (the Yamanouchi basis, labelled by
the sequence of intermediate spins $S_1=\frac12,\dots,S_L=S$); it is orthonormal by construction, so no null space or
orthogonalisation is needed, and it is expressed directly in the $U(1)$ basis, where the block of an invariant
operator is $A_S=B_S^*A_NB_S$. Clebsch--Gordan coefficients are computed exactly in rational arithmetic (Racah's
formula~\cite{racah1942}, Condon--Shortley phases). By the Wigner--Eckart theorem~\cite{edmonds1957} the reduced
matrix elements of the site spins follow from highest-weight vectors alone, and $\mathbf S_i\cdot\mathbf S_j$ is
assembled from them as a sum over intermediate spins, so spin correlations need no sum over magnetic quantum numbers.
For $L=14$ the largest $U(1)$ block, $\binom{14}{7}=3432$, becomes $m_1=m_2=1001$, and the cubic cost of diagonalising
all blocks, $\sum_N\binom LN^3$ against $\sum_Sm_S^3$, drops by a factor of about $48$; for $L=16$ the largest block
shrinks from $12870$ to $3640$. The bases of all spins hold $\sum_Sm_S\binom L{L/2+S}$ numbers ($1.0\cdot10^8$ for
$L=16$), which limits dense SU(2) computations to about $L=16$. The tests compare the full spectra of open and periodic
chains, with a field and next-nearest-neighbour couplings, for $L\le12$ with dense diagonalisation (without next-nearest-neighbour
couplings also with the $U(1)$ sectors; to $10^{-10}$), and reproduce the ring energies of Table~\ref{tab:ring} for $L\le14$ (to $10^{-9}$, with total
spin $0$), the theorem of Lieb and Mattis~\cite{lieb1962} for open chains and even rings with $L\le12$ (the ground
state has $S=0$ for even and $S=\frac12$ for odd $L$, and the lowest energy $E_0(S)$ increases strictly with $S$), the
two exactly degenerate singlets with $E=-3L/8$ at the Majumdar--Ghosh point $J_2=J/2$~\cite{majumdar1969,majumdar1969b}, textbook
and SymPy values of the Clebsch--Gordan coefficients, and Wigner--Eckart matrix elements and spin correlations against
direct evaluation.


\noindent Section~\ref{sec:cs} uses the conditional expectations of this section to compute the Connes--St{\o}rmer
relative entropy of two maximal abelian subalgebras of $M_n$.

\section{Use cases}\label{sec:usecases}

Each of the following computations asks a question from the literature, is validated against independent
computations, and states what it reproduces and what is new. The first turns the characterisations of the trace into
quantitative statements: batched, certified optimisation over the density and the operators gives the size of the
failure of each inequality as a function of the Michelson contrast, and the sector structure shows that only the
contrast inside a block is detected. The second uses the inclusions of Section~\ref{sec:symmetries} to settle the
Connes--St{\o}rmer entropy of two MASAs of $M_2$ in closed form and to find, for $n\ge3$, partitions of unity that beat
the orthonormal bases. The third evolves a boundary time crystal with local dissipation in the permutation-invariant
algebra, with maps between sectors, to obtain finite-$N$ lifetimes; the fourth computes symmetry-resolved entanglement
in the XXZ chain directly from the sector blocks of the reduced state.

% Use case (short): quantitative failure of trace characterisations. Full account with proofs: man/trace_envelopes.
% Numbers: scripts/trace_characterizations.py -> generated/trace_characterizations.tex (macros \TC...).
\subsection{Use case: how strongly trace characterisations fail}\label{sec:tc}

Many inequalities that hold for the trace characterise it: a positive functional $\varphi$ that satisfies one of them for
all operators is a trace (Gardner~\cite{gardner1979}, and a long series of results surveyed in~\cite{bikchentaev2025survey}).
On $M_n$ every positive functional is $\varphi_W=\Tr(W\,\cdot\,)$, and we ask how strongly such an inequality $L\le R$ fails
for a non-tracial $\varphi_W$, measured by the envelope $V(\delta)=\sup\{(L-R)/N:\Delta(W)=\delta\}$ over all operator
variables, with a scale-invariant normalisation $N$. For \IfFileExists{generated/trace_characterizations.tex}{\TCItems}{25}
characterisations from the literature the library computes $V$ in one batched optimisation per inequality: the density
is parametrised with an exact contrast (\texttt{optimize.SpectrumParam}), the operators by \texttt{PositiveParam},
\texttt{SelfAdjointParam} and \texttt{UnitaryParam} (projections of every rank included), the lattice operations
$P\wedge Q$, $P\vee Q$ by functional calculus with a Daleckii--Krein gradient, and the best point of every contrast is
recomputed in mpmath with \IfFileExists{generated/trace_characterizations.tex}{\TCDigits}{40} digits, so that every reported
violation is certified. \IfFileExists{generated/trace_characterizations.tex}{The whole computation takes \TCSeconds~s on
\TCWorkers\ CPU processes.}{}

Figure~\ref{fig:tc} shows the result. The envelopes vanish at the trace to different orders, $\delta$, $\delta^2$,
$\delta^4$ and $\delta^8$; for sixteen inequalities the computed values agree with closed forms proved in the companion
paper~\cite{nn_trace} (the Kantorovich constant $\kappa-1$ for Gardner's inequality and six relatives, $\kappa^2-1$,
$h-1$, the Specht ratio and others), for four with explicit lower bounds, and five are known only numerically. On
$\bigoplus_cM_{k_c}$ only the contrast inside the blocks is detected: a density that varies only between the blocks
violates none of the inequalities that hold for every trace. The companion paper also prescribes the contrast of the
operators: the violation vanishes when either contrast is zero, for Gardner's inequality it equals
$(1-\Delta(W)^2\Delta(|X|)^2)^{-1/2}-1$ in even dimension, and the largest violations are attained, with few exceptions,
at projections or at two-level operators $\ell+(1-\ell)P$.

\begin{figure}[h]
\centering\fig{trace_characterizations.pdf}{\textwidth}
\caption[]{Certified violation envelopes $V(\delta)$ on $M_2$ (symbols) and the closed forms or lower bounds of~\cite{nn_trace}
(lines), grouped by the order of vanishing: (a) $\delta$; (b) $\delta^2$, sharp; (c) $\delta^2$, lower bounds and items
without closed form; (d) $\delta^4$ and $\delta^8$. (e) Search divided by the proved value, $n=2,3,4$. (f) Items that
depend on $n$.}
\label{fig:tc}
\end{figure}


% Use case (short): the Connes--Stormer entropy of two MASAs. Full account with proofs: man/connes_stormer.
% Numbers: scripts/connes_stormer.py -> generated/connes_stormer.tex (macros \CS...).
\paragraph{The Connes--St{\o}rmer entropy of two MASAs}\label{sec:cs} \cite{nn_cs}. The relative entropy
$H(\Delta_n|U\Delta_nU^*)$ of the diagonal algebra and a rotated copy is a supremum over partitions of unity; with the
conditional expectations of \texttt{Inclusion} and partitions parametrised by isometries, many unitaries and starts form
one batch. For $n=2$ the computation suggested a closed form, proved in~\cite{nn_cs}; for $n=3,4,5$ it finds partitions
that beat every orthonormal basis found.


\subsection{Use case: a boundary time crystal under local dissipation}\label{sec:btc}

\paragraph{The question.} $N$ two-level systems with a collective drive and collective decay,
\begin{equation}\label{eq:btc}
  \dot\rho=-i[\omega J_x,\rho]+\frac{\kappa}{S}\,\mathcal D[J_-]\rho
  +\gamma\sum_{i=1}^N\mathcal D[\sigma_-^{(i)}]\rho+\gamma_\phi\sum_{i=1}^N\mathcal D[\sigma_z^{(i)}]\rho,\qquad
  \mathcal D[L]\rho=L\rho L^*-\tfrac12\{L^*L,\rho\},
\end{equation}
with $S=N/2$ and $J_a=\frac12\sum_i\sigma_a^{(i)}$. For $\gamma=\gamma_\phi=0$ this is the \emph{boundary time
crystal} of Iemini et al.~\cite{iemini2018}: for $\omega>\kappa$ the magnetisation $m_z=\langle J_z\rangle/S$
oscillates, the oscillations are damped at finite $N$ with a lifetime that grows linearly with
$N$~\cite{liu2026,nemeth2025}, and they persist as
$N\to\infty$ (a limit cycle of the mean-field equations, which become exact as $N\to\infty$ at fixed
times; Carollo and Lesanovsky~\cite{carollo2021} prove this for open Dicke models and note that the proof adapts to
all-to-all spin models and to individual spin dissipation). Real emitters also decay and dephase individually. Such
local terms do not conserve the total spin and, in the mean-field equations below, turn the limit cycle into a damped
oscillation, as Piccitto et al.~\cite{piccitto2021} found in mean field for jump operators acting collectively on
blocks of $N_s<N$ spins (damping rate $\propto1/N_s$, in a related $p,q$-spin model); robustness of
boundary time crystals has been studied in multi-level settings~\cite{prazeres2021}, and Tucker et
al.~\cite{tucker2018} studied single-particle dephasing in a related synchronised phase. We ask how long the
oscillations live as a function of $N$ and of the local rate, and how the two damping mechanisms---the collective
one, which vanishes as $N\to\infty$, and the local one, which does not---combine.

\paragraph{The algebra.} By Schur--Weyl duality, $(\C^2)^{\otimes N}=\bigoplus_jV_j\otimes\C^{m_j}$ with $\dim
V_j=2j+1$, $j=N/2,N/2-1,\dots\ge0$, and $m_j=\binom N{N/2-j}-\binom N{N/2-j-1}$, and the operators that commute with
the permutations of the sites form
\[
  \M=\bigoplus_jM_{2j+1}(\C)\otimes1_{m_j},\qquad \Tr_{(\C^2)^{\otimes N}}(A)=\sum_jm_j\Tr A_j .
\]
The right-hand side of~\eqref{eq:btc} commutes with the permutations, so a permutation-invariant state stays one:
$\rho=\bigoplus_j\sigma_j\otimes1_{m_j}/m_j$, where $\sigma=\bigoplus_j\sigma_j$ is the density with respect to
$\Tr_{\mathrm{blunt}}$ in the sense of Section~\ref{sec:math} (the ``Dicke-basis'' density of PIQS~\cite{shammah2018}).
The collective operators act inside the blocks, as the spin-$j$ matrices. A local term is a map between sectors: for
$L\in\{\sigma_-,\sigma_+,\sigma_z\}$ the map $\Phi(X)=\sum_iL^{(i)}XL^{(i)*}$ is normal and completely positive on
$\M$, and restricting to the last site and branching $N\to N-1$ ($\C^{m_j}$ is the direct sum of the multiplicity
spaces of $j_1=j\pm\frac12$ for $N-1$ sites, which $L^{(N)}$ preserves, and the $N$ sites contribute equally) gives
its Kraus operators on blunt densities,
\begin{equation}\label{eq:btc-kraus}
  \Phi(\sigma)_{j'}=\sum_j\sum_{j_1}K^{j'j}_{j_1}\sigma_jK^{j'j*}_{j_1},\qquad
  K^{j'j}_{j_1}=\Bigl(\frac{N\,m^{(N-1)}_{j_1}}{m^{(N)}_j}\Bigr)^{1/2}
  \bigl\langle j_1,\tfrac12;j'\bigr|\,1\otimes L\,\bigl|j_1,\tfrac12;j\bigr\rangle ,
\end{equation}
with $|j'-j|\le1$, $j_1\in\{j\pm\frac12\}\cap\{j'\pm\frac12\}$, and the coupled multiplets of $V_{j_1}\otimes\C^2$ in
the last factor; each $K$ has one non-zero diagonal, and $\Phi^*(1)=\sum_iL^{(i)*}L^{(i)}$ is $S+J_z$, $S-J_z$ or $N$.
(Closed forms of these coefficients are given in~\cite{chase2008,baragiola2010,shammah2018}; Eqs.~(40)--(46)
of~\cite{baragiola2010}, whose $\hat\sigma_z$ is half the Pauli matrix and whose $\rho_{J,M;J,M'}$ are the entries of
the blunt density, agree with~\eqref{eq:btc-kraus}. Permutational symmetry with individual decay and dephasing is also
exploited by the library PsiQuaSP~\cite{gegg2017} and by the method of Kirton and Keeling~\cite{kirton2017}.) In the library,
\texttt{permutational.PermutationalSystem(N)} is the algebra with these sectors; \texttt{local\_jump} returns $\Phi$ as an
\texttt{InterSectorChannel} of rank two, so \texttt{dynamics.lindblad\_evolve} handles~\eqref{eq:btc} with the
collective jumps as operators and the local ones between sectors, unchanged. Because every operator in~\eqref{eq:btc}
is banded in the basis $|j,m\rangle$, \texttt{PermutationalLindbladian} evaluates the same generator by shifted
elementwise products, $O(N^3)$ operations per state instead of $O(N^4)$ for block products, with all rates batched.
We drive about $y$ instead of $x$: the rotation $e^{-i\pi J_z/2}$ commutes with every dissipator, and the generator
becomes real, so the computation runs in the real algebra $\bigoplus_jM_{2j+1}(\mathbb R)$.
\IfFileExists{generated/permutational.tex}{For $N=\PMNmax$ a state has $\sum_j(2j+1)^2=\PMDimNmax$ entries, a density
matrix on $(\C^2)^{\otimes N}$ has $4^N\approx\PMDimFullNmax$.}{}

\noindent The program \repofile{permutational.py} makes one batched run over four local emission rates for $N=30$
in the real algebra $\bigoplus_jM_{2j+1}(\mathbb R)$, from all spins down, and returns $m_z(t)$ and the distribution
of $j$ (\texttt{PermutationalSystem}, \texttt{PermutationalLindbladian}, \texttt{coherent\_state},
\texttt{dynamics.rk4}, \texttt{m\_values}).

\paragraph{Validation.} \IfFileExists{generated/permutational.tex}{%
(i)~The banded generator, with all nine terms (drive about $x,y,z$, collective and local emission, pumping and
dephasing), equals the exact Lindblad generator on $(\C^2)^{\otimes N}$, compressed to $\M$ through Dicke bases built
independently from $J^2$ and $J_z$, to $\PMFullErr$ for $N\le\PMFullNmax$; an RK4 trajectory computed on
$(\C^2)^{\otimes\PMTrajN}$ gives the same blocks to $\PMTrajErr$.
(ii)~The generic route (\texttt{lindblad\_evolve} with \texttt{InterSectorChannel} jumps) agrees with the banded one
to $\PMLibErr$ ($N=\PMLibN$; $\PMLibSecGeneric$~s against $\PMLibSecBanded$~s).
(iii)~Against \texttt{qutip.piqs}~\cite{shammah2018} (QuTiP \PMPiqsVersion): the generators agree to a relative
$\PMPiqsGenErr$ (PIQS stores its rates in single precision, and its local dephasing rate is $4\gamma_\phi$); trajectories
of~\eqref{eq:btc} for $N=\PMPiqsNs$ agree to $\PMPiqsTrajErr$ in $m_z$ and to $\PMPiqsBlockErr$ in the blocks; the
steady states for $N=\PMPiqsSsNs$ to $\PMPiqsSsErr$.
(iv)~Pure collective decay $\mathcal D[J_-]$ (rate $\Gamma=1$) from $|S,S\rangle$ is a pure death process with rates
$\Gamma(S+m)(S-m+1)$, so $\int_0^\infty P(m>0)\,dt=\sum_{m=1}^S[\Gamma(S+m)(S-m+1)]^{-1}=H_N/(\Gamma(N+1))$ with
the harmonic number $H_N$; Table~\ref{tab:btc-sr} reproduces it to $\PMSRErr$ for $N\le\PMSRNmax$, together with the
superradiant delay $t_D\approx\ln N/(\Gamma N)$ and the reduction of the peak of the mean intensity below the
mean-field value $\Gamma N^2/4$ by the fluctuations of the delay~\cite{gross1982}.
(v)~Local emission, pumping and dephasing alone keep a product state a product of identical qubit states, so
$\langle J_z\rangle$, $\langle J_x\rangle$ and $\langle J^2\rangle$ follow from the Bloch equations of one qubit,
and the spin distribution is $p_j=m_j\sum_{m=-j}^j\lambda_+^{S+m}\lambda_-^{S-m}$ with the eigenvalues $\lambda_\pm$ of
the one-qubit state. For $N=\PMProdN$ the moments agree to a relative $\PMProdErr$ and $\sum_j|p_j-p_j^{\rm
exact}|=\PMProdPjErr$; this tests the multiplicities in~\eqref{eq:btc-kraus} far from the reach of the full space.
(vi)~Halving the time step changes $m_z$ by at most $\PMConvMz$ and $\tau$ by $\PMConvTau$ ($N=\PMConvN$).%
}{Run \texttt{scripts/permutational.py} to generate the numbers.}

\begin{table}[ht]
\centering\small
\caption{Pure collective decay from $|S,S\rangle$ ($\Gamma=1$): the integral of $P(m>0)$ against the exact value
$H_N/(N+1)$, the time $t_D$ of the peak of the mean intensity $\Gamma\langle J_+J_-\rangle$ in units of $\ln N/N$, and
the peak in units of $N^2/4$ (mean field: $1$ and $1$ for $N\to\infty$).}\label{tab:btc-sr}
\IfFileExists{generated/permutational.tex}{%
\begin{tabular}{rccccc}
\toprule
$N$ & $\int P(m>0)$ & $H_N/(N+1)$ & rel.\ dev. & $t_DN/\ln N$ & $I_{\max}/(N^2/4)$\\
\midrule
\PMSRRows
\bottomrule
\end{tabular}}{\fbox{\small generated/permutational.tex not found}}
\end{table}

\paragraph{The sweep.} We take $\omega=\IfFileExists{generated/permutational.tex}{\PMOmega}{2}\kappa$, all spins down
at $t=0$, and define the lifetime $\tau$ as the first time at which the half-swing of $m_z$ between consecutive
extrema falls below $1/e$ of the first one (log-linear interpolation between the midpoints of the swings). For every
$N$ one batched run holds all local rates at once
(\IfFileExists{generated/permutational.tex}{$\gamma\in\{\PMGe\}$, $\gamma_\phi\in\{\PMGd\}$, $\PMBatch$ parameter sets
for $N=\PMNs$}{ten parameter sets}); a batch element is dropped once its lifetime is settled, and the run ends at
$\kappa t=\IfFileExists{generated/permutational.tex}{\PMT}{45}$ (censored entries are marked $>$). Without local
terms only the sector $j=S$ is populated, and that run uses the single block $M_{N+1}$
(\texttt{PermutationalSystem(N, j\_min=N/2)}). The mean-field limit $\tau_{\mathrm{MF}}(\gamma)$ is the same lifetime for
the equations $\dot m_x=\omega m_z+\kappa m_xm_z-\Gamma_2m_x$, $\dot m_y=\kappa m_ym_z-\Gamma_2m_y$, $\dot
m_z=-\omega m_x-\kappa(m_x^2+m_y^2)-\gamma(1+m_z)$, $\Gamma_2=\gamma/2+2\gamma_\phi$ (drive about $y$).

\begin{table}[ht]
\centering\small
\caption{Lifetime $\kappa\tau$ of the oscillations of $m_z$ (parameters of the text): without local terms ($0$), with
local emission $\gamma/\kappa$ (columns 3--7) and with local dephasing $\gamma_\phi/\kappa$ (columns 8--12); last row:
mean field ($N\to\infty$). $>$: no lifetime before the end of the run.}\label{tab:btc}
\IfFileExists{generated/permutational.tex}{%
\begin{tabular}{r|c|ccccc|ccccc}
\toprule
\PMTauHeader\\
\midrule
\PMTauRows
\bottomrule
\end{tabular}}{\fbox{\small generated/permutational.tex not found}}
\end{table}

\begin{figure}[ht]
\centering
\fig{permutational.pdf}{\textwidth}
\caption{(a) $m_z(t)$ for the largest $N$ without and with local emission, and the mean-field curves (dotted).
(b) $\tau(N,\gamma)$ for local emission; dotted: $\tau_{\mathrm{MF}}(\gamma)$, dashed: the additive
rule~\eqref{eq:btc-rule}. (c) All lifetimes with a finite mean-field limit, $\tau/\tau_{\mathrm{MF}}$ against
$x=\tau(N,0)/\tau_{\mathrm{MF}}$; the dashed curve $x/(1+x)$ is~\eqref{eq:btc-rule}. (d) The relative deviation from
the rule; points of equal rate are joined, $N$ increasing to the right (filled: emission, open: dephasing).}\label{fig:btc}
\end{figure}

\paragraph{Results.} \IfFileExists{generated/permutational.tex}{%
Table~\ref{tab:btc} and Figure~\ref{fig:btc}. Without local terms the lifetime grows linearly,
$\kappa\tau(N,0)\approx\PMSlope\,N+\PMIntercept$ for $\PMFitNmin\le N\le\PMNmax$ (least squares). The growth is
linear, as expected: for $\omega\gg\kappa$ the oscillations of $m_z$ from all spins down decay at the rate
$3\kappa/(2N)$~\cite{liu2026,nemeth2025}, which corresponds to $\kappa\tau\approx2N/3$; our $\omega=\PMOmega\kappa$ lies
outside that perturbative regime, and the slope is smaller. With local
terms it saturates at the mean-field value (last row of Table~\ref{tab:btc}), and the two damping mechanisms combine
approximately like independent decay rates,
\begin{equation}\label{eq:btc-rule}
  \frac1{\tau(N,\gamma)}\approx\frac1{\tau_{\mathrm{rule}}}=\frac1{\tau(N,0)}+\frac1{\tau_{\mathrm{MF}}(\gamma)} ,
\end{equation}
for local emission and local dephasing alike: in the variables $x=\tau(N,0)/\tau_{\mathrm{MF}}$ and
$\tau/\tau_{\mathrm{MF}}$ all data lie close to the curve $x/(1+x)$ (Figure~\ref{fig:btc}c), and the crossover size at
which the local damping takes over ($x=1$) is $N^*\approx\kappa\tau_{\mathrm{MF}}(\gamma)/\PMSlope$. Over the
$\PMRuleCount$ entries with finite $\tau$ and $\tau_{\mathrm{MF}}$, $\tau/\tau_{\mathrm{rule}}-1$ lies between
$\PMRuleMin$ and $\PMRuleMax$, and between $\PMRuleBigMin$ and $\PMRuleBigMax$ for $N\ge\PMRuleNmin$. The deviation has a
structure (Figure~\ref{fig:btc}d): for the larger rates $\tau$ exceeds $\tau_{\mathrm{rule}}$, by an amount that
decreases at large $N$; for the smallest rates the deviation decreases through zero as $N$ grows, so that there
$\tau<\tau_{\mathrm{rule}}$ at the largest $N$; at equal $x$ local dephasing deviates more than local emission.
So~\eqref{eq:btc-rule} describes the lifetimes to within the range just stated; it is not an asymptotic law that
our data establish.%
}{Run \texttt{scripts/permutational.py}.}

\paragraph{What is reproduced and what is new.} Without local terms, the growth of the lifetime with $N$ is known:
Iemini et al.~\cite{iemini2018} report a decay rate that vanishes as a power of $N$, Piccitto
et al.~\cite{piccitto2021} state that it decreases as $N^{-1}$ for this model, and for $\omega\gg\kappa$ Liu
et al.~\cite{liu2026} and Nemeth et al.~\cite{nemeth2025} derive the rate $3\kappa/(2N)$; only the slope
\IfFileExists{generated/permutational.tex}{\PMSlope}{0.45} at
$\omega=\IfFileExists{generated/permutational.tex}{\PMOmega}{2}\kappa$ is, as far as we know, new. The damping of the
mean-field limit cycle by non-collective dissipation was found by Piccitto et al.~\cite{piccitto2021} (for local
decay see also~\cite{wang2025}), and for local terms it follows directly from the mean-field equations above; that
local dephasing and local losses wash out the oscillations at finite $N$ was shown with PIQS for
$N=30$~\cite{shammah2018}. The superradiance checks reproduce textbook results~\cite{dicke1954,gross1982}. The
finite-$N$ lifetimes with local
emission or dephasing for $N$ up to \IfFileExists{generated/permutational.tex}{\PMNmax}{80} (Table~\ref{tab:btc}), the
approximate additivity~\eqref{eq:btc-rule} of the two damping rates and the structure of its deviation are, as far as
we know, new numbers. They are observations for this model, these rates and this definition of the lifetime; we have
no proof of~\eqref{eq:btc-rule}.

\paragraph{Cost.} \IfFileExists{generated/permutational.tex}{%
The batched run for $N=\PMNmax$ ($\PMBatch$ parameter sets) took $\PMSecNmax$~s, the whole sweep $\PMSecSweep$~s and
the script with all checks $\PMSeconds$~s on one CPU thread (\PMMode\ mode, \PMDevice; the full mode, with $N$ up to
$200$, is meant for a GPU, the cost per step growing as $N^3$ and the number of steps as $N$). For $N=\PMPiqsNmax$
one trajectory of the validation took $\PMPiqsSecOurs$~s here and $\PMPiqsSecBuild+\PMPiqsSecSolve$~s with PIQS
(Liouvillian, \texttt{mesolve}) at matched accuracy, and the generic route with the dense
\texttt{InterSectorChannel} needs $2C^2k_{\max}^2$ Kraus entries ($C=N/2+1$ sectors, $k_{\max}=N+1$), which limits it
to a few tens of spins. On $(\C^2)^{\otimes N}$ the state alone would have $\PMDimFullNmax$ entries for
$N=\PMNmax$.%
}{}


% Use case (short): symmetry-resolved entanglement in the XXZ chain. Full account: man/btc_sre (Part B).
% Numbers: scripts/symmetry_resolved.py -> generated/symmetry_resolved.tex (macros \SR...).
\subsection{Use case: symmetry-resolved entanglement in the XXZ chain}\label{sec:sre}

For a pure state invariant under a $U(1)$ symmetry, the reduced state $\rho_A$ commutes with the charge $Q_A$ of the
subsystem and splits into sector blocks, $\rho_A=\bigoplus_qp(q)\rho_A^{(q)}$, and $S(\rho_A)=\sum_qp(q)S(q)+H(p)$
\cite{laflorencie2014,goldstein2018,xavier2018}. In critical chains the resolved entropies $S(q)$ are equipartitioned to
leading order~\cite{xavier2018}; the corrections and their asymptotics for free fermions are due to Bonsignori, Ruggiero
and Calabrese~\cite{bonsignori2019}, and the variance of $Q_A$ grows as $(K/\pi^2)\log\ell$ with the Luttinger parameter
$K$~\cite{song2012}. In the library the reduced state of a charge-conserving state \emph{is} an operator of
$\bigoplus_qM_{d_q}$, so the resolved quantities are read off its blocks without projecting.

For the XXZ chain at $L=\IfFileExists{generated/symmetry_resolved.tex}{\SRL}{24}$, the ground states for eight anisotropies
come from one batched Lanczos run in the sector $S^z=0$, and the resolved entropies, the charge statistics and the widths
$B_n$ of the charge distributions of $\rho_A^n$ follow for all subsystem sizes (Figure~\ref{fig:sre}). The charge
statistics give the Luttinger parameter to within
\IfFileExists{generated/symmetry_resolved.tex}{$\SRKFcsDevPct\,\%$ for $\SRKFcsMidRange\le\Delta\le\SRKFcsMidRangeHi$}{about
one per cent} (about $4\%$ at $\Delta=-0.75$, and with the known logarithmic corrections at the isotropic point); the
$1/n$ law of the widths $B_n$ is tested and the amplitude $G_1(\Delta)$ of the breaking of equipartition measured; for free
fermions the asymptotic value $1/(2\pi^2)$ of~\cite{bonsignori2019} is recovered and its slow approach shown. Details,
validation and tables are in the companion paper~\cite{nn_btc}.

\begin{figure}[h]
\centering\fig{symmetry_resolved.pdf}{\textwidth}
\caption{Periodic chain, $L=\SRL$. Left: $nB_n$ against $\log\ell_c$, $n=1$ (filled circles), $2$ (squares), $3$
(triangles); lines: slope $K/\pi^2$ (Bethe ansatz). Middle: $K$ from the charge statistics and from the widths, against
the Bethe-ansatz value. Right: $S(q)-S(\ell/2)$ at $\ell=\SRHalf$ for $\Delta\in\{\SRDeltas\}$ (dark to light); dotted: the
Gaussian form, exact at $q-\ell/2=\pm1$.}
\label{fig:sre}
\end{figure}


\section{Comparison with NumPy/SciPy}\label{sec:baselines}

To separate what the library adds from what plain array code already gives, the same computations were written
three ways (\texttt{scripts/baselines.py}; Table~\ref{tab:baselines}). (A)~The Monte Carlo of Exp.~1: $B$ pairs of
positive operators with prescribed contrast and Haar-orthogonal eigenvectors in every sector, a Haar-orthogonal $U\in\M$,
and $z$. The direct NumPy/SciPy formulation loops over samples and sectors with \texttt{scipy.stats.ortho\_group};
a hand-vectorised NumPy version uses stacked QR decompositions with the phase correction; the library needs three
calls. (B)~Ground states of sector Hamiltonians: \texttt{scipy.sparse.linalg.eigsh} (ARPACK~\cite{lehoucq1998}), one
call per Hamiltonian, against the batched Lanczos solver. The three Monte Carlo versions sample the same distribution:
their means differ by at most \IfFileExists{generated/baselines.tex}{\BaseMaxSigma}{a few} standard errors, a two-sample
Kolmogorov--Smirnov test of the library against vectorised NumPy gives $p\ge\IfFileExists{generated/baselines.tex}{\BaseMinKS}{0.01}$,
and the ground-state energies agree to round-off. All times are medians of up to five runs after a warm-up.

The comparison is instructive in both directions. On a CPU the gain over the loop depends on the block size: it is
large for small blocks, where the Python overhead of the loop dominates, and small for large blocks, where both spend
their time in the same LAPACK routines. Against carefully vectorised NumPy the library is not systematically faster on a
CPU (Table~\ref{tab:baselines}): both call the same kernels. Its value there is that the vectorised, correctly
sampled, block-aware version is the default rather than an expert rewrite, and that the same code runs on a GPU
and can be differentiated. The NumPy versions work in double precision, so the library's CPU and GPU columns do too;
the last column repeats the GPU run in single precision, since double precision is 32--64 times slower on consumer
GPUs. For ground states, ARPACK is the reference for a single Hamiltonian on a CPU; our Lanczos solver
re-orthogonalises fully for robustness and is batched over Hamiltonians, which is what lets disorder averages and GPUs
pay off (Table~\ref{tab:baselines}).

\begin{table}[h]
\centering\small
\begin{tabular}{rrrrrrr}
\toprule
$C$ & $k$ & NumPy loop & NumPy vectorised & library (CPU) & library (GPU) & GPU, FP32\\
\midrule
\BaseMCRows
\bottomrule
\end{tabular}\\[2mm]
\begin{tabular}{rrrrrrr}
\toprule
$L$ & sector dim & Hamiltonians & \texttt{eigsh} [s] & library, CPU [s] & library, GPU [s] & $\max|\Delta E_0|$\\
\midrule
\BaseGSRows
\bottomrule
\end{tabular}
\caption[]{Top: time per Monte Carlo sample [$\mu$s] for (A); bottom: ground states of Heisenberg rings
(one Hamiltonian) and random-field chains ($W=3$, several Hamiltonians) in the sector $N=L/2$, total time for all
Hamiltonians (``--'': no GPU in this run)
\IfFileExists{generated/baselines.tex}{(\BaseMode\ run on \BaseDevice)}{}.}
\label{tab:baselines}
\end{table}

\section{Performance}\label{sec:performance}

Table~\ref{tab:benchmark} and Figure~\ref{fig:benchmark} report the wall-clock time per sample of library
operations (not of isolated kernels), in single precision, with batch $B=\min\bigl(4096,\,\IfFileExists{generated/benchmark.tex}{2^{\BenchTargetLog}}{2^{22}}/(Ck^2)\bigr)$ (shown in the table)\IfFileExists{generated/benchmark.tex}{ (\BenchMode\ run on \BenchDevice;
CPU with \BenchCpuThreads\ threads; PyTorch \BenchTorch, CUDA \BenchCuda, cuSOLVER \BenchCusolver)}{}. The versions matter:
batched Hermitian eigensolvers on GPUs cover all block sizes only in recent cuSOLVER and PyTorch releases (older
releases fall back to a loop over the matrices for blocks larger than $32$), and batched SVD on GPUs is limited to blocks
of size at most $32$, which is why the library uses the Hermitian eigensolver wherever the operator is self-adjoint
(and the eigensolver of the Hermitian dilation $\bigl(\begin{smallmatrix}0&A\\A^*&0\end{smallmatrix}\bigr)$ for singular values of
larger blocks). Batched QR on GPUs factorises small matrices one at a time; Haar sampling therefore uses a batched
Gram--Schmidt orthonormalisation with reorthogonalisation on GPUs, which gives the same matrices as QR with the phase
correction. On the CPU, PyTorch's batched eigen-, singular-value and QR decompositions loop over the matrices on one
core; the library splits such batches into one chunk per thread (identical results), so the CPU column uses all cores.
The operations are: generating random operators with prescribed spectra, $f(A)$ by eigendecomposition, $|A|$ (singular
values), the Michelson contrast, applying a channel, the Lindblad channel $e^{t\mathcal L}$, and one gradient step of
the optimisation\ifnum\BenchHasGPU=1; the table gives the CPU and GPU times and their ratio\fi.

Besides speed, the structure of the computation is checked directly (\texttt{scripts/gpu\_check.py} on a GPU,
\texttt{tests/test\_sync\_free.py} on a CPU): the operations read no values back to the host (each such read
synchronises a GPU), and their number of kernel launches does not grow with the number of sectors $C$ at a fixed total
size. Per-sector work of this kind made the random generation of an earlier version of the library slower on a GPU
than on a CPU; the check would have found it.

The first version of this work reported GPU speedups for raw kernels on a P100 and found them only for part of the
configurations. That observation reflects where GPUs help: a GPU needs enough independent work per kernel launch.
For computations in~\eqref{eq:M} this means many samples, many sectors or large blocks, and matrix products rather
than decompositions of small matrices. The design follows from it: everything is batched, so that the batch
dimension supplies the parallel work, and the cost module warns before configurations that would be slow. For
workloads that are small or dominated by a single eigenproblem, the CPU (or SciPy, Section~\ref{sec:baselines}) is
the right tool, and the library runs there unchanged.

\begin{table}[p]
\centering\footnotesize\renewcommand{\arraystretch}{0.85}
\ifnum\BenchHasGPU=1
\begin{tabular}{lrrrrrr}
\toprule
operation & $k$ & $C$ & batch & GPU [$\mu$s] & CPU [$\mu$s] & speedup\\
\else
\begin{tabular}{lrrrr}
\toprule
operation & $k$ & $C$ & batch & time per sample [$\mu$s]\\
\fi
\midrule
\BenchmarkRows
\bottomrule
\end{tabular}
\caption[]{Wall-clock time per sample of library operations for $C$ sectors of size $k$ and the batch size shown\IfFileExists{generated/benchmark.tex}{ (\BenchMode\ run on \BenchDevice, CPU with \BenchCpuThreads\ threads)}{}
(single precision; ``exp($tL$)'' uses a smaller batch because its superoperators have $k^4$ entries). With a GPU the
table also gives the CPU time and the speedup CPU/GPU.}
\label{tab:benchmark}
\end{table}

\begin{figure}[h]
\centering\fig{benchmark.pdf}{0.8\textwidth}
\caption{\ifnum\BenchHasGPU=1 Speedup CPU/GPU (logarithmic axis; above 1 the GPU is faster)\else Time per sample
(logarithmic axis)\fi\ of the operations of Table~\ref{tab:benchmark}, one bar per configuration $(k,C)$.}
\label{fig:benchmark}
\end{figure}

\section{Limitations and roadmap}\label{sec:limitations}

\begin{itemize}
\item \emph{Layout.} Blocks are padded to $k_{\max}$. For very different sector sizes (binomial sectors of
spin chains) one works in single sectors, and \texttt{SU2Chain} builds the block of a single total spin without
padding; operators acting on all sectors at once would need a bucketed layout, which is not implemented yet.
\item \emph{Symmetries.} Abelian charges with an arbitrary fusion rule and lattice translations are
supported, and for spin-$\frac12$ chains the SU(2) multiplets reduce the blocks further (\texttt{su2}; dense up to
$L\approx16$, Section~\ref{sec:symmetries}). Not yet implemented: SU(2) combined with momentum sectors, chains of
spins larger than $\frac12$, reflections and spin flip, and sparse (Lanczos) solvers in the SU(2) blocks. The
numerical Wedderburn decomposition handles any symmetry given by matrices, but with dense linear algebra; it
identifies the real, complex and quaternionic summands of real algebras, but the core algebra class stores only real
and complex blocks (quaternionic blocks are kept as complex blocks together with their $J_c$). Dense methods reach about 16--18 spins
in the $U(1)$ sectors, the sparse Lanczos and
Krylov solvers handle the half-filled sector of 24 spins (dimension $2.7\cdot10^6$; Table~\ref{tab:krylov} in the
full run), and about 26 spins on a GPU with tens of gigabytes (an estimate from the memory of the Krylov basis); tensor-network methods are out of
scope.
\item \emph{Optimisation.} The search is local; envelopes are lower bounds for suprema (upper bounds for
infima). Gradients of the functional calculus are finite at degenerate spectra, but gradients through the
eigenvectors returned by \texttt{eigh} are not;
second derivatives of $f(A)$ are correct at non-degenerate spectra only.
\item \emph{Divergences.} For $\alpha<1$ the sandwiched divergence depends on all eigenvalues of
$\sigma^g\rho\,\sigma^g$, $g=(1-\alpha)/(2\alpha)$, whose condition number grows without bound as $\alpha\to0$; for small $\alpha$ and ill-conditioned states the computed value is not accurate (hence the restriction
$\alpha\ge0.2$, $\rho,\sigma\ge\varepsilon/d$ in Section~\ref{sec:renyi}).
\item \emph{Scale.} Choi matrices, superoperators and $e^{t\mathcal L}$ work with $k_c^2\times k_c^2$
matrices and are limited to blocks of a few tens; RK4 integration scales to larger blocks.
\item \emph{Scope.} Finite dimensions and Type~I only; no tensor-network methods, and sparse methods only for
sector Hamiltonians (Lanczos and Krylov).
\end{itemize}

\section{Availability and reproducibility}

The code is available under the MIT licence at
\begin{center}\url{https://github.com/andre-a-no/von_neumann_type_i}.\end{center}
\texttt{bash scripts/run\_all.sh full cuda} regenerates every table and figure of this text on a GPU;
\texttt{bash scripts/run\_all.sh} runs a reduced version on a CPU in a few minutes (written to
\texttt{.quick\_output/}, not into the paper's tables), and \texttt{python scripts/gpu\_check.py} checks a GPU in a few
minutes (correctness against the CPU, absence of per-sector work, timings). The tables state the device and the mode
they were produced with; the JSON file next to each table also records the commit of the code.

\appendix
\section{Proof of Theorem~\ref{thm:allk}}\label{app:envelopes}

Throughout, $a=\ell=(1-\delta)/(1+\delta)$ and $b=\ell'=(1-\delta')/(1+\delta')$ are the smallest eigenvalues of the operators
with contrasts $\delta$ and $\delta'$; it suffices to prove the upper bounds in the complex algebra, since the
maximisers are real. A \emph{vertex argument} is used twice: a convex function on a face of a cube $[s,t]^n$ attains its
maximum at a vertex of the face, i.e.\ at a point with all coordinates in $\{s,t\}$.

\noindent\emph{Fact} (two projections; Halmos~\cite{halmos1969}). For projections $P,Q$ on a finite-dimensional
Hilbert space $H$, $H$ is the orthogonal sum of subspaces of dimension $1$ or $2$ invariant under $P$ and $Q$; on each
two-dimensional one there is an orthonormal basis $f,f'$ with $P=ff^*$ and $Q=qq^*$,
$q=\cos\theta\,f+e^{i\psi}\sin\theta\,f'$.

\begin{lemma}\label{lem:exp3}
Let $X=X^*$ be invertible, $P$ the projection onto the positive part of $X$, $S=2P-1$, and $Y\ge0$. In Exp.~3:
(a) if $|X|\le|X'|$, $X'=S|X'|$ and $|X'|$ commutes with $P$, then $z(X,Y)\le z(X',Y)$; (b) for fixed $X$, $z$ is
a convex function of $Y^2$.
\end{lemma}
\begin{proof}
With $A=YX_+Y$ and $B=YX_-Y$, $z=\Tr A+\Tr B-\|A-B\|_1$. (a) $A=Y|X|^{1/2}P|X|^{1/2}Y$ is non-decreasing in $|X|$ among
operators commuting with $P$, and so is $B$; by the triangle inequality $\|A+R-B\|_1\le\|A-B\|_1+\Tr R$ for $R\ge0$,
so $z$ does not decrease when $A$ or $B$ increases. (b) As $\|H\|_1=\max_{-1\le T\le1}\Tr(TH)$ for self-adjoint $H$, with
$T=1-2E$ we get $z=2\min_{0\le E\le1}\Tr\bigl(EA+(1-E)B\bigr)$. For invertible $Y$, $N=|X|^{1/2}Y$ is invertible and
$E\mapsto F=NEN^*$ maps $\{0\le E\le1\}$ onto $\{0\le F\le G\}$, $G=NN^*=|X|^{1/2}Y^2|X|^{1/2}$, with $\Tr(EA)=\Tr(PF)$ and
$\Tr((1-E)B)=\Tr((1-P)(G-F))$. So $z=2\min\{\Tr(PF+(1-P)(G-F)):0\le F\le G\}$ is the partial minimum of a linear
function over a convex set of pairs $(F,G)$, hence convex in $G$, which is affine in $Y^2$; by continuity also at
singular $Y$.
\end{proof}

\emph{Exp.~1.} By the trace-norm duality $|\Tr(XUY)|\le\|YX\|_1$, with equality for the polar unitary of $YX$, which
lies in $\M$; so $\sup_Uz=F(X,Y):=\|YX\|_1-\Tr(XY)$. For a fixed eigenbasis of $X$, $F$ is a convex function of the
eigenvalues (a norm of a linear map minus a linear function), and the admissible eigenvalues form the union of the
faces $\{x_i=1,\ x_j=a\}$ of $[a,1]^n$; by the vertex argument $F(X,Y)\le F(X',Y)$ for some $X'=a+(1-a)P$ with a
projection $P$ in $\M$, and likewise $Y'=b+(1-b)Q$. By the two-projection Fact, applied in every sector, $X'$, $Y'$ are
direct sums over subspaces of dimension $1$ or $2$, and so are $\|Y'X'\|_1$ and $\Tr(X'Y')$. One-dimensional summands
contribute $0$; on a two-dimensional one $X'=\mathrm{diag}(1,a)$ and $Y'$ has the eigenvalues $1,b$, so it contributes
at most $\delta^2\delta'^2/((1+\delta)(1+\delta'))$ by Proposition~\ref{prop:k2}(a); a sector of size $k_c$ holds at most
$\lfloor k_c/2\rfloor$ of them. Equality: $N$ blocks at the angle of Proposition~\ref{prop:k2}(a), $X=Y=1$ on the
remaining one-dimensional blocks, and the polar unitary.

\emph{Exp.~3.} Let $\delta<1$, so that $X$ is invertible; $z(-X,Y)=z(X,Y)$. Since $|X|$ commutes with $P$, it has a
unit eigenvector $e$ for its eigenvalue $a$ in one sector and in $\mathrm{ran}\,P$ or in $\ker P$; replacing $X$ by
$-X$ if necessary, $e\in\mathrm{ran}\,P$. Then $|X|\le|X'|:=1-(1-a)ee^*$, and $X'=S|X'|=S-(1-a)ee^*$ has contrast $\delta$;
by Lemma~\ref{lem:exp3}(a), $z(X,Y)\le z(X',Y)$. By Lemma~\ref{lem:exp3}(b) and the vertex argument for the eigenvalues
of $Y^2$ in $[b^2,1]$, $z(X',Y)\le z(X',Y')$ with $Y'=b+(1-b)Q$, $Q\ne0,1$ a projection in $\M$. Apply
the two-projection Fact to $P$ and $Q$ in every sector. In a sector without $e$, $X'=S$ is a direct sum over the
subspaces, a one-dimensional summand contributes $0$, and a two-dimensional one contributes at most $(1-b)^2$ by
Proposition~\ref{prop:k2}(b) ($|X|=1$ there): together at most $\lfloor k_c/2\rfloor(1-b)^2$. In the sector of $e$,
Lemma~\ref{lem:exp3}(a) with $|X'|\le1$ gives the same bound for $z$ restricted to it, and, if its size $k_c$ is odd or
it contains fewer than $k_c/2$ two-dimensional subspaces, at most $\lfloor k_c/2\rfloor(1-b)^2$, resp.\
$(k_c/2-1)(1-b)^2$. Otherwise its $m=k_c/2$ two-dimensional subspaces fill it, with bases $f_j,f_j'$, $f_j\in\mathrm{ran}\,P$;
the unitary that multiplies $f_j'$ by $e^{-i\psi_j}$ commutes with $X'$ (which is $-1$ on $\ker P$) and removes the
phases, and in the basis $(f_1,\dots,f_m,f_1',\dots,f_m')$ we are in the setting of Proposition~\ref{prop:mblocks} with
$X_+=1-(1-a)ee^*$, whose eigenvalues are $a,1,\dots,1$: the sector contributes at most
$(1-b)^2\bigl(\frac{2a}{1+a}+m-1\bigr)=(k_c/2-\delta)(1-b)^2$. Summing over the sectors, $z\le(N-\delta)(1-b)^2$ if every $k_c$ is
even and $z\le N(1-b)^2$ otherwise. Equality: if every $k_c$ is even, $X=\mathrm{diag}(a,-1)\oplus\mathrm{diag}(1,-1)^{\oplus(N-1)}$
and $Y$ with eigenvalues $1,b$ in every block at the angle of Proposition~\ref{prop:k2}(b); if some $k_c$ is odd,
$X=\mathrm{diag}(1,-1)$ on $N$ blocks, the eigenvalue $a$ on a one-dimensional block of an odd sector, $1$ on the other
one-dimensional blocks, and $Y$ as before ($Y=1$ on the one-dimensional blocks). \qed

% Appendix to the use case "quantitative failure of trace characterisations" (usecases/trace_characterizations.tex).
% \input at the end of main.tex. Numbers: generated/trace_characterizations.tex (macros \TC...).
\section{Trace characterisations: full table and proofs}\label{app:tc}

\begin{table}[h]
\centering\scriptsize
\setlength{\tabcolsep}{1.8pt}
\begin{tabular}{rlccccccrc}
\toprule
 & inequality $L\le R$ & ref. & $V(0.1)$ & $V(0.5)$ & $V(0.9)$ & $V^\star(0.5)$ & order & slope & $n$-dep.\\
\midrule
\IfFileExists{generated/trace_characterizations.tex}{\TCRows}{\multicolumn{10}{c}{\emph{run scripts/trace\_characterizations.py}}\\}
\bottomrule
\end{tabular}
\caption[]{The inequalities and their certified search values on $M_2$ at $\delta=0.1,0.5,0.9$. $N=R$ except
$N=\varphi(P)+\varphi(Q)$ in 9 and~21, $N=(\det A)^{1/n}$ in~11 and $N=\varphi(A+B)$ in~13. $V^\star$ is the sharp
value (Propositions~\ref{prop:tcK}, \ref{prop:tcB}, \ref{prop:tcP}, \ref{prop:tcL}); $^a$ marks a value proved sharp
for $n=2$ only, $^b$ a lower bound (Propositions~\ref{prop:tcC}, \ref{prop:tcW}), $^c$ a sharp value that is
approached but not attained. ``order'': leading order of $V^\star$ as $\delta\to0$; ``slope'':
$\log_2(V(0.2)/V(0.1))$ of the search; ``$n$-dep.'': the largest $|V_n/V_2-1|$ for $n=3,4$ and $\delta\ge0.1$; for
items 7 and 13 it reflects search shortfalls; items 19 and 20 differ for odd $n$ (Proposition~\ref{prop:tcP}). In~9, $p=1$ is
the strongest member of the family $\varphi(PQ+QP)\le2\varphi((QPQ)^p)$, $p\in(0,1]$. In~7, $r=2$ and $q=1$. In~8 we
use an extension of Young's inequality ($p=q=2$) to arbitrary $A,B$, modelled on Ando's trace form. In~11, the AM--GM
analogue. In 12 and~13, $f(t)=t^{1/2}$ (13: $f=g=t^{1/2}$); in 14, $\alpha=3$; in 15 and~16, the symmetric forms; in
24, $p=2$, $q=1$; in 25, $r=1$. ``n.r.'': no violation resolved by the search
\IfFileExists{generated/trace_characterizations.tex}{(\TCMode\ run on \TCDevice)}{}.}
\label{tab:tc}
\end{table}

\paragraph{Sources and forms.} We saw several of the cited papers only as abstracts or listings; the catalogue
\texttt{usecases/trace\_characterizations\_catalogue.md} records what we saw. The following forms are checked against
the full texts. Item~4 is~\cite[Theorem~1]{petz1988} with $k=2$; there $|\varphi(A^k)|\le\varphi(|A|^k)$ for all $A$ and
one fixed $k\ge1$ characterises the trace on $M_n$ ($k=1$ is Gardner's inequality, item~1). Item~7
is~\cite[Lemma~2]{bikchentaev2011pb} with $r=2$, $p=1$. Item~10 is the case $r=0$ of~\cite[Theorems~4, 5]{hoa2010},
which also state it on $M_n$~\cite{bikchentaev2007mono,sano2006}; item~25 is the case $r=1$, on $M_n$ due
to~\cite[Theorem~6]{hoa2007}. Item~11 is~\cite[Theorem~3(iii)]{bikchentaev2020det}, $\det(A)^{1/n}\le\frac1n\varphi(A)$
for $A\ge0$ and $\varphi(1)=n$, and item~24 is~\cite[Theorem~3(vii)]{bikchentaev2020det} with $p=2$, $q=1$; it needs
$\varphi$ faithful, which holds for $W>0$. Item~12 is~\cite[Corollary~1]{tikhonov2005} with the concave
$f(t)=t^{1/2}$. Item~13 is~\cite[Lemma~3.1 and Remark~4]{hoa2013} with $f(t)=t^{1/2}$ ($s=\frac12$ in their (10)).
Items 14, 15, 20, 21 and~22 are~\cite[Theorem~2(i), (iv), (v), (vi), (vii)]{bikchentaev2011mn}, item~14 with
$\alpha=3$. Item~16 is~\cite[Lemma~1]{bikchentaev2011pb}, stated there for positive $H,K$; both sides scale by $e^c$
under $H\mapsto H+c$, $K\mapsto K+c$, so the envelope over self-adjoint arguments is the same. Items 17, 18 and~19
are~\cite[Theorem~3.4(iv), (v), (iii)]{bikchentaev2010smj} (announced in~\cite{bikchentaev2009rm}), and item~23
is~\cite[Theorem~4]{petz1988}: a positive functional on a von Neumann algebra is tracial iff it is subadditive on the
projection lattice. In the proof of~\cite[Theorem~3.4]{bikchentaev2010smj}, (iv)$\Rightarrow$(vi), the density
$\mathrm{diag}(\frac12+s,\frac12-s)$ and $p=r_{(1,1/2+t)}$, $q=r_{(1,1/2-t)}$ give
$\varphi(pqp)-\varphi(q)=2t(2s-t-4st^2)$; the paper writes $h(t)=s-t-2st^2$, which has the same sign for $t<s/2$, so
the conclusion stands (checked in the tests). The published form of item~8 is
$\varphi(|AB|)\le\frac1p\varphi(A^p)+\frac1q\varphi(B^q)$ for $A,B\ge0$~\cite{bikchentaev2005young}; the table uses the
extension to arbitrary $A,B$, and the value $V=h-1$ of Proposition~\ref{prop:tcB}(a) refers to this extension. For
the published form with $p=q=2$, an independent search (L-BFGS from 25 random starts, $n=2,3$) gives only
$V\approx\kappa-1$: $0.1546$ at $\delta=0.5$ and $1.294$ at $\delta=0.9$, against $\kappa-1=0.1547$ and $1.294$. This
is an observation, not a proof; the linear failure of item~8 is a feature of the extension.

\subsection*{Proof of Proposition~\ref{prop:tcP}}

\emph{Step 1: reduction to rank-one projections on $M_2$.} By Lemma~\ref{lem:twoproj} (see also~\cite{halmos1969}),
$\C^n=\bigoplus_jH_j$ with $\dim H_j\in\{1,2\}$, every $H_j$ invariant under $P$ and $Q$. Every operator in $L$, $R$
and $N$ (namely $PQP$, $Q$, $|P-Q|$, $P+Q$, $1+P\wedge Q$, $e^{QPQ}$, $e^P$, $\mathrm{Re}\,PQ$, $|PQ|$, $P\vee Q$)
is then block diagonal, with block $F(P_j,Q_j)$; for the lattice operations because $\mathrm{ran}\,P\cap\mathrm{ran}\,Q
=\bigoplus_j(\mathrm{ran}\,P_j\cap\mathrm{ran}\,Q_j)$, the block projections commuting with $P$ and $Q$. Hence
$L=\sum_jL_j$ etc., where $L_j$ is the same expression for $P_j,Q_j$ and the functional $\Tr(W_j\,\cdot\,)$ on
$H_j$, with the compression $W_j$ of $W$; its spectrum lies in $[b,a]$, so $\Delta(W_j)\le\delta$. On a
one-dimensional $H_j$, $P_j,Q_j\in\{0,1\}$ commute and $L_j\le R_j$ in all seven items (e.g.\ $\max(p,q)\le p+q$,
$p+q\le1+pq$, $e^{pq}\le e^p$). On a two-dimensional $H_j$, $P_j\ne Q_j$ are rank-one projections. Since $N=\sum_jN_j$
with $N_j\ge0$, and $L_j\le R_j$ wherever $N_j=0$, we get $(L-R)/N\le\max_j(L_j-R_j)/N_j$, and the ratio does not
change when $W_j$ is scaled. It therefore suffices to bound $(L-R)/N$ for $W=\frac12(1+\tau\sigma_3)$, $0\le\tau\le\delta$,
and rank-one $P=\frac12(1+\mu\cdot\sigma)$, $Q=\frac12(1+\nu\cdot\sigma)$ with unit vectors $\mu\ne\nu$ in $\mathbb R^3$
(Bloch vectors). Moreover, for fixed $\mu,\nu$ the ratio is a linear-fractional function of $\tau$ (all three
quantities are affine in $W$), with positive denominator, hence monotone on $[-\delta,\delta]$; and conjugation
with $\sigma_1$ maps $\tau$ to $-\tau$ and $(\mu,\nu)$ to reflected vectors. So the bound for $\tau=\delta$ covers all
$\tau\le\delta$, and the $M_2$ envelope bounds the envelope in every dimension. In 17, 18 and 21--23, $N$ vanishes
on blocks where $P=Q=0$, so the values below, attained (or approached) on $M_2$, are attained on $M_n$ with $P,Q$
supported on the span of $e_1,e_2$. In 19 and 20, $N$ contains $\varphi(1)$. For even $n$ take $W$ with the
eigenvalues $a,b$, each $n/2$ times, and $n/2$ copies of the $M_2$ extremiser, one in each block $\mathrm{span}
\{e_a,e_b\}$; this attains the $M_2$ value. For odd $n$ at least one one-dimensional block is left over. In 19, a
two-dimensional block contributes $L_j-R_j<\lambda_1(W_j)-\lambda_2(W_j)$ (approached) and $N_j=\Tr W_j$, and a
one-dimensional one $L_j-R_j\le0$; by Ky Fan's maximum principle the sums of the $\lambda_1(W_j)$ and of the
$\lambda_2(W_j)$ over $k'\le k=(n-1)/2$ blocks are at most the sum of the $k'$ largest and at least the sum of the $k'$
smallest eigenvalues of $W$, so $(L-R)/N<k(a-b)/(k(a+b)+b)=(n-1)\delta/(n-\delta)$, approached for the spectrum
$a^{(k)},b^{(k+1)}$ with $P=1$, $Q=0$ on the remaining eigenvector. In~20 the same configuration, with $P=Q=0$ on the
remaining eigenvector, adds $b/k$ per block to $N$, i.e.\ $(1-\delta)/(n-1)$ to $\beta$ below (normalisation
$\Tr W_j=1$); this is a lower bound, and the $M_2$ value an upper bound.
\IfFileExists{generated/trace_characterizations.tex}{For $n=3$, $\delta=0.5$ the two are $\TCExpOddHalf$ and
$\TCRefExpqpqHalf$; the search on $M_3$ finds the former.}{}

\emph{Step 2: formulas on $M_2$.} Write $\varphi(x_0+x\cdot\sigma)=x_0+\delta x_3$, $\mu\cdot\nu=\cos2\theta$
($0<\theta\le\pi/2$) and $c=\cos\theta=|\langle u,v\rangle|$ for unit vectors $u,v$ with $P=uu^*$, $Q=vv^*$. Then
$PQP=c^2P$, $(P-Q)^2=\sin^2\theta\cdot1$, $P\wedge Q=0$, $P\vee Q=1$, $\mathrm{Re}\,PQ=\frac14(1+\mu\cdot\nu+(\mu+\nu)\cdot
\sigma)$, $|PQ|=(QPQ)^{1/2}=cQ$, $e^{QPQ}=1+(e^{c^2}-1)Q$ and $e^P=1+(e-1)P$, and $|\mu+\nu|=2\cos\theta$. We use:
\emph{for $\beta>|z|$, $\max_{|\xi|=1}(\alpha+y\cdot\xi)/(\beta+z\cdot\xi)$ is the unique zero $\lambda$ of
$\alpha-\lambda\beta+|y-\lambda z|$}, since $(\alpha+y\cdot\xi)-\lambda(\beta+z\cdot\xi)\le\alpha-\lambda\beta+|y-\lambda z|$
with equality for $\xi\parallel y-\lambda z$, and the right side decreases strictly in $\lambda$.

(17) $L/R=(1+\mu\cdot\nu)(1+\delta\mu_3)/(2(1+\delta\nu_3))$. For fixed $\mu$, the lemma with $\alpha=\beta=1$, $y=\mu$,
$z=\delta e_3$ gives $\max_\nu(1+\mu\cdot\nu)/(1+\delta\nu_3)=2(1-\delta\mu_3)/(1-\delta^2)$, so
$L/R\le(1-\delta^2\mu_3^2)/(1-\delta^2)\le\kappa^2$, with equality for $\mu_3=0$ and $\nu=(\mu-\lambda\delta e_3)/(\lambda-1)$,
$\lambda=2/(1-\delta^2)$.

(18) $L/R=\sin\theta/(1+\frac\delta2(\mu_3+\nu_3))\le\sin\theta/(1-\delta\cos\theta)\le\kappa$, the maximum at
$\cos\theta=\delta$, $\mu+\nu\parallel-e_3$. (Also $\varphi(|P-Q|)\le\kappa\varphi(P+Q)$ is the triangle inequality of
Proposition~\ref{prop:tcK} with $X=P$, $Y=-Q$.)

(19) $(L-R)/N=\varphi(P+Q)-1=\frac\delta2(\mu_3+\nu_3)<\delta$, approached as $\mu,\nu\to e_3$.

(20) For fixed $g=2\theta$, the pair $(\mu_3,\nu_3)$ ranges over the image of the unit sphere under
$\xi\mapsto(\mu\cdot\xi,\nu\cdot\xi)$, a filled ellipse, and the linear-fractional ratio is maximal on its boundary,
$\mu_3=\cos x$, $\nu_3=\cos(x+g)$. With $K=e-1$ and $E=e^{\cos^2(g/2)}-1$,
\[
\frac{L-R}N=\frac{\alpha+y\cdot\xi_x}{\beta+z\cdot\xi_x},\quad \alpha=\tfrac{E-K}2,\ y=\tfrac\delta2(E\cos g-K,\,-E\sin g),
\ \beta=1+\tfrac K2,\ z=(\tfrac{\delta K}2,0),
\]
$\xi_x=(\cos x,\sin x)$; by the lemma its maximum $\lambda(\delta,g)$ is the larger root of
$(\beta^2-|z|^2)\lambda^2-2(\alpha\beta-y\cdot z)\lambda+\alpha^2-|y|^2=0$. For small $\delta$ the maximiser has
$g=2K\delta/e+O(\delta^2)$, $x=\pi/2$, and $L-R=-\frac e8g^2+\frac{K\delta}2g+O(\delta^3)$, which gives
$V=K^2\delta^2/(e(e+1))+O(\delta^3)$.

(21) With $s=\mu_3+\nu_3\ge-2\cos\theta$, $(L-R)/N=-(2\cos^2\theta+\delta s)/(4(1+\frac\delta2s))$, which decreases in $s$,
so $(L-R)/N\le(\delta C-C^2)/(2(1-\delta C))$, $C=\cos\theta$. The maximum over $C$ is at
$\delta C^2-2C+\delta=0$, $C=\tan(\psi/2)$, with value $C^2/2=\frac12\tan^2(\psi/2)$.

(22) $\varphi(\mathrm{Re}\,PQ)=\mathrm{Re}\,\varphi(PQ)\le|\varphi(PQ)|\le\kappa\varphi(|PQ|)$ by
Proposition~\ref{prop:tcK}, in every dimension. Equality: $v=\sqrt{b/(a+b)}\,e_1+\sqrt{a/(a+b)}\,e_2$, $u=Wv/\|Wv\|$;
then $PQ=\langle u,v\rangle uv^*$ has the polar factor of Proposition~\ref{prop:tcK}.

(23) $L/R=1/(1+\frac\delta2(\mu_3+\nu_3))<1/(1-\delta)$, approached as $\mu,\nu\to-e_3$. \qed

\subsection*{Proof of Proposition~\ref{prop:tcL}}
For $A\ge0$, $\varphi(A^2)=\Tr(A^{1/2}WA^{1/2}A)\le\|A\|\varphi(A)\le\Tr(A)\varphi(A)\le\varphi(A)^2/w_{\min}$, with
equality for the projection onto an eigenvector for $w_{\min}$. With $\Tr W=n$ and the extreme eigenvalues in the
ratio $h$, $w_{\min}=nb/\sum_iw_i\ge nb/(b+(n-1)a)=n\ell/(\ell+n-1)$, with equality for the spectrum
$\{a,\dots,a,b\}$. For $n=2$, $(\ell+1)/(2\ell)=1/(1-\delta)$. \qed

\subsection*{Proof of Proposition~\ref{prop:tcW}}
If $A=\alpha uu^*$ is of rank one and $A\le B$, then $u^*Bu\ge\alpha$ and $AB^2A=\alpha^2\|Bu\|^2uu^*\ge\alpha^2(u^*Bu)^2
uu^*\ge A^4$, so every positive $\varphi$ satisfies the inequality. For the construction, take $W$ with $\Tr W=1$,
eigenvalues $(1\pm\delta)/2$, and, in an orthonormal basis at the angle $\vartheta$ to its eigenbasis,
$A=\mathrm{diag}(1,a)$ and $B=A+tvv^*$ with $v=(\cos\phi,\sin\phi)$; there $W$ has the entries
$w_{11,22}=\frac12(1\pm\delta\cos2\vartheta)$ and $w_{12}=\frac\delta2\sin2\vartheta$. Then $R-L=tT_1+t^2T_2$ with
$T_1=\varphi(A^2PA+APA^2)$, $T_2=\varphi(APA)$, $P=vv^*$, and the best $t=-T_1/(2T_2)$ gives $(L-R)/R=G/(L-G)$,
$G=T_1^2/(4T_2)$, whenever $T_1<0$. Let $m_{25}(\delta)$ be the maximum over $(\vartheta,a,\phi)$; it is attained,
and by Step~1 above (applied to $\C^2\subset\C^n$) it is a lower bound in every dimension. For the leading order put
$\vartheta=\pi/4$, $\sqrt a=\sigma\delta$ and $\tan\phi=-u\,a^{-3/2}$. Then
$T_1=\cos^2\phi\,(1+u^2-(1+a)u/\sigma)$, $T_2=\cos^2\phi\,(\frac12+\frac{u^2}{2a}-\frac u\sigma)$ and
$\cos^2\phi=a^3/(a^3+u^2)$, so that
\[
\frac{L-R}R=\bigl(\sigma^2(s-s^2)-\sigma^4\bigr)^2\delta^8+O(\delta^{10}),\qquad s=\sigma/u,
\]
which is maximal for $s=\frac12$, $\sigma^2=\frac18$: $m_{25}(\delta)\ge\delta^8/4096+O(\delta^{10})$. The numerical
maximisation gives $m_{25}(0.1)/(0.1^8/4096)=\IfFileExists{generated/trace_characterizations.tex}{\TCWmonLead}{1.02}$.
The same scaling for general $r\ge0$ gives the order $\delta^{4(r+1)}$; for $r=0$ the supremum of its leading
coefficient is $\frac1{16}$, that of item~10. \qed

\paragraph{Items 12--16.} These have no closed form so far. For~14 a short expansion gives the leading order.
Write $W=\frac12(1+\delta D)$ with $D=\mathrm{diag}(1,-1)$, take real unit vectors $u,v$ at angles $m\pm\varepsilon/2$
in the eigenbasis of $W$, and set $P=uu^*$, $Q=vv^*$. Then
$L-R=\frac12(-\varepsilon^2-\delta(1-\frac2p)\varepsilon\sin2m)+O(\delta^3)$ for $\varepsilon=O(\delta)$.
Hence $V_{14}\ge(1-\frac2p)^2\delta^2/4+O(\delta^3)$, which is $\delta^2/36$ for $p=3$, while for $p=2$
every positive $\varphi$ satisfies the inequality, in agreement with the remark after~\cite[Theorem~2]{bikchentaev2011mn}
that the case $\alpha=2$ does not characterise traces. For item~13 the natural normalisation $N=R$
would give $V=\infty$ for every $\delta>0$: for rank-one $A=uu^*$, $B=vv^*$ with overlap $c\to0$,
$L\approx2c\,|u^*Wu_\perp|$ while $R=2c^2u^*Wu$. We therefore take $N=\varphi(A+B)$. The search gives:
\begin{itemize}
\item all five items fail quadratically as $\delta\to0$, with $V\approx0.5\,\delta^2$ (12, 13), $\delta^2/36$ (14),
$0.064\,\delta^2$ (15) and $0.09\,\delta^2$ (16);
\item item~15 decreases with $n$, because the identity part of $e^{P/2}e^Qe^{P/2}$ enters $N$; the others give the
same values for $n=2,3,4$, up to search shortfalls of item~13 at $\delta\ge0.9$;
\item in item~16 the best points have large $\|A\|$, $\|B\|$, and the supremum is approached as the scale grows.
\end{itemize}
An independent SciPy search on $M_2$ (plain NumPy, BFGS from
\IfFileExists{generated/trace_characterizations.tex}{\TCSciStarts}{many} random starts, at $\delta=0.3,0.5,0.9$)
gives values whose ratio to the library's lies in
\IfFileExists{generated/trace_characterizations.tex}{$[\TCSciRatioMin,\TCSciRatioMax]$}{(not generated)}. Ratios
above~1 come from item~13, where BFGS stops earlier than the batched search.

\paragraph{Items 17--25 in the search.} For the projection items the search on $M_n$, $n=2,3,4$, reproduces the values
of Proposition~\ref{prop:tcP}, including the odd-$n$ values of 19 and~20
\IfFileExists{generated/trace_characterizations.tex}{(for 17, 18 and 21--23 on $M_3$ and $M_4$ to $\TCNDevNew$ of the
$M_2$ values)}{}. In 19 and~23 the suprema are approached as $P$ and $Q$ merge, where $P\wedge Q$ and $P\vee Q$ jump;
the search therefore maximises $(L-R)/N-10^{-12}/\|P-Q\|_{\mathrm F}^2$, which keeps $P$ and $Q$ about $10^{-3}$ apart,
while the reported and certified values are those of $(L-R)/N$. The search resolves item~25 for
\IfFileExists{generated/trace_characterizations.tex}{$\delta\ge\TCWmonFirst$, reaches $\TCWmonRatioHalf$ of $m_{25}$ at
$\delta=0.5$ and $m_{25}$ itself (to $10^{-3}$) for $\delta\ge\TCWmonReach$; $m_{25}(0.9)=\TCRefWmonNine$, search
$\TCVWmonNine$.}{large $\delta$.} A BFGS search over $A,C\ge0$ on $M_2$ (tests) also reaches $m_{25}(0.9)$ and does not
exceed it.


{\sloppy
\begin{thebibliography}{99}
\bibitem{v1} I.~Nikolaeva, A.~Novikov, ``Finite-dimensional Type I von Neumann algebras in PyTorch: a
  GPU-accelerated framework for random block-diagonal operators'', arXiv:2606.15882v1 (2026).
\bibitem{gardner1979} L.~T.~Gardner, ``An inequality characterizes the trace'', Canad.\ J.\ Math.\ 31
  (1979), no.~6, 1322--1328.
\bibitem{petz1988} D.~Petz, J.~Zem\'anek, ``Characterizations of the trace'', Linear Algebra Appl.\ 111
  (1988), 43--52.
\bibitem{novikov2015} A.~A.~Novikov, O.~E.~Tikhonov, ``Characterization of central elements of operator
  algebras by inequalities'', Lobachevskii J.\ Math.\ 36 (2015), no.~2, 208--210; arXiv:1502.01267.
\bibitem{virosztek2017} D.~Virosztek, ``Connections between centrality and local monotonicity of certain
  functions on C*-algebras'', J.\ Math.\ Anal.\ Appl.\ 453 (2017), 221--226.
\bibitem{abed2024} S.~A.~Abed, I.~A.~Nikolaeva, A.~A.~Novikov, ``Generalisation of Michelson contrast for
  operators and its properties'', Lobachevskii J.\ Math.\ 45 (2024), no.~8, 3835--3848; arXiv:2010.10130.
\bibitem{johansson2013} J.~R.~Johansson, P.~D.~Nation, F.~Nori, ``QuTiP 2: A Python framework for the
  dynamics of open quantum systems'', Comput.\ Phys.\ Commun.\ 184 (2013), 1234--1240.
\bibitem{lambert2026} N.~Lambert et al., ``QuTiP~5: The Quantum Toolbox in Python'', Phys.\ Rep.\ 1153 (2026),
  1--62; arXiv:2412.04705.
\bibitem{qiskit2019} A.~Javadi-Abhari et al., ``Quantum computing with Qiskit'', arXiv:2405.08810 (2024).
\bibitem{dynamiqs} P.~Guilmin, A.~Bocquet, \'E.~Genois, D.~Weiss, R.~Gautier, ``Dynamiqs: an open-source Python
  library for GPU-accelerated and differentiable simulation of quantum systems'', in preparation (2025),
  \url{https://github.com/dynamiqs/dynamiqs}.
\bibitem{mezzadri2007} F.~Mezzadri, ``How to generate random matrices from the classical compact
  groups'', Notices Amer.\ Math.\ Soc.\ 54 (2007), 592--604.
\bibitem{mehta2004} M.~L.~Mehta, \emph{Random Matrices}, 3rd ed., Elsevier, 2004.
\bibitem{page1993} D.~N.~Page, ``Average entropy of a subsystem'', Phys.\ Rev.\ Lett.\ 71 (1993),
  1291--1294.
\bibitem{foong1994} S.~K.~Foong, S.~Kanno, ``Proof of Page's conjecture on the average entropy of a subsystem'',
  Phys.\ Rev.\ Lett.\ 72 (1994), 1148.
\bibitem{sanchezruiz1995} J.~S\'anchez-Ruiz, ``Simple proof of Page's conjecture on the average entropy of a
  subsystem'', Phys.\ Rev.\ E 52 (1995), 5653--5655.
\bibitem{lubkin1978} E.~Lubkin, J.\ Math.\ Phys.\ 19 (1978), 1028.
\bibitem{collins2006} B.~Collins, P.~\'Sniady, ``Integration with respect to the Haar measure on unitary,
  orthogonal and symplectic group'', Commun.\ Math.\ Phys.\ 264 (2006), 773--795.
\bibitem{zyczkowski2001} K.~\.Zyczkowski, H.-J.~Sommers, ``Induced measures in the space of mixed quantum
  states'', J.\ Phys.\ A 34 (2001), 7111--7125.
\bibitem{peschel2003} I.~Peschel, ``Calculation of reduced density matrices from correlation functions'',
  J.\ Phys.\ A 36 (2003), L205--L208.
\bibitem{oganesyan2007} V.~Oganesyan, D.~A.~Huse, ``Localization of interacting fermions at high temperature'',
  Phys.\ Rev.\ B 75 (2007), 155111.
\bibitem{atas2013} Y.~Y.~Atas, E.~Bogomolny, O.~Giraud, G.~Roux, ``Distribution of the ratio of consecutive
  level spacings in random matrix ensembles'', Phys.\ Rev.\ Lett.\ 110 (2013), 084101.
\bibitem{takesaki2002} M.~Takesaki, \emph{Theory of Operator Algebras I}, Springer, 2002.
\bibitem{wick1952} G.~C.~Wick, A.~S.~Wightman, E.~P.~Wigner, ``The intrinsic parity of elementary particles'',
  Phys.\ Rev.\ 88 (1952), 101--105.
\bibitem{bartlett2007} S.~D.~Bartlett, T.~Rudolph, R.~W.~Spekkens, ``Reference frames, superselection rules, and
  quantum information'', Rev.\ Mod.\ Phys.\ 79 (2007), 555--609.
\bibitem{zurek2003} W.~H.~Zurek, ``Decoherence, einselection, and the quantum origins of the classical'',
  Rev.\ Mod.\ Phys.\ 75 (2003), 715--775.
\bibitem{sandvik2010} A.~W.~Sandvik, ``Computational studies of quantum spin systems'', AIP Conf.\ Proc.\ 1297
  (2010), 135--338.
\bibitem{weinberg2017} P.~Weinberg, M.~Bukov, ``QuSpin: a Python package for dynamics and exact diagonalisation of
  quantum many body systems part I: spin chains'', SciPost Phys.\ 2 (2017), 003.
\bibitem{goldstein2018} M.~Goldstein, E.~Sela, ``Symmetry-resolved entanglement in many-body systems'',
  Phys.\ Rev.\ Lett.\ 120 (2018), 200602.
\bibitem{xavier2018} J.~C.~Xavier, F.~C.~Alcaraz, G.~Sierra, ``Equipartition of the entanglement entropy'',
  Phys.\ Rev.\ B 98 (2018), 041106.
\bibitem{zanardi1997} P.~Zanardi, M.~Rasetti, ``Noiseless quantum codes'', Phys.\ Rev.\ Lett.\ 79 (1997), 3306--3309.
\bibitem{knill2000} E.~Knill, R.~Laflamme, L.~Viola, ``Theory of quantum error correction for general noise'',
  Phys.\ Rev.\ Lett.\ 84 (2000), 2525--2528.
\bibitem{beny2007} C.~B\'eny, A.~Kempf, D.~W.~Kribs, ``Generalization of quantum error correction via the Heisenberg
  picture'', Phys.\ Rev.\ Lett.\ 98 (2007), 100502.
\bibitem{hadfield2019} S.~Hadfield, Z.~Wang, B.~O'Gorman, E.~G.~Rieffel, D.~Venturelli, R.~Biswas, ``From the quantum
  approximate optimization algorithm to a quantum alternating operator ansatz'', Algorithms 12 (2019), 34.
\bibitem{dyson1962} F.~J.~Dyson, ``The threefold way. Algebraic structure of symmetry groups and ensembles in quantum
  mechanics'', J.\ Math.\ Phys.\ 3 (1962), 1199--1215.
\bibitem{cohen2016} T.~Cohen, M.~Welling, ``Group equivariant convolutional networks'', Proc.\ ICML 2016, PMLR 48,
  2990--2999.
\bibitem{finzi2021} M.~Finzi, M.~Welling, A.~G.~Wilson, ``A practical method for constructing equivariant multilayer
  perceptrons for arbitrary matrix groups'', Proc.\ ICML 2021, PMLR 139, 3318--3328.
\bibitem{martin2021} C.~H.~Martin, M.~W.~Mahoney, ``Implicit self-regularization in deep neural networks: evidence from
  random matrix theory and implications for learning'', J.\ Mach.\ Learn.\ Res.\ 22 (2021), 1--73.
\bibitem{harris2020} C.~R.~Harris et al., ``Array programming with NumPy'', Nature 585 (2020), 357--362.
\bibitem{virtanen2020} P.~Virtanen et al., ``SciPy 1.0: fundamental algorithms for scientific computing in Python'',
  Nat.\ Methods 17 (2020), 261--272.
\bibitem{lehoucq1998} R.~B.~Lehoucq, D.~C.~Sorensen, C.~Yang, \emph{ARPACK Users' Guide}, SIAM, 1998.
\bibitem{fishman2022} M.~Fishman, S.~R.~White, E.~M.~Stoudenmire, ``The ITensor software library for tensor network
  calculations'', SciPost Phys.\ Codebases 4 (2022).
\bibitem{rams2025} M.~M.~Rams, G.~W\'ojtowicz, A.~Sinha, J.~Hasik, ``YASTN: Yet another symmetric tensor networks; A
  Python library for Abelian symmetric tensor network calculations'', SciPost Phys.\ Codebases 52 (2025);
  arXiv:2405.12196.
\bibitem{hauschild2018} J.~Hauschild, F.~Pollmann, ``Efficient numerical simulations with tensor networks: Tensor
  Network Python (TeNPy)'', SciPost Phys.\ Lect.\ Notes 5 (2018).
\bibitem{fukuda2019} M.~Fukuda, R.~K\"onig, I.~Nechita, ``RTNI -- a symbolic integrator for Haar-random tensor
  networks'', J.\ Phys.\ A 52 (2019), 425303.
\bibitem{higham2008} N.~J.~Higham, \emph{Functions of Matrices: Theory and Computation}, SIAM, 2008.
\bibitem{trefethen2014} L.~N.~Trefethen, J.~A.~C.~Weideman, ``The exponentially convergent trapezoidal rule'',
SIAM Rev. 56 (2014) 385--458.
\bibitem{daleckii1965} Ju.~L.~Daleckii, S.~G.~Krein, ``Integration and differentiation of functions of Hermitian
operators and applications to the theory of perturbations'', Amer. Math. Soc. Transl. (2) 47 (1965) 1--30.
\bibitem{bhatia1997} R.~Bhatia, \emph{Matrix Analysis}, Graduate Texts in Mathematics 169, Springer, 1997.
\bibitem{kasim2020} M.~F.~Kasim, ``Derivatives of partial eigendecomposition of a real symmetric matrix for
degenerate cases'', arXiv:2011.04366 (2020).
\bibitem{petz1986} D.~Petz, ``Quasi-entropies for finite quantum systems'', Rep.\ Math.\ Phys.\ 23 (1986), 57--65.
\bibitem{petz1986suff} D.~Petz, ``Sufficient subalgebras and the relative entropy of states of a von Neumann
  algebra'', Commun.\ Math.\ Phys.\ 105 (1986), 123--131.
\bibitem{petz1988chan} D.~Petz, ``Sufficiency of channels over von Neumann algebras'', Quart.\ J.\ Math.\ Oxford 39
  (1988), no.~1, 97--108.
\bibitem{mullerlennert2013} M.~M\"uller-Lennert, F.~Dupuis, O.~Szehr, S.~Fehr, M.~Tomamichel, ``On quantum R\'enyi
  entropies: a new generalization and some properties'', J.\ Math.\ Phys.\ 54 (2013), 122203.
\bibitem{wilde2014} M.~M.~Wilde, A.~Winter, D.~Yang, ``Strong converse for the classical capacity of
  entanglement-breaking and Hadamard channels via a sandwiched R\'enyi relative entropy'', Commun.\ Math.\ Phys.\ 331
  (2014), 593--622.
\bibitem{frank2013} R.~L.~Frank, E.~H.~Lieb, ``Monotonicity of a relative R\'enyi entropy'', J.\ Math.\ Phys.\ 54
  (2013), 122201.
\bibitem{beigi2013} S.~Beigi, ``Sandwiched R\'enyi divergence satisfies data processing inequality'', J.\ Math.\
  Phys.\ 54 (2013), 122202.
\bibitem{berta2017} M.~Berta, O.~Fawzi, M.~Tomamichel, ``On variational expressions for quantum relative
  entropies'', Lett.\ Math.\ Phys.\ 107 (2017), 2239--2265.
\bibitem{zhang2020} H.~Zhang, ``From Wigner--Yanase--Dyson conjecture to Carlen--Frank--Lieb conjecture'', Adv.\
  Math.\ 365 (2020), 107053.
\bibitem{junge2018} M.~Junge, R.~Renner, D.~Sutter, M.~M.~Wilde, A.~Winter, ``Universal recovery maps and approximate
  sufficiency of quantum relative entropy'', Ann.\ Henri Poincar\'e 19 (2018), 2955--2978.
\bibitem{jones1983} V.~F.~R.~Jones, ``Index for subfactors'', Invent.\ Math.\ 72 (1983), 1--25.
\bibitem{ghj1989} F.~M.~Goodman, P.~de~la~Harpe, V.~F.~R.~Jones, \emph{Coxeter Graphs and Towers of Algebras}, MSRI
  Publications 14, Springer, 1989.
\bibitem{pimsner1986} M.~Pimsner, S.~Popa, ``Entropy and index for subfactors'', Ann.\ Sci.\ \'Ecole Norm.\ Sup.\ (4)
  19 (1986), 57--106.
\bibitem{connes1975} A.~Connes, E.~St{\o}rmer, ``Entropy for automorphisms of II$_1$ von Neumann algebras'', Acta Math.\
  134 (1975), 289--306.
\bibitem{choda} M.~Choda, ``Relative entropy for maximal abelian subalgebras of matrices and the entropy of
  unistochastic matrices'', Internat.\ J.\ Math.\ 19 (2008), 767--776; arXiv:0803.0591.
\bibitem{bakshi} K.~C.~Bakshi, S.~Guin, ``Relative position between a pair of spin model subfactors'', J.\ Aust.\
  Math.\ Soc.\ 119 (2025), no.~1, 1--38; arXiv:2212.07197.
\bibitem{petz2009} D.~Petz, A.~Sz\'ant\'o, M.~Weiner, ``Complementarity and the algebraic structure of 4-level
  quantum systems'', Infin.\ Dimens.\ Anal.\ Quantum Probab.\ Relat.\ Top.\ 12 (2009), 99--116; arXiv:0802.2391.
\bibitem{ghosh2026} I.~Ghosh, S.~Kumar, ``Kadison--Kastler distance and interior angle for masas: connections with
  entropy and probabilistic index'', arXiv:2609.38206 (2026).
\bibitem{greub1959} W.~Greub, W.~Rheinboldt, ``On a generalization of an inequality of L.~V.~Kantorovich'', Proc.\
  Amer.\ Math.\ Soc.\ 10 (1959), 407--415.
\bibitem{furuta2005} T.~Furuta, J.~Mi\'ci\'c Hot, J.~Pe\v{c}ari\'c, Y.~Seo, \emph{Mond--Pe\v{c}ari\'c Method in Operator
  Inequalities}, Element, Zagreb, 2005.
\bibitem{watatani1990} Y.~Watatani, ``Index for $C^*$-subalgebras'', Mem.\ Amer.\ Math.\ Soc.\ 83 (1990), no.~424.
\bibitem{frank1998} M.~Frank, E.~Kirchberg, ``On conditional expectations of finite index'', J.\ Operator Theory 40
  (1998), 87--111.
\bibitem{bdh1988} M.~Baillet, Y.~Denizeau, J.-F.~Havet, ``Indice d'une esp\'erance conditionnelle'', Compositio
  Math.\ 66 (1988), 199--236.
\bibitem{murota2010} K.~Murota, Y.~Kanno, M.~Kojima, S.~Kojima, ``A numerical algorithm for block-diagonal
  decomposition of matrix $*$-algebras with application to semidefinite programming'', Japan J.\ Indust.\ Appl.\
  Math.\ 27 (2010), 125--160.
\bibitem{maehara2010} T.~Maehara, K.~Murota, ``A numerical algorithm for block-diagonal decomposition of matrix
  $*$-algebras with general irreducible components'', Japan J.\ Indust.\ Appl.\ Math.\ 27 (2010), 263--293.
\bibitem{maehara2011} T.~Maehara, K.~Murota, ``Algorithm for error-controlled simultaneous block-diagonalization of
  matrices'', SIAM J.\ Matrix Anal.\ Appl.\ 32 (2011), 605--620.
\bibitem{fulton1991} W.~Fulton, J.~Harris, \emph{Representation Theory: A First Course}, Graduate Texts in
  Mathematics 129, Springer, 1991.
\bibitem{racah1942} G.~Racah, ``Theory of complex spectra.~II'', Phys.\ Rev.\ 62 (1942), 438--462.
\bibitem{edmonds1957} A.~R.~Edmonds, \emph{Angular Momentum in Quantum Mechanics}, Princeton University Press, 1957.
\bibitem{lieb1962} E.~Lieb, D.~Mattis, ``Ordering energy levels of interacting spin systems'', J.\ Math.\ Phys.\ 3
  (1962), 749--751.
\bibitem{majumdar1969} C.~K.~Majumdar, D.~K.~Ghosh, ``On next-nearest-neighbor interaction in linear chain.~I'',
  J.\ Math.\ Phys.\ 10 (1969), 1388--1398.
\bibitem{majumdar1969b} C.~K.~Majumdar, D.~K.~Ghosh, ``On next-nearest-neighbor interaction in linear chain.~II'',
  J.\ Math.\ Phys.\ 10 (1969), 1399--1402.
\bibitem{paszke2019} A.~Paszke et al., ``PyTorch: An imperative style, high-performance deep learning
  library'', NeurIPS 2019.
\bibitem{iemini2018} F.~Iemini, A.~Russomanno, J.~Keeling, M.~Schir\`o, M.~Dalmonte, R.~Fazio, ``Boundary time
  crystals'', Phys.\ Rev.\ Lett.\ 121 (2018), 035301.
\bibitem{carollo2021} F.~Carollo, I.~Lesanovsky, ``Exactness of mean-field equations for open Dicke models with an
  application to pattern retrieval dynamics'', Phys.\ Rev.\ Lett.\ 126 (2021), 230601.
\bibitem{piccitto2021} G.~Piccitto, M.~Wauters, F.~Nori, N.~Shammah, ``Symmetries and conserved quantities of
  boundary time crystals in generalized spin models'', Phys.\ Rev.\ B 104 (2021), 014307.
\bibitem{liu2026} Z.~Liu, Y.~Li, Z.~Fei, X.~Wang, ``Boundary time crystals beyond mean-field theory: a stroboscopic
  rotating-wave approximation'', Phys.\ Rev.\ A 114 (2026), 012211; arXiv:2510.03028.
\bibitem{nemeth2025} D.~Nemeth, A.~Principi, A.~Nazir, ``Analytical solution of boundary time crystals via the
  superspin basis'', arXiv:2507.06998 (2025).
\bibitem{wang2025} Z.~Wang, R.~Gao, X.~Wu, B.~Bu\v{c}a, K.~M{\o}lmer, L.~You, F.~Yang, ``Boundary time crystals
  induced by local dissipation and long-range interactions'', Phys.\ Rev.\ Lett.\ 135 (2025), 230401.
\bibitem{tucker2018} K.~Tucker, B.~Zhu, R.~J.~Lewis-Swan, J.~Marino, F.~Jimenez, J.~G.~Restrepo, A.~M.~Rey,
  ``Shattered time: can a dissipative time crystal survive many-body correlations?'', New J.\ Phys.\ 20 (2018),
  123003.
\bibitem{prazeres2021} L.~F.~dos~Prazeres, L.~da~Silva~Souza, F.~Iemini, ``Boundary time crystals in collective
  $d$-level systems'', Phys.\ Rev.\ B 103 (2021), 184308.
\bibitem{chase2008} B.~A.~Chase, J.~M.~Geremia, ``Collective processes of an ensemble of spin-1/2 particles'',
  Phys.\ Rev.\ A 78 (2008), 052101.
\bibitem{baragiola2010} B.~Q.~Baragiola, B.~A.~Chase, J.~Geremia, ``Collective uncertainty in partially polarized
  and partially decohered spin-1/2 systems'', Phys.\ Rev.\ A 81 (2010), 032104.
\bibitem{shammah2018} N.~Shammah, S.~Ahmed, N.~Lambert, S.~De~Liberato, F.~Nori, ``Open quantum systems with local
  and collective incoherent processes: Efficient numerical simulations using permutational invariance'',
  Phys.\ Rev.\ A 98 (2018), 063815.
\bibitem{gegg2017} M.~Gegg, M.~Richter, ``PsiQuaSP -- a library for efficient computation of symmetric open
  quantum systems'', Sci.\ Rep.\ 7 (2017), 16304.
\bibitem{kirton2017} P.~Kirton, J.~Keeling, ``Suppressing and restoring the Dicke superradiance transition by
  dephasing and decay'', Phys.\ Rev.\ Lett.\ 118 (2017), 123602.
\bibitem{dicke1954} R.~H.~Dicke, ``Coherence in spontaneous radiation processes'', Phys.\ Rev.\ 93 (1954), 99--110.
\bibitem{gross1982} M.~Gross, S.~Haroche, ``Superradiance: an essay on the theory of collective spontaneous
  emission'', Phys.\ Rep.\ 93 (1982), 301--396.
\bibitem{stolyarov2002} A.~I.~Stolyarov, O.~E.~Tikhonov, A.~N.~Sherstnev, ``Characterization of normal traces on
  von Neumann algebras by inequalities for the modulus'', Math.\ Notes 72 (2002), no.~3--4, 411--416.
\bibitem{bikchentaev2005young} A.~M.~Bikchentaev, O.~E.~Tikhonov, ``Characterization of the trace by Young's
  inequality'', J.\ Inequal.\ Pure Appl.\ Math.\ 6 (2005), no.~2, Art.~49.
\bibitem{bikchentaev2007mono} A.~M.~Bikchentaev, O.~E.~Tikhonov, ``Characterization of the trace by monotonicity
  inequalities'', Linear Algebra Appl.\ 422 (2007), no.~1, 274--278.
\bibitem{bikchentaev2011pb} A.~M.~Bikchentaev, ``The Peierls--Bogoliubov inequality in C$^*$-algebras and
  characterization of tracial functionals'', Lobachevskii J.\ Math.\ 32 (2011), no.~3, 175--179.
\bibitem{bikchentaev2011mn} A.~M.~Bikchentaev, ``Commutation of projections and trace characterization on von
  Neumann algebras.~II'', Math.\ Notes 89 (2011), no.~3--4, 461--471.
\bibitem{bikchentaev2010smj} A.~M.~Bikchentaev, ``Commutativity of projections and characterization of traces on von
  Neumann algebras'', Sib.\ Math.\ J.\ 51 (2010), no.~6, 971--977.
\bibitem{cho2009} K.~Cho, T.~Sano, ``Young's inequality and trace'', Linear Algebra Appl.\ 431 (2009), no.~8,
  1218--1222.
\bibitem{bikchentaev2019abed} A.~M.~Bikchentaev, S.~A.~Abed, ``Projections and traces on von Neumann algebras'',
  Lobachevskii J.\ Math.\ 40 (2019), no.~9, 1260--1267.
\bibitem{bikchentaev2019metric} A.~M.~Bikchentaev, ``Metrics on projections of the von Neumann algebra associated
  with tracial functionals'', Sib.\ Math.\ J.\ 60 (2019), 952--956.
\bibitem{bikchentaev2020det} A.~M.~Bikchentaev, ``Inequalities for determinants and characterization of the
  trace'', Sib.\ Math.\ J.\ 61 (2020), no.~2, 248--254.
\bibitem{bikchentaev2026moslehian} A.~M.~Bikchentaev, M.~S.~Moslehian, ``Trace inequalities and characterizations
  of tracial functionals in operator algebras'', Positivity 30 (2026), Art.~24; doi:10.1007/s11117-026-01180-0.
\bibitem{bikchentaev2025survey} A.~M.~Bikchentaev, M.~S.~Moslehian, ``Characterizations of tracial functionals on
  matrix and operator algebras'', Bull.\ Belg.\ Math.\ Soc.\ Simon Stevin 32 (2025), no.~3, 343--361.
\bibitem{bikchentaev2026sgm} A.~M.~Bikchentaev, T.~H.~Dinh, A.~V.~Le, M.~S.~Moslehian, ``Spectral geometric mean
  and trace characterizations'', Positivity 30 (2026), Art.~46.
\bibitem{laflorencie2014} N.~Laflorencie, S.~Rachel, ``Spin-resolved entanglement spectroscopy of critical spin
  chains and Luttinger liquids'', J.\ Stat.\ Mech.\ (2014), P11013.
\bibitem{bonsignori2019} R.~Bonsignori, P.~Ruggiero, P.~Calabrese, ``Symmetry resolved entanglement in free fermionic
  systems'', J.\ Phys.\ A 52 (2019), 475302.
\bibitem{fraenkel2020} S.~Fraenkel, M.~Goldstein, ``Symmetry resolved entanglement: exact results in 1D and beyond'',
  J.\ Stat.\ Mech.\ (2020), 033106.
\bibitem{estienne2021} B.~Estienne, Y.~Ikhlef, A.~Morin-Duchesne, ``Finite-size corrections in critical
  symmetry-resolved entanglement'', SciPost Phys.\ 10 (2021), 054.
\bibitem{song2012} H.~F.~Song, S.~Rachel, C.~Flindt, I.~Klich, N.~Laflorencie, K.~Le~Hur, ``Bipartite fluctuations as
  a probe of many-body entanglement'', Phys.\ Rev.\ B 85 (2012), 035409.
\bibitem{lukin2019} A.~Lukin, M.~Rispoli, R.~Schittko, M.~E.~Tai, A.~M.~Kaufman, S.~Choi, V.~Khemani, J.~L\'eonard,
  M.~Greiner, ``Probing entanglement in a many-body-localized system'', Science 364 (2019), 256--260.
\bibitem{kiefer2020} M.~Kiefer-Emmanouilidis, R.~Unanyan, M.~Fleischhauer, J.~Sirker, ``Evidence for unbounded growth
  of the number entropy in many-body localized phases'', Phys.\ Rev.\ Lett.\ 124 (2020), 243601.
\bibitem{chen2026} K.~Chen, W.~Jia, X.~Li, R.~Meyer, J.~Zhao, ``Detecting symmetry-resolved entanglement: a quantum
  Monte Carlo approach'', arXiv:2604.02307 (2026).
\bibitem{tikhonov2005} O.~E.~Tikhonov, ``Subadditivity inequalities in von Neumann algebras and characterization of
  tracial functionals'', Positivity 9 (2005), no.~2, 259--264.
\bibitem{hoa2007} Dinh Trung Hoa, O.~E.~Tikhonov, ``Weighted trace inequalities of monotonicity'', Lobachevskii J.\
  Math.\ 26 (2007), 63--67.
\bibitem{hoa2010} Dinh Trung Hoa, O.~E.~Tikhonov, ``Weighted monotonicity inequalities for traces on operator
  algebras'', Math.\ Notes 88 (2010), no.~1--2, 177--182; Mat.\ Zametki 88 (2010), no.~2, 193--200.
\bibitem{sano2006} T.~Sano, T.~Yatsu, ``Characterizations of tracial property via inequalities'', J.\ Inequal.\ Pure
  Appl.\ Math.\ 7 (2006), no.~1, Art.~36.
\bibitem{hoa2013} Dinh Trung Hoa, H.~Osaka, Ho Minh Toan, ``On generalized Powers--St{\o}rmer's inequality'', Linear
  Algebra Appl.\ 438 (2013), no.~1, 242--249.
\bibitem{moslehian2025jmaa} M.~S.~Moslehian, A.~M.~Bikchentaev, ``Characterizations of the canonical trace on full
  matrix algebras'', J.\ Math.\ Anal.\ Appl.\ 552 (2025), no.~2, 129764.
\bibitem{halmos1969} P.~R.~Halmos, ``Two subspaces'', Trans.\ Amer.\ Math.\ Soc.\ 144 (1969), 381--389.
\bibitem{bikchentaev2009rm} A.~M.~Bikchentaev, ``Commutativity of projectors and trace characterization on von Neumann
  algebras.~I'', Russian Math.\ (Iz.\ VUZ) 53 (2009), no.~12, 68--71.
\end{thebibliography}
}

\end{document}
